\chapter{Sāṅkhya-yoga}
\chaptersubtitle{The Eternal Nature of the Self}

In the last chapter, we saw how Arjuna feels confused and how depressed he has become; because when the mind is focused and concentrated on the negative, everything becomes negative. We will see how Lord Kṛṣṇa gives Arjuna the knowledge of the Gita and how He starts to talk to him about the knowledge of the Self, Sankhya Yoga. Kṛṣṇa is a great psychologist; He doesn’t straight away go and gives Arjuna everything. He goes slowly with him, step by step. He doesn’t change him, but He gradually allows Arjuna to transform himself.

This is very important to know: God doesn’t force anybody to change their life. He has placed everything inside you. He has placed this great wisdom inside you. And He has given you feeling! He has made you try to discriminate between what you hear and see outside and what you hear and see inside. This is what we will discover together with Arjuna! We will travel the same path as Arjuna. Because when Kṛṣṇa is talking to Arjuna, He is not only talking to Arjuna, He is talking to everybody! He is calling everyone to awaken. If you choose to sleep, He can’t push you! He can go gradually, giving you the divine nectar, but it’s up to you to drink it and transform yourself.

In this chapter, Kṛṣṇa becomes the \textit{guru}. All this time, Arjuna has been looking at Kṛṣṇa as his cousin and his friend. But now, seeing this terrible state of Arjuna, this pitiful state of Arjuna, Kṛṣṇa starts to teach. At this moment, Kṛṣṇa becomes the \textit{guru} and Arjuna, the disciple. But Arjuna’s willingness to listen to Kṛṣṇa is very important! How can Kṛṣṇa tell something to Arjuna, if Arjuna doesn’t want to listen? In the first chapter, you saw that Kṛṣṇa doesn’t say much: He keeps quiet and listens to what Arjuna has to say. But when Arjuna is tired of talking, then Kṛṣṇa reveals to him the Gita.

Only when you can listen, can you change! If you can’t listen, how will you change? You will never change. Life is like this: if you don’t learn to listen and observe, you will not transform. That’s why in the previous chapter, we talked about self-analysing, where we see and observe ourselves. We see and observe our day, and we see and observe what we can change in our life. Arjuna is getting ready to listen to Bhagavān Kṛṣṇa, to listen to God, and learn how he can transform and change.

Before we start, we will chant a mantra:

ajñāna timirān dhasya jñānāñjana śalākayā cakṣur unmīlitaṁ yena tasmai śrī gurave namaḥ

This verse is praising the Master. “I was born in the darkness of ignorance, and my spiritual Master opened my eyes with the torch of knowledge. I offer my respectful obeisances unto Him.”

Knowledge is very important. Knowledge is not something that stays in the mind, it’s not intelligence. Knowledge is something deeper; it touches the core of your heart. In the next chapters, you will see that Kṛṣṇa says, “Above all, knowledge is the greatest.” However, further on He will say, “Above knowledge, Bhakti is the greatest!” First He emphasises that knowledge is important. But what kind of knowledge? Not what you go to school and study: that makes you become a robot! It makes you become a slave. No! True knowledge is the knowledge of the Self, knowledge which unites you with the Divine. That is why you incarnate from time to time – to realize this. That is why you are sitting here now! It’s not just to listen to what I have to say, or to what Kṛṣṇa has to say. It’s because there is a deep longing for God inside you and that’s why God has put you here. Previously, in Chapter 1, we heard the story about the parrot and the prostitute. If one hears a chapter of the Gita, whether one understands it or not, is not important. The prostitute didn’t understand anything in the Gita, but she was meditating on it. She didn’t have intellectual knowledge about it, but she had the inner knowledge. Inner knowledge has nothing to do with what is in your mind. Inner knowledge has to do with what is in your heart.

In this story from Chapter 1, I emphasised that the good merit, the \texti{puṇya}, which the lady offered to the ox when he was dying, came from her heart, not from her mind. The mind will make you feel pity and become like those you see in the outside world. But what comes from the heart, that’s what is important. When you receive spiritual knowledge from the great sages and the holy scriptures – you take it into you. But how you transform it, what you do with it inside of you, and how you bring it out as a service to humanity, this is very important!  The strength of what comes from your heart is what makes an effect in life. If you take something positive inside, yet what comes out of your heart is negative and not productive – what good can come from it? No good!

It’s how you are taking it in that matters. And that will bear its fruits: if you have taken this knowledge inside you and have meditated on it inside your heart, then, when you offer it to the Lord, it will bear a million times more fruit than just superficially offering whatever. That’s why we say, “Kṛṣṇa Arpanam,” which means ‘I offer everything to Kṛṣṇa.’ Everything comes from Him, and we offer everything back to Him. But you can’t offer it superficially. It has to emerge from deep within you: then it bears fruit, then it is productive. This verse says that everybody is born in ignorance: they don’t know the Self. Later on, Kṛṣṇa says to Arjuna, “Our relationship is much older than this lifetime! It’s just that you don’t know it. I know!”

vande ōhaṁ śrī-guroḥ śrī-yuta-pada-kamalaṁ śrī-gurūn vaiṣṇavāṁś ca

“I offer my respectful obeisances unto the Lotus Feet of my spiritual Master, and unto the Feet of all Vaishnavas.”

he kṛṣṇa karuṇā sindhu dīna bandhu jagat pate gopeśa gopikā kānta rādhā kānta namostute

“O my dear Kṛṣṇa! You are the friend of the distressed and the source of creation. You are the Master of the gopis and the lover of Rādhārānī. I praise you, I offer my respectful obeisances unto You.”

God is a friend, you know? He is the only One who is seated in the core of the heart. He is the source of all creation.

For the one who offers the devotion of a servant, Kṛṣṇa chooses to be the Master who cares and looks after that person. But the Master is not to be served. He is serving! He is caring! He is looking after the welfare of the people who love Him. He is the lover of Rādhārānī. Rādhārānī stands for all who have bhakti inside their hearts, who love Him and have a relationship with Him. But this love is not just superficially proclaiming, “Oh, I love Him, I serve Him.” No! He who has a deep love inside, who loves Him as a wife loves her husband or as a husband loves his wife, goes blindly into that love. In that state, you are the lover of God.

Many saints have loved Him like that: Chaitanya Mahaprabhu, Meerabai, Tukārām. In that state, there is no man or woman. When the \textit{ātmā} loves, it’s not the outer gender that counts. Love is beyond gender! When people are in the mind and hang onto gender, there is much judgment. But when one transcends this and sees deeply inside one’s core, one sees that the heart knows no difference. You can love and that’s what matters. This love is what matters.

tapata-kāñcana gaurāṅgī rādhe vṛndāvaneśvarī vṛṣabhānu sute devī praṇamāmi harī priye

“I offer my respects to Rādhārānī, whose bodily complexion is like molten gold and who is the Queen of Vrindavan. You are the daughter of King Vrishbhanu, and You are very dear to Lord Kṛṣṇa.”

vāñchā-kalpatarūbhyaś ca kṛpā-sindhubhya eva ca patitānāma pāvanebhyo vaiṣṇave(bhyo) namo namaḥ

“I offer my respectful obeisances unto all the Vaishnava devotees of the Lord. They can fulfil the desires of everyone, just like the Kalpataru, the wishfulfilling tree, and they are full of compassion for everyone, even for the fallen souls.”

This verse explains the qualities of a Vaishnava, a devotee of Viṣṇu. There is a beautiful song that Gandhi always loved to sing, the Narsinh Mehta’s bhajan “Vaiṣṇava jana to”:

vaiṣṇava jana to tene kahiye je pīḍa parāyī jāṇe re

In this song, Narsinh Mehta is saying, “Know the quality of a Vaishnava! The quality of a Vaishnava is to praise.” A Vaishnava is completely different from a Saivite, a Shakta, or others, because everything for a Vaishnava is based on devotion and love. That’s why Lord Kṛṣṇa says, “The servant of My devotee is My devotee.” A true Vaishnava cares for everyone. A true Vaishnava has love for all human beings, irrespective of their caste, colour, creed or anything else. A true Vaishnava represents the Lord Himself and how the Lord loves. When one calls oneself a Vaishnava, one has to love all. There is no criticising. This compassion is the sign of a Vaishnava. Whoever associates with a Vaishnava will feel the Love of God and their desires will be fulfilled. Whatever your soul deeply longs for, you will receive. Then each one will become like a wishfulfilling tree for others. When one drinks of the nectar of the \texit{Śrīmad Bhagavad Gītā} and puts it into practice in one’s life, one becomes like a wishfulfilling tree for others.

Like I said, your life is not only for yourself; your life has been given to you to spread the Love of God. To do that, you have to feel the Love for God inside you. You have to manifest this Love, from deep within your heart. And compassion: you are full of compassion. In Chapter 1, we see that Arjuna has much compassion. It is out of compassion that he feels pity for his relatives. It is out of this compassion that he says, “I would rather not fight.” If it was out of pride and arrogance, he would not care. In the verses 38-39 of Chapter 1, Arjuna says, “Their minds are clouded and they can’t see. But we can see. That is why I don’t want to fight.” And that is the beauty of someone who has love inside his heart. When love arises, there is also understanding, there is also non-judgment.

\section{Śrī Gītā Ṭīkā}

Now we’ll start the second chapter, Sankhya, which is filled with Vedic knowledge.

Arjuna, overwhelmed with grief, in a state of complete distress, overtaken by faintheartedness, by looking toward the outside, grief, as grief has no place in that situation when the battle is yet to start.

You see, if you sit and cry, “What am I going to lose?”, you can’t fight. If you look backward on the spiritual path, what is changing in your life, you will never move forward. Because you will sit and cry already for what you are losing. You will sit and cry because the change is happening.

The feet are made to look – to walk to the front, no? To move forward, not to go backward. That’s why what is finished, what is past, is past. You have to let go. Sitting and crying for the past will not change anything. Sitting and crying about what will be in the future will not change anything. Being now, in the moment, this is when you can change something.

So Kṛṣṇa, considering, looking and hearing, looking how he is, His friend is, hearing all that he is saying, He saw that this is a lack of knowledge. Due to a lack of knowledge, he is saying all these words. And on Arjuna’s part, he is talking, he is giving good reasons for not fighting, yet he doesn’t understand what he is saying. Lack of depth of understanding on the part of Arjuna.

You see, many people say many things, can talk a lot, but their word doesn’t have any power. Word is power. Word is strength. All the holy scriptures start with the word, no? The Veda, the first word of the Veda, is “word.” The power of the word is very important. Bhagavān Kṛṣṇa noticed that there is no understanding in what Arjuna was saying, crying. So, He says, “I must - now it is time.”

Meanwhile, like I said earlier, Arjuna was just talking and talking and talking, and Bhagavān was listening, not showing any emotion. You know, when somebody sits and cries with you, “Poor you!” No? You try to pity the person. But you are not helping the person by telling the person – He is already poor and then you add your “poor you” upon that person. It doesn’t change. That’s why I was giving the example of a doctor. When you go to see the doctor, the doctor wants to heal you, but to heal somebody, you have to be in a neutral state. You can’t side with that person’s side, because that’s what the person wants you to do. That’s what Arjuna was doing with Bhagavān Kṛṣṇa, saying, “Poor me, You have to listen to me, and it must be the way I want.” But Bhagavān is the doctor of all doctors, no? He is the Supreme physician. He knows.

When a problem arises, and you know it is not real, and you take it to be real -- in the last chapter, Arjuna started believing in his fantasy -- then you start looking for solutions left and right. Already, the problem is there. When you have that problem, the root of that problem, it is already a problem itself. The root of that problem is fake. It’s not real. Whatever solution you will get from that will be problematic. That’s why the \textit{guru} doesn’t get involved in drama. That is why Bhagavān is not getting involved in the drama of Arjuna.

One takes a rope to be a snake. Then the problem will still be with you, because you will have to suffer for that. But when ignorance is removed, then you have a clear understanding. But if you are still sitting with your ignorance and sitting and pitying yourself, how would you awaken when you see and sit and cry, “Why is it like this? Why is it like that?” Would you change? No, it will not! It will change only when you take shelter and let go, take shelter at the Feet of the Master and let go of that worry and that pain.

You see, everybody wants to be free. Tell me who doesn’t want to be free? Everybody is doing everything to free themselves from pain and sorrow. But merely just chanting, “I want to be free,” and not doing anything to be free, it doesn’t help. It makes you weak, actually. But when you make your effort, then you don’t need to do much, because you will have that inner strength inside of you to move forward. Once that ignorance is gone from you, then truly you can be free. If you still hang on to ignorance, you will sit and cry like Arjuna.

\Verse[2.1]
{% --- 1. The Verse ---
  \hspace*{1em}sañjaya uvāca \\
  taṁ tathā kṛpayāviṣṭam aśru-pūrṇākulekṣaṇam \\
  viṣīdantam idaṁ vākyam uvāca madhusūdanaḥ
}
{% --- 2. The Word-for-Word (Clean text, no formatting) ---
  \textit{sañjaya uvāca}---Sañjaya said; 
  \textit{tam}---to him (Arjuna) [who was]; 
  \textit{tathā}---thus; 
  \textit{viṣīdantam}---despondent; 
  \textit{kṛpayā-āviṣṭam}---[and] overwhelmed with pity; 
  \textit{ākula-īkṣaṇam}---[his] eyes full of confusion; 
  \textit{aśru-pūrṇa}---[and] filled with tears; 
  \textit{madhusūdanaḥ}---Madhusūdana (Kṛṣṇa); 
  \textit{uvāca}---spoke; 
  \textit{idam}---these; 
  \textit{vākyam}---words
}
{% --- 3. The Translation ---
  \hspace*{1em}Sañjaya said: \\
  To Arjuna, who was thus despondent and overwhelmed with pity, his eyes filled with confusion and tears, Madhusūdana spoke these words:
}

In this state, Arjuna is completely distressed and overtaken by faint-heartedness. Looking at his friends and relatives, he starts to cry. He is not just talking, he is crying terribly. He is overtaken by grief and fear that the family will be destroyed. That will bring great sin and negative \textit{karma} to him, to his family, and to their descendants. It will be complete destruction. In this verse, it is very important to understand why Sanjaya doesn’t call the Lord, ‘Kṛṣṇa’. Instead he refers to Lord Kṛṣṇa as madhusūdanaḥ, which means the slayer of Madhu. Madhu was a very fierce demon who was killed by Kṛṣṇa.

Sanjaya is reminding Dhritarashtra, the evil-minded, blind king, that Lord Madhusudana Himself, Sri Kṛṣṇa Himself who had slain the demon Madhu, is there as Arjuna’s charioteer. He is warning Dhritarashtra saying, “He killed this terrible demon -- so what about your sons who are on the battlefield? They’re nothing compared to the demon Madhu! This great demon had been terrorising not only the people, not only the world, he had even been terrorising the devas in Indra Loka! How powerful Kṛṣṇa is! And look at your sons, what are they? Sorry, they are a bunch of idiots. Even when Kṛṣṇa went to see them as a messenger of peace, they chose to be blind. They have such stupidity, they don’t see. They don’t see that the great Madhusudana is here as Arjuna’s charioteer.”

Everybody, including the Kauravas, knew what Kṛṣṇa had done in His youth in Vrindavan. They knew about all the demons He had killed. But they still chose to be ignorant. Sanjaya continues, “The slayer of that great demon is here. Surely the Pandavas will be victorious! The Pandavas will win! Don’t count on your sons to win!” It is this Lord – the One who killed the demon Madhu – it is Him, who will now talk to Arjuna, who is depressed and unwilling to fight. Sanjaya warns, “Dhritarashtra, be ready, because all these tyrants, all these persecutors, will be removed. If the Lord Himself has chosen to be present here, how could your sons be victorious? They can’t be!” Sanjaya is reminding him that even if Arjuna is overtaken by distress and faintheartedness, Kṛṣṇa is there to change everything.

Sanjaya was revealing this to Dhritarashtra after the tenth day of the war. He said, “If you want peace, tell your children to stop. There is still time! Your hundred sons have not yet been slain. As long as Bhishma Pita was still fighting, nothing could happen to any of them, because he was a big shield for them. But now, what will happen? Now it’s imminent that your hundred sons will die.”

You see, when people start on the spiritual path, when they feel this inner strength, even if they look outside and feel pity for others, the negativity of the outside will not affect them. They will be strengthened; they will have a shield around them and feel protected.


\textbf{2016}---Sanjaya is saying to the king: “To Arjuna, who was thus despondent and overwhelmed with pity, his eyes filled with confusion and tears, Madhusūdana spoke these words.” Sanjaya was eager, waiting. Arjuna was also waiting to hear what Bhagavān Kṛṣṇa will say to him. In his mind, he wanted that Kṛṣṇa says what he wants, but a good doctor doesn’t give you the medicine you want. Otherwise, why would you go to the doctor? You would just buy the medicine and drink it yourself, no? Like we have self-doctors nowadays. In each family, you have somebody who knows better than the doctor themselves.

Sanjaya is reminding the blind king that the Lord Himself, who killed the demon Madhusudana, who killed the demon Madhu, is there with Kṛṣṇa. This demon Madhu was a very fierce demon, but Kṛṣṇa killed him. This terrifying monster was killed by Kṛṣṇa. So here, Sanjaya is reminding how powerful Kṛṣṇa is. He is reminding the blind king that, “Don’t be a fool thinking that Kṛṣṇa is just a normal human being. Here, He who is the charioteer of Arjuna, I’m giving you a warning now, He is not an ordinary person. He has killed the demon Madhu, who was terrorizing all the three worlds, the people, the devas. He was so powerful, but He is with the Pandavas.” Here he is saying, “Look at your sons, if they don’t realize who is in front of them, they are the worst of idiots, the worst of their kind.”

You have different degrees of idiocy. Sometimes you find a very high degree. Because, look, Bhagavān Kṛṣṇa went there, tried to convince them, He said, “Look, I came very peacefully to you. I am with them.” He even showed His Cosmic Form. They saw Him, who He is. Some of them, at least. Those who are blind will stay blind. But due to their degree of stupidity, they don’t want to see.

So He is there sitting, do you have any chance against Him? No, you don’t! They all have heard all the stories about what Kṛṣṇa did in His childhood. How all the demons who have been sent by his uncle Kamsa, how He got rid of them. That’s why He didn’t fight. Imagine Kṛṣṇa would fight, what would happen? The fight would not even move. But, in reality, at the end, you will see it is all a big drama. That’s why He is saying that, “Look, Arjuna is still distressed due to his faintheartedness, but Kṛṣṇa will change everything.”

\Verse[2.2]
{% --- 1. The Verse ---
  \hspace*{1em}śrī-bhagavān uvāca \\
  kutas tvā kaśmalam idaṁ viṣame samupasthitam \\
  anārya-juṣṭam asvargyam akīrti-karam arjuna
}
{% --- 2. The Word-for-Word (Clean text, no formatting) ---
  \textit{śrī-bhagavān uvāca}---Bhagavān Kṛṣṇa said; 
  \textit{arjuna}---O Arjuna; 
  \textit{kutaḥ}---from where; 
  \textit{samupasthitam}---has come; 
  \textit{tvā}---to you; 
  \textit{idam}---this; 
  \textit{kaśmalam}---weakness; 
  \textit{viṣame}---in this hour of danger; 
  \textit{anārya-juṣṭam}---[it is] unworthy of noble men; 
  \textit{asvargyam}---not leading to Heaven; 
  \textit{akīrti-karam}---causing dishonor
}
{% --- 3. The Translation ---
  \hspace*{1em}Bhagavān Kṛṣṇa said: \\
  O Arjuna, from where has this weakness come in this hour of danger? It is unworthy of noble persons and will not lead to the attainment of Heaven. Instead, it causes dishonor.
  }

Here Lord Kṛṣṇa Himself starts speaking. He says that Arjuna’s heaviness of heart and low spirits are due to over-sentimentality. He says, “You have made yourself weak. It is your willingness to become weak!”

Kṛṣṇa continues, “I am taken aback! I am shocked to see that here, in the middle of the battlefield, you are in such a state! This state is very dangerous and difficult! We are in the middle of the battlefield and now you feel like this? If you would have felt like this before coming on the battlefield, I would understand. But now we are here. This is not a noble feeling. It is not the feeling of a noble man. There is no place for sadness, cowardice or pain on the battlefield. Elsewhere, yes! There is a place for everything. But this is completely out of place here! How can you fight all these great heroes if you are so fainthearted? With a heart like this, it is very difficult to fight. Where did you get this faintheartedness? Where did it come from?” Kṛṣṇa has never known Arjuna to be like this – very sentimental and soft-hearted. Calling this dejection and faint-heartedness of Arjuna unworthy of a noble soul, He says, “If you are like this, how can you even talk about gaining Heaven?” If you are weak, you can’t expect God-realization! If you are weak, you can’t expect God to act through you! He will not reveal your \textit{dharma} to you! God knows that there is no use in throwing pearls to pigs. They will not know what to do with them. They will trample over them; they will not value them. Arjuna’s faint heartedness is not worthy of a warrior on the battlefield.

If you are weak on your path in life, you can’t do anything. If you are weak on your path towards God-realization, how can you attain God? Kṛṣṇa says, “I am really amazed. O Arjuna, this is not the way cherished by the noble man: this mood came not from Heaven, nor can it lead to Heaven. This weak kind of mood, this weakness will not lead you anywhere. You are noble.” “Noble” here doesn’t mean that one is of the noble or royal class. By saying to Arjuna, “You are noble,” Kṛṣṇa means, “You are strong, you are powerful! You are \textit{dharma}! You have the knowledge of the Self. This is not the place for weakness. This kind of thing will not inspire others. All know you as a great warrior.”

Do you remember before the war, when King Virata’s son was going to fight the Kauravas? Arjuna was still disguised as the woman, Brihannala. Though he was dressed as a woman, everyone was always talking about how strong he was. When Arjuna told the son of King Virata, “Go and get my bow, my Gandiva,” he asked, “How will I know what your bow looks like?” Arjuna said, “There is a tree in front of me. Go! You will find my bow there. It is as tall as a palm tree.” That’s how tall the bow of Arjuna was! So you can imagine how tall the Pandavas were! How strong and powerful they were!

Lord Kṛṣṇa continues, “This weak kind of mood will not lead you to Moksha. There are four types of noble qualities in life: Moksha (salvation); \textit{dharma} (virtue); artha (wealth); and kama (enjoyment). This is not how to achieve them! To attain these noble qualities, you have to be strong in life! How can you talk about attaining salvation when you are in this pitiful state? Especially at this odd hour, when you are on the battlefield, My dear!” In Chapter 1, I said -- what if you go to the doctor and the doctor just starts crying with you – how can the doctor help you?

Here Kṛṣṇa becomes the \textit{guru} and says, “Look at this world! These excuses are nothing. Your \textit{dharma} is to fight!” This could appear very terrible! Some people would say, “But Kṛṣṇa is God, He could change everything!” Of course He could change everything! He can change everything by His will. But, no. He planned to give this knowledge to the world, to remind the world of this knowledge.

\section{Śrī Gītā Ṭīkā}
“Bhagavān Kṛṣṇa said: O Arjuna, from where has this weakness come in this hour of danger?” “Arjuna, when did you become like that? I heard that you are so powerful, so strong, but I don’t see in front of me somebody who is strong and powerful. But I see a very weak person standing in front of problems and crying. How would you be free from that? How would you attain peace and happiness? When has this come upon you?”

“O Arjuna, this is not the way cherished by the noble man. You are a noble person. You are a prince, and a prince should not sit and cry like that.”

“It is unworthy of noble persons.” Here Bhagavān Kṛṣṇa said, “Don’t think that this mood is given by God to you.”

“And will not lead to the attainment of Heaven.” By sitting and crying, thinking, “Oh, I will sacrifice myself, I will go to Heaven.” He said, “No, it will not lead you to Heaven. And on Earth, if you are looking for glory, forget it, they will consider you a coward. It will not lead you anywhere. If you are weak on the path of life, you can’t achieve anything. If you are weak on your path toward God-Realization, do you think God will reveal Himself to you? No! Here Bhagavān Himself said, “I am really amazed to see you in such a pitiable condition. O Arjuna, this is not the way you should be. You should be strong. Because when you are weak, you don’t achieve anything.”

In life itself is like that. If you want to achieve something, you have to be strong, you know? Bhagavān Kṛṣṇa reminds him, “You are strong, you are very powerful, you are very dharmic, you have the knowledge of the Self. This is not the place, this is the wrong place to be weak. Once you have been born on Earth, you should not be weak, especially if you are second-born.

You see, first-born means you are born from the mother, second-born means you are born from your spiritual path. All scriptures say that. Religion also says that. In Hinduism, they say you are born two times; the first time when you come from the mother, the second time when you receive the blessing of the Sat\textit{guru}.

Here Bhagavān Kṛṣṇa is reminding him, “You should never be weak, but because you have something greater, you have the knowledge of the Self. Sitting in your weakness, you know that you are wrong, and this will not help anybody. It will not help you. It will not help anyone, because everybody knows you as being a great warrior.”

He said, “You are from the Aryan clan. Anybody from that clan has great Vedic culture. And through this Vedic culture, you must be very strong.” Kṛṣṇa was a bit cheeky and started smiling, and said, “You are so great, everybody knows you as being a great warrior who has fought so many battles.” With a smile, He asked Arjuna, “How are you running away now? This is the wrong place. You could have run away before! Why are you so hopeless right now? How did you discover that you have such a weakness at such a wrong time?” A very big problem for himself. He is reminding him that problems come in this walking state.

When you start, you have a certain perception that you can’t achieve something. Due to this perception that you can’t achieve something, it makes you weak, and it brings more problems. It’s like you are entering a state, like being in a coma. When you are in a state of coma, in a dead state, vegetable state, can you do something? No, you can’t do anything!

So, He said, “Awake. Rise. Remind yourself of being the Spirit. If you choose to sleep in a deep state, nothing can come out of you. And if you think that you are just the body, then you will be in big problems. Awake, Arjuna. This is not the right place to sleep. You have already made the first step of coming to this war, now go till the end. But not in weakness, but in strength.”

Here, Bhagavān Kṛṣṇa revealed to him that He is the Supreme \textit{guru} Himself. Now He will teach him. He saw that Arjuna had this knowledge, but he is tainted by other things which need to be polished. You see, when you leave some things for a long time, it catches dust, no? Or sometimes gets rusty. And if you want to reuse it, you have to dust it and polish it again. Due to the connection that Arjuna had to Kṛṣṇa, not only in this life itself, but in many lives, that connection is still there, but sometimes it gets dusty and rusty. Probably it was a bit rusty that time, so Bhagavān had to polish him a little bit. He’s getting a good polishing cream.

\Verse[2.3]
{% --- 1. The Verse ---
  klaibyaṁ mā sma gamaḥ pārtha naitat tvayy upapadyate \\
  kṣudraṁ hṛdaya-daurbalyaṁ tyaktvottiṣṭha parantapa
}
{% --- 2. The Word-for-Word (Clean text, no formatting) ---
  \textit{mā sma}---do not; 
  \textit{gamaḥ}---go into; 
  \textit{klaibyam}---cowardice; 
  \textit{etat}---this; 
  \textit{na upapadyate}---it is not befitting; 
  \textit{tvayi}---in you; 
  \textit{tyaktvā}---throw away; 
  \textit{kṣudram}---[this] lowly; 
  \textit{hṛdaya-daurbalyam}---weakness of heart; 
  \textit{uttiṣṭha}---arise; 
  \textit{parantapa}---O vanquisher of enemies
}
{% --- 3. The Translation ---
  Do not give in to this cowardice, O Arjuna, it does not befit you. Throw away this lowly weakness of heart and arise, O vanquisher of enemies!
}

Lord Kṛṣṇa addresses Arjuna as Paartha, which means son of Pritha (Kunti). He says, “You are the son of a brave lady. You are the son of Kunti, who is a true Kshatriya, a true warrior. Your mother is so strong! You are a hero! So don’t go into this state of cowardice; don’t let yourself be overtaken by weakness. You possess great valour and are skilled in fighting. Let go of this ‘faintheartedness and arise!’ You strike fear into the hearts of your enemies; you strike terror into the hearts of great heroes. The moment they hear of you, they are scared! The moment they think of you, their hair stands on end! ‘Scourge of your enemy!’ The moment they hear your Gandiva, they have a powerful image of you. But if you are trembling and your Gandiva is falling from your hand, what will they think of you? This is not worthy of you, My dear.”

Then Kṛṣṇa addresses Arjuna saying, tyaktvottiṣṭha parantapa, “Attack them!” Here Kṛṣṇa did not only mean to attack the outer enemy. He also meant to attack and remove one’s negative inner qualities. The enemy is not only outside; it is inside each of us. Look at your inner enemy. Remove it! Attack it! Don’t be afraid or weak. Remove all this stuff inside! Be in your power! Kṛṣṇa addresses Arjuna as parantapa. Arjuna had this name because in the past he had killed many demons. Kṛṣṇa says, “All the demons you have killed, all the other big, big baddies -- you have terrified them all. So get up! Stand up! Arise! This is not the place for weakness! Fight!”

Here Kṛṣṇa is going very slowly with Arjuna. When one goes on the spiritual path, one always looks at what one will lose, and thinks, “Okay. I’m losing this and I’m losing that.” But if you want to go deeper to attain God Consciousness, to attain the Grace of God, there should not be weakness in your spiritual \textit{sādhana}. You have to go forward with strength, with full trust; and with that trust, you know that God is with you. Here, God Himself is reminding Arjuna, “I am with you, all the time. And not only now. I was always with you – even when you were killing all the big, big demons, all the evil people. I was there with you. I am ever with you. You are My creation. You are in Me, and I am in you. We are not separate. We are never far away from each other.”

\textbf{2016}---“Do not give in to this cowardice.” He said, “Don’t, you are more, you are a great warrior. You come from a brave lineage of warriors. Don’t be weak. Remind yourself of where you come from.”

“It does not befit you. Throw away this lowly weakness of heart and arise, O vanquisher of enemies!”

Bhagavān is saying to Arjuna, “You are a warrior and you are actively participating in this war.” You see, when you incarnate, you are actively participating in coming here, no? It’s your willingness to come here to do what you have to do. You should not be fainthearted. When he was on the battlefield, so depressed that he wanted to leave, the charioteer had to remind him that he is a warrior. The companions who you keep around you must always be for your progression on your spiritual path. If you associate yourself with people who will bring you down, they will not wish good for you. A true friend will always encourage you to move forward in life and will always encourage you to keep away from things which will harm you. Because if friends, people you associate yourself with, encourage you to do wrong, they're not your friends. Today, they encourage you to do wrong to others, tomorrow they will do it to you.

Here, that’s what you see, Bhagavān is a charioteer, and He is reminding. Arjuna represents the soul, and Bhagavān Kṛṣṇa represents the inside of the soul, which is reminding the soul why you have come here. It’s not to get yourself entangled in this drama, which you call life, but you have come here to make life happen. You have to take counsel from the inner, take advice from Bhagavān Himself. Bhagavān is reminding Arjuna, He takes that function.

Bhagavān Kṛṣṇa said, “I will not fight.” He never said that He will not give counsel. Giving counsel is not in the agreement which they settled. His assignment as a charioteer, He has to look after it. He is doing His job. He is boosting the morale of Arjuna. He used the word, “O vanquisher of enemies, Paranthapa.” “You come from a great dynasty. You are the son of a brave lady, Kunti, who is a true fighter, a true warrior. Your mother is strong, and you are a hero, so don’t be in this state of dejection. Don’t let yourself be weak, because you are very strong. You are very skilled in fighting. Let go of these weaknesses. Don’t hold onto it, don’t make it yours. In reality, it is not you who should fear them, but they should fear you.” That’s why He addressed him as Paranthapa. “You strike fear in your enemy. Then the enemies hear your voice, they get scared, and now, when they see you, what would they think? ‘Look at him.’ Of course, they would already claim victory, glory and happiness, especially what you are saying, how you are acting. This is not the right place for it.”

Bhagavān keeps giving him good counseling, reminding him where he comes from. As a teacher, Bhagavān Kṛṣṇa is trying to reason with him, put him on the right path.

\Verse[2.4]
{% --- 1. The Verse ---
  \hspace*{1em}arjuna uvāca \\
  kathaṁ bhīṣmam ahaṁ saṅkhye droṇaṁ ca madhusūdana \\
  iṣubhiḥ pratiyotsyāmi pūjārhāv ari-sūdana
}
{% --- 2. The Word-for-Word (Clean text, no formatting) ---
  \textit{arjuna uvāca}---Arjuna said; 
  \textit{madhusūdana}---O Madhusūdana; 
  \textit{ari-sūdana}---destroyer of enemies; 
  \textit{katham}---how; 
  \textit{aham pratiyotsyāmi}---shall I fire against; 
  \textit{iṣubhiḥ}---arrows; 
  \textit{bhīṣmam}---at Bhīṣma; 
  \textit{ca}---and; 
  \textit{droṇam}---Droṇa; 
  \textit{saṅkhye}---in battle; 
  \textit{pūjā-arhau}---[who are] worthy of worship
}
{% --- 3. The Translation ---
  \hspace*{1em}Arjuna said: \\
  O Madhusūdana, destroyer of enemies, how can I fire arrows against Bhīṣma and Droṇa in battle, who are worthy of worship?
}

Kathaṁ means ‘I am surprised’. Arjuna says, “I am surprised! How can I kill these people? They are not demons! They are not really my enemies! On the contrary, they are the great elders who are worthy to be put on the altar and be worshipped. This is terrible! It would be a big sin if I attack them! We should have respect for our elders.” As I said in Chapter 1, the elders were well-respected at that time. When you talk to elders, you should speak to them respectfully. Using harsh or abusive words is a sin. That’s why Arjuna says, “These elders are worthy to be worshipped and we shouldn’t do anything disrespectful. We don’t even dare look in their eyes!”

Here again, Arjuna addresses Kṛṣṇa as Madhusudana, “the slayer of enemies.” He says, “How powerful and great You are! How will I pierce them with my arrows?” He is surprised that the Lord Himself is telling him to fight them -- because he knows that these people were the gurus, the teachers. Arjuna says, “We can’t attack them!”

\Verse[2.5]
{% --- 1. The Verse ---
  gurūn ahatvā hi mahānubhāvān \\
  śreyo bhoktuṁ bhaikṣyam apīha loke \\
  hatvārtha-kāmāṁs tu gurūn ihaiva \\
  bhuñjīya bhogān rudhira-pradigdhān
}
{% --- 2. The Word-for-Word (Clean text, no formatting) ---
  \textit{hi}---indeed; 
  \textit{śreyaḥ}---[it would be] better; 
  \textit{bhoktum}---[and] live; 
  \textit{iha}---in this; 
  \textit{loke}---world; 
  \textit{api}---even; 
  \textit{bhaikṣyam}---by begging; 
  \textit{ahatvā}---instead of killing; 
  \textit{mahā-anubhāvān}---[these] noble; 
  \textit{gurūn}---venerable elders; 
  \textit{tu}---but; 
  \textit{bhuñjīya}---our enjoying; 
  \textit{artha-kāmān}---[of] avaricious; 
  \textit{bhogān}---pleasures; 
  \textit{eva}---indeed; 
  \textit{hatvā}---after killing; 
  \textit{gurūn}---the elders; 
  \textit{iha}---here; 
  \textit{rudhira-pradigdhān}---stained with blood
}
{% --- 3. The Translation ---
  It would be better to live in this world, even by begging, than to kill these noble, venerable elders. Our enjoyment of avaricious pleasures, after killing them here, would certainly be stained with blood.
}

Here Arjuna is referring to Kripacharya, and Dronacharya: “These people are worthy to be praised and worshipped. How will I slay them? How will I be free from that sin? Even though these noble characters have fallen low, I can’t dishonour them. Even though I am a Kshatriya, even though I am a warrior, I would rather live on alms than commit this great sin for enjoyment and pleasure.”

Here, we have to see that Arjuna had very high knowledge. For one who is realized, the ‘good’ and the ‘bad’ become one. However, Realisation also means that one has to let go of everything. Arjuna was wise, because he knew this. But here he is saying, “How can I let go of all this? I can’t let go of it! It’s not right, if out of revengefulness, I massacre these elders who are worthy to be praised. How will I enjoy what I gain? I will not get salvation or good merit from this and will not even enjoy the kingdom which afterwards will be mine. I will feel guilty.”

As in Chapter 1, Arjuna is still continuing to try to reason with Kṛṣṇa. He is still very attached to the world, though not just blindly. He is saying, “I want to fight and I don’t want to fight.” He is still in this unsettled state of mind, with, one foot in and one foot out of the war.

\Verse[2.6]
{% --- 1. The Verse ---
  na caitad vidmaḥ kataran no garīyo \\
  yad vā jayema yadi vā no jayeyuḥ \\
  yān eva hatvā na jijīviṣāmas \\
  te ’vasthitāḥ pramukhe dhārtarāṣṭrāḥ
}
{% --- 2. The Word-for-Word (Clean text, no formatting) ---
  \textit{na vidmaḥ}---we don't know; 
  \textit{etat}---this; 
  \textit{ca}---as well; 
  \textit{katarat}---which of the two; 
  \textit{garīyaḥ}---[would be] better; 
  \textit{naḥ}---for us; 
  \textit{yat vā}---whether; 
  \textit{jayema}---we should defeat them; 
  \textit{yadi vā}---or whether; 
  \textit{jayeyuḥ}---they should conquer; 
  \textit{naḥ}---us; 
  \textit{yān}---those; 
  \textit{dhārtarāṣṭrāḥ}---sons of Dhṛtarāṣṭra; 
  \textit{hatvā}---after killing [whom]; 
  \textit{na jijīviṣāmaḥ}---we would no longer wish to live; 
  \textit{te}---they; 
  \textit{avasthitāḥ}---are standing; 
  \textit{eva}---verily; 
  \textit{pramukhe}---before us
}
{% --- 3. The Translation ---
  We do not know which of the two is better for us—defeating them or being conquered by them. Those sons of Dhṛtarāṣṭra after killing whom we would no longer wish to live, are verily standing here before us.
}

Here again Arjuna is showing that he is in a state of confusion, because he is still very attached to the outside. You see, if somebody is not fully dedicated, they always go into a state of confusion. They are not making themselves ready. Arjuna is in this state and says, “I don’t yet know if we will be the winners or they will be the winners. Here are the sons of Dhritarashtra on the battlefield in front of us.” He is still trying to find ways to not fight them. He is indicating that he has the choice: either to fight or not to fight. If he doesn’t fight, it will be a great offence of a Kshatriya, of a fighter, to not fight and to run away from the battlefield. How will he face his family? How will he face his people? If he runs away and doesn’t fight, it will also be a great sin and an evil act. This state of confusion makes him unsure.

In the next part of the verse, Arjuna says, “Those sons of Dhṛtarāṣṭra after killing whom we would no longer wish to live, are verily standing here before us.” Arjuna knows what to do, but at the same time, he is expressing confusion. In the state of confusion there are always two parts in it. There is a knowing, there is a selfconfidence, that, “Yes, I feel what is right to do!” The feeling of the heart is there and you know that it is right. But then the mind starts to reason. There is a fight between what you feel in your heart and what your head is saying to you: this makes the state of confusion. If your feet are more in the outside world, of course, then the mind will win. But if the feeling is stronger, then you will try to take the right decision. The state of confusion is the beginning of one’s path, the beginning of one’s search. In our confusion, we keep asking, “Am I right to be on this path or not? Am I right to do this or not?” Even the word of the \textit{guru}, at that moment, won’t mean much to the disciple, because he is still searching. He is still wandering. And the outside samskaras are still very powerful. The attachment to the outside world is still so strong that one can’t really think or see.

\Verse[2.7]
{% --- 1. The Verse ---
  kārpaṇya-doṣopahata-svabhāvaḥ \\
  pṛcchāmi tvāṁ dharma-sammūḍha-cetāḥ \\
  yac chreyaḥ syān niścitaṁ brūhi tan me \\
  śiṣyas te ’haṁ śādhi māṁ tvāṁ prapannam
}
{% --- 2. The Word-for-Word (Clean text, no formatting) ---
  \textit{kārpaṇya-doṣa-upahata-svabhāvaḥ}---with my nature overcome by the stain of weakness; 
  \textit{dharma-sammūḍha-cetāḥ}---[and] my mind confused about my duty; 
  \textit{pṛcchāmi}---I urge; 
  \textit{tvām}---You; 
  \textit{brūhi}---[to] tell [me]; 
  \textit{niścitam}---clearly; 
  \textit{tat}---that; 
  \textit{yat}---which; 
  \textit{syāt}---would be; 
  \textit{śreyaḥ}---better; 
  \textit{me}---for me; 
  \textit{aham}---I [am]; 
  \textit{te}---Your; 
  \textit{śiṣyaḥ}---disciple; 
  \textit{tvām prapannam}---who has taken refuge in You; 
  \textit{śādhi}---[please] instruct; 
  \textit{mām}---me
}
{% --- 3. The Translation ---
  With my nature overcome by the stain of weakness and my mind confused about my duty, I urge You to tell me clearly what would be better for me. I am Your disciple who has taken refuge in You. Please instruct me.
}

The first word here, kārpaṇya, means “the stain of weakness.” “It is poorness of spirit that has smitten away from me, my true heroic nature.” Arjuna knows that his noble nature, his powerful nature, has fled. 

“... and my mind confused about my duty” He is seeing what is right and what is wrong from the point of view of the mind. His mind is showing him what is right and what is wrong – not his heart, not his consciousness.

“I am Your disciple who has taken refuge in You. Please instruct me.” Here Arjuna is making himself ready. He says, “Look, I am truly confused. I feel terrible. I don’t know where I stand.”

In the Gita, Kṛṣṇa often calls materialistic people who are attached to the world ‘kripana’. They miss the purpose of life and depart from this Earth without knowing God, without Realisation. Arjuna is neither greedy nor miserly. It is not out of greed that he is in this state of confusion. His attachment to the world is not for his personal benefit. When you want to find God, if it is for personal gain, and done with pride, it will not lead you anywhere. But if it is to surrender to a greater cause, even if it is not done correctly, if it is not for your personal gain, it will bear its fruit. Arjuna had a great nature. He didn’t want to enjoy himself. He was not thinking, “Oh, I will fight this war for my own enjoyment.”

Here he is saying, “I don’t know what is right, You are Kṛṣṇa, You are God Himself! But You are telling me the opposite of what I am thinking.” Very often the \textit{guru} says something which is the opposite of what the devotee thinks, and the devotee says, “Ah! How is this possible?” But then you must ask: who is the \textit{guru} and who is the disciple? Here it’s the same thing. Arjuna is expressing something and expecting Kṛṣṇa to say nice words to him. But Kṛṣṇa doesn’t say any nice words to him, so he is taken aback, he is shocked! This also humbles Arjuna so that he is able to accept the Lord as his \textit{guru}. He surrenders to Lord Kṛṣṇa and says, “Enlighten me.” This shows that Arjuna is very wise: but he is not yet enlightened.

Arjuna was Kṛṣṇa’s best friend. Both saw each other as, “We are best friends.” They always had great respect for each other, but it was mostly as friends, as cousins. They did not have a Guru-disciple relationship. But now Arjuna has changed. In the past, Arjuna always felt equal to Kṛṣṇa. When they sat and ate, even when they slept, they would sleep in the same room, on the same bed, for they were friends.

However, Arjuna’s confused state makes him realize that even if one thinks of one’s self as being equal to God, God is superior. Even if His actions are the same, in all His actions, He is far superior to the human way of understanding. That’s why Arjuna says, “I am Your disciple who has taken refuge in You.” Arjuna doesn’t say, “Because You are my best buddy, I am going to listen to You now!” No! To take refuge at the Lotus Feet of the Lord, to take refuge at the lotus feet of the \textit{guru}, one has to put away one’s own pride and one’s own thinking. Arjuna says, “Advise me! Counsel me now! Instruct me! Before I was thinking of You as being the same as me, but now I realize that You are far superior to me.” That’s why Lord Kṛṣṇa also says, “My devotees, who are completely surrendered to Me, are far superior, because I live in them, and they live in Me.”

Arjuna has now realized that Kṛṣṇa’s advice is very important and that he needs to be guided by a true guide: and this true guide, is Kṛṣṇa Himself. Only through that, will he be taken out of his grief and be able to attain the Divine, be able to attain this inner strength so that he can move fearlessly and fight. Who else there can be his best teacher? It is not Dronacharya, not Kripacharya, not Bhishma. It is God Himself.

Even though Kṛṣṇa has been Arjuna’s best friend, now He becomes his \textit{guru}. The instruction of the \textit{guru} is like a seed which is planted in fertile land. If the disciple’s heart is like a stone, nothing will grow from it. 

Here, Arjuna and the Lord are playing the drama of the \textit{guru} and the disciple. When one is ready deep within one’s heart, when the heart becomes like fertile soil, then one is ready to receive: before that, no. Therefore Arjuna says, “Lord, I am Your disciple. Take me as Your disciple. I take refuge at Your Feet.”

There are many types of disciples and many types of devotees. It also depends on one’s self-effort and how one is centred: is one centred in the ego, is one centred in pride, or is one centred in the spirit? This verse also says that whoever is not ready will look elsewhere for things to please themselves. They will not look for the Sat\textit{guru}. But when they are ready, they will look for the Sat\textit{guru}. And once they find the Sat\textit{guru}, they will have the benefit of having the \textit{guru} always with them.

Here, therefore, Arjuna declares his discipleship and surrenders to Kṛṣṇa. He says, “I am not only Your disciple, but I have taken refuge, I have surrendered to You.” What does it mean, ‘to surrender’, ‘to take refuge’? You can be a disciple, or you can even be a part-time disciple who is not surrendered. But Arjuna says, “More than just as a disciple, more than just as a devotee, I surrender myself to You.” In the last part of the verse, the word ‘prapannam’, means to take refuge, to surrender to God. It is not just the willingness to surrender, it is also accepting that He is greater. When a disciple takes refuge in the \textit{guru}, the disciple must completely accept the superiority of the Master. This is what Arjuna is accepting. He says, “I am surrendering to You, and I accept Your superiority: only You can take me out of this ignorance. I throw myself at Your Feet. You are God.” With this, Arjuna’s eyes are brightened and his mind is completely transformed. His whole attitude has changed. With folded hands, he knows that by taking this role of the disciple, that the Lord, the God Almighty, the omniscient and the knower of all hearts, has taken the form of the supreme Master. And that He, who is full of Love, greatness, virtue, knowledge, He who is non-attached, who can’t be touched by karma, by anything, is his dear friend. And from being his best friend, He has become the supreme Guide and the supreme Divinity.”

This is the greatness of the \textit{guru}. The \textit{guru} does not have just one role. He is not just a teacher. He takes different forms, a multitude of aspects. Because the life of the \textit{guru} is not for Himself, but for others. The help and support, the knowledge, the power, the affection that the \textit{guru} has for the devotee is amazing and exquisite.

Here Arjuna is also expressing his desire to know more, and he also has in his heart the attitude of surrender to the Lord. When one has this desire to know more, to go deeper and deeper, to not just know in the mind, but deep within the heart, one will happily surrender their life to God. One knows that the only aim in life is Him, is to completely surrender to Him. Then one is ready to receive the true knowledge. When one is not ready, when one still dwells on the outside knowledge, when one still dwells on things which one thinks are good for oneself, one will always say, “I want this, I want that; I want this and I want that.” They will never ask, “What does God want from me?” In their delusion, they will always think that they are right and that it is God’s Will for them. But this is just an image that one creates in one’s mind.

Later on, we will see how Kṛṣṇa enlightens Arjuna and makes him receptive so that he can receive this gift of surrender. He explains to Arjuna what it means to surrender all his duty to Him.

\textbf{2016}---After seeing his weakness, Bhagavān Kṛṣṇa has reminded him of who he is, where he comes from, and he realizes he can’t do it by himself. “I need help. Who can really free me? Who can really help me to go out of this delusion? Who can truly help me to see what is truly right and wrong? Who can challenge this mind, that mind which is already polluted in the poorness of spirit that has smitten away from me?” The mind is already poor. The mind is already weak. What can come out of that weak mind? Nothing. But only when that mind takes shelter at the Supreme Lord Himself, then that mind can elevate itself.

Those who are attached to the world, a person in a state of confusion, is quite likely to indulge themselves in adharma. For them, due to the ignorance inside of them, the ignorance pervades the mind, controls everything, and stops them from thinking. How to truly change that mind? Bhagavān Kṛṣṇa here said, “Look, to change that crazy mind you have to have true understanding, you have to have true knowledge. Only true knowledge can dissolve this ignorance. Without true knowledge, impossible! Only when you truly understand, you have discrimination, can you be free. If you don’t have discrimination, you will never be free. Then you stay in a state of coma. Then you don’t need to do anything, you will just say, ‘I can’t change, I can’t make any effort, I don’t want to understand.’” If somebody is in that state, it’s very difficult. But when you take shelter at the Feet of the Master and when you truly want to know yourself, to realize the Self, then the freedom, the free nature itself which lies within the Self will guide you, will make it possible for you to attain Self-Realization in this life itself. That’s what I was saying earlier. This is when the soul calls for help. When the soul calls for help, it sends a message into the cosmos and says – sends the message to \textit{guru}, “Ah, that one is ready! But where is he or she? How to go there? What must be the excuse to go there?” All this gets prepared in action, and all this happens in nature. Outside, you think, “Yes, yes, I am doing something.” No, it is the Divine which is doing it. But yet, still, this action of you doing something will never free you, because that mind is still there, no? If that mind is not changed, how would you make an effort? That action will never help you to be free.

Here, Arjuna is saying to Bhagavān Kṛṣṇa, “Please, help me!” He knows, Arjuna realizes, “I don’t know what is right, what is wrong. You are the Supreme Lord. I trust in You. Even if I want to, my mind is so busy, so distracted, so weak; if You don’t help me, I can’t do anything. You are my best friend, and I see, due to that friendship, I am asking You for help. Destroy this ignorance from me.” Who can destroy that ignorance? Only the \textit{guru}. So now he is taking shelter at the Feet of Kṛṣṇa, “I take refuge in You as a disciple, so please, enlighten me! I give myself to you. If I stay in this state of confusion, it will not lead me anywhere. I realized now through what You said to me about my mind. You have not said something to please me. I wanted to hear something nice from You, but You have not said it. I see the reality that my mind is wrong, and I accept it. This willingness to surrender to the Feet of the Master, that’s what Sharanagati stands for. I take refuge at the Feet of the Master, and when I take refuge, I forget about myself. I give myself completely to You. Do as it pleases You. I am like this diamond. I know I have something great inside, but when you see a diamond, I look just like a normal stone. But in the hands of a diamond polisher, it will start shining.” In the hands of the \textit{guru}, the disciple will start to shine.

Here, all the friendship which was there, disappeared. It’s a different relationship that starts. It’s a \textit{guru} relationship.

He asked Bhagavān to have mercy upon him, to take pity upon him and that he is willing to change something. He is willing to change and accept whatever Bhagavān Kṛṣṇa is saying to him. And Bhagavān knows that he is ready for that. He knows that his heart is ready to plant whatever He is going to plant inside of him. He is surrendered. He is ready to receive. That’s why he said, “I am Your disciple, please take me, and I take refuge at Your Feet.”

We don’t say, “Yes, I take refuge in your hands.” Or, “I take refuge in your arms.” No, He said, “If you want to surrender, you have to humble yourself. You have to make yourself low. Like Arjuna, who has reduced to the lowest state, from being a great warrior to becoming a coward who wanted to run away from the battlefield. He must have great humility to bow down to Bhagavān Kṛṣṇa and take shelter in Him, no? And he did. He did ask Bhagavān, “I want to surrender to You. I am surrendering to You, and I declare that I am Your disciple. I know that I belong to You.” It’s not about the \textit{guru} coming and reminding. The \textit{guru} will make ways that you are reminded, but you also have to have the feeling inside of you, “Yes, I belong to you.” If you don’t have the sense of belonging, then it’s very difficult. Then you don’t belong to anybody. Your \textit{guru} is yours. It’s more than your father, more than your mother, more than anybody, because he is the only one you can claim, “Yes, that’s mine.” Nobody else you can claim that it is yours. But you have to have that feeling. And rhe \textit{guru} will have the same feeling, “You are mine.”

Here, it was not Bhagavān Kṛṣṇa who said to Arjuna, “Take shelter at my feet!” If Bhagavān Kṛṣṇa would have said to him, “Take shelter at my feet”, do you think Arjuna would listen? You know how the mind is rebellious, it would never happen. But it comes from himself. Inside of him, he felt this call, and then he surrendered. This willingness to surrender made him fall down at the Feet of Bhagavān Kṛṣṇa and ask Him, “Please Lord, take me and guide me!”

He was his best friend, He was his cousin, but deep inside he knew that He is different, He has a certain divinity inside of Him, He has the Grace. Probably he was thinking He is a great saint, a great sage. But Arjuna, all of them, all the Pandavas, they considered Him differently. But Arjuna had a deeper relationship with Kṛṣṇa. They were always together. That’s why, when you have the feeling of the \textit{guru} inside of you, that’s why you can say your \textit{guru} belongs to you, because it’s not only one time that you are having this relationship with the \textit{guru}. It is an eternal relationship.

\Verse[2.8]
{% --- 1. The Verse ---
  na hi prapaśyāmi mamāpanudyād \\
  yac chokam ucchoṣaṇam indriyāṇām \\
  avāpya bhūmāv asapatnam ṛddhaṁ \\
  rājyaṁ surāṇām api cādhipatyam
}
{% --- 2. The Word-for-Word (Clean text, no formatting) ---
  \textit{api}---even; 
  \textit{avāpya}---upon attaining; 
  \textit{ṛddhaṁ}---a prosperous; 
  \textit{rājyaṁ}---kingdom; 
  \textit{asapatnam}---unchallenged; 
  \textit{bhūmāv}---on this Earth; 
  \textit{vā}---or; 
  \textit{cādhipatyam}---lordship; 
  \textit{surāṇām}---over the gods; 
  \textit{aham}---I; 
  \textit{na hi}---cannot; 
  \textit{prapaśyāmi}---see; 
  \textit{yac}---what; 
  \textit{apanudyād}---could remove; 
  \textit{idam}---this; 
  \textit{chokam}---grief; 
  \textit{ucchoṣaṇam}---drying up; 
  \textit{mama}---my; 
  \textit{indriyāṇām}---senses
}
{% --- 3. The Translation ---
  Even the attainment of a prosperous, unchallenged kingdom on this Earth—or lordship over the gods—I cannot see what could remove this grief that dries up my senses.
}

Here Arjuna asks, “What am I doing in this battle where I will gain only earthly things? Even if I attain ‘sovereignty of the Gods’, lordship over the Gods, and become like Indra, that will not remove this grief, which is drying up my senses. Please Lord, help me! Remove this grief! Give me everlasting happiness and bliss. That’s what I’m longing for. I’m not longing for material gain and I’m not longing for Heaven. I long to hear You. I long to be out of Maya. I want to be free from this delusion.”

\Verse[2.9]
{% --- 1. The Verse ---
  \hspace*{1em}sañjaya uvāca \\
  evam uktvā hṛṣīkeśaṁ guḍākeśaḥ parantapaḥ \\
  na yotsya iti govindam uktvā tūṣṇīṁ babhūva ha
}
{% --- 2. The Word-for-Word (Clean text, no formatting) ---
  \textit{sañjaya uvāca}---Sañjaya said; 
  \textit{evam}---having thus; 
  \textit{uktvā}---spoken; 
  \textit{hṛṣīkeśam}---to Hṛṣīkeśa; 
  \textit{guḍākeśaḥ}---Arjuna, the conqueror of sleep; 
  \textit{parantapaḥ}---[and] vanquisher of enemies; 
  \textit{uktvā}---said; 
  \textit{iti}---thus; 
  \textit{govindam}---to Govinda; 
  \textit{na yotsya}---\enquote{I will not fight}; 
  \textit{ha}---[and] indeed; 
  \textit{babhūva}---became; 
  \textit{tūṣṇim}---silent
}
{% --- 3. The Translation ---
  \hspace*{1em}Sañjaya said: \\
  Having thus spoken to Hṛṣīkeśa, Arjuna, the conqueror of sleep and vanquisher of enemies, said to Govinda, \enquote{I will not fight} and became silent.
}

Here Arjuna makes his decision, and says, “I’m not ready to fight! I’m happy with my normal, mundane life. I don’t want anything else. I don’t want go deeply into spirituality. I don’t want to discover my true Self. I’m happy living life; waking up in the morning; having my breakfast; going to work; coming back from work; enjoying the TV and sleeping.” Again, next day the same routine. “Then I look forward to the weekend: I will go to the pub, I will drink; I will lose my mind; I won’t even know what happened. I will go completely into delusion; then I will come back home,” – if you come back home, or maybe you will sleep on the road; or wake up in somebody else’s bed or somewhere else, not even knowing what has happened. Sunday, I will relax; Monday, the same routine – and, “I’m happy with that.” This is the delusion in the world. Who wants to change when people have created their own reality of happiness? No one wants to change.

So here, Arjuna says, “I don’t want to fight! I want to listen to you.” That’s why he becomes silent. He is praying to the Lord. Sanjaya says to the king, “After taking refuge at the Lotus Feet of the Lord, after surrendering to the Lotus Feet of the Lord, Arjuna waits for His guidance and instruction. And he becomes silent.” This is what Kṛṣṇa has been waiting for, that Arjuna silence himself a bit and be ready to listen.

Here Sanjaya refers to Kṛṣṇa as Govinda. The truth about God is known only through the Vedas. Kṛṣṇa is Govinda, the Lord of the Vedas. He is the incarnation of the Vedas, the living Vedas. It is through the study of the Vedas that He is to be known. The study of the Vedas has only one aim: to attain the Grace of Nārāyaṇa Kṛṣṇa.

Arjuna has become silent. He has made himself ready. Sanjaya is excited inside and is eagerly waiting to see what will happen next. He says, “Okay, now Arjuna is in silence after expressing himself and going through all this delusion.” Sanjaya himself is getting excited. You know, it’s like when you are watching a nice movie and you’re wondering what will happen next and at that moment they put on a big advertisement and you have to wait. Such is the state of Sanjaya. He has been closely watching this big movie, he sees the climax coming soon and asks, “Ah, what will be next?” He has very deep anticipation; he is waiting for what will come next.

\Verse[2.10]
{% --- 1. The Verse ---
  tam uvāca hṛṣīkeśaḥ prahasann iva bhārata \\
  senayor ubhayor madhye viṣīdantam idaṁ vacaḥ
}
{% --- 2. The Word-for-Word (Clean text, no formatting) ---
  \textit{bhārata}---O descendant of Bharata; 
  \textit{viṣīdantam}---as he was grieving; 
  \textit{madhye}---between; 
  \textit{ubhayoḥ}---the two; 
  \textit{senayoḥ}---armies; 
  \textit{hṛṣīkeśaḥ}---Kṛṣṇa; 
  \textit{iva}---slightly; 
  \textit{prahāsann}---smiling; 
  \textit{uvāca}---spoke; 
  \textit{idam}---the following words; 
  \textit{tam}---to him
}
{% --- 3. The Translation ---
  O descendant of Bharata, as he was grieving between the two armies, Hṛṣīkeśa, slightly smiling, spoke the following words to him.
}

Here, Sanjaya is expressing interest in seeing the Lord smile. The smile of the Lord is saying, “Hey, now you are ready!” It has been Arjuna himself who has requested Lord Kṛṣṇa to bring the chariot between the two armies. Arjuna has first been in the spirit of challenging: but when he sees the two opposing armies, he falls into a depression. He becomes very disturbed and grief-stricken. It is this Arjuna who now asks the Lord to guide him and to remove from him the pain which is overwhelming his mind. So Kṛṣṇa smiles.

When the \textit{guru} sees that the disciple is really ready, he has the same smile! The disciple can ten thousand times ask, ask, ask, ask, and the \textit{guru} will give whatever He has to give. But when the disciple is truly ready, even without asking, all will be given. This is the smile that Lord Kṛṣṇa gives to Arjuna. He says, “Ah! I see that now you are ready.”

When one is confused, there is only one thing one has to do: hold tightly to the Feet of God. This is a natural state, you know? Whenever something goes wrong, who do you call? You say, “Oh my God!” You don’t say, “Oh my mother! “Oh my father!” or “Oh my friend!” You don’t do that! When your mind is battling, when you see that you are confused, when you see that something inside is coming up, the first thing you say is, “Aaaahh! My God!” This is a natural reaction – you just do it; maybe you don’t even realize that you are doing it. But this is when the soul reacts and says, “God, you are the creator! This is a reminder that You are the Ultimate. In every situation, You are the Ultimate! In good situations, You are the Ultimate! In bad situations, You are the Ultimate! And beyond that, You are still the Ultimate!” This is the state that Arjuna is expressing. The Lord is looking at this mental state of Arjuna and laughing. He says, “Yes, My dear, I can now see that you are ready to receive the greater knowledge of the Self. You are ready to fight and transform yourself.”

\textbf{2016}---Sanjaya carried on explaining. Here, Sanjaya is talking to the blind king, and expressing, “Hṛṣīkeśa, slightly smiling, spoke the following words to him.” He didn’t say, “You stupid!” Paf, paf! No, with a smile!

Sanjaya is saying, “O descendant of Bharata, as he was grieving between the two armies.” When the chariot was in between the two armies, Bhagavān Kṛṣṇa was fully relaxed, He was not stressed, looking at His stressed devotees. He was smiling. Bhagavān said, “Hey, now you are ready.” The changes started to happen, and Bhagavān was very happy when Arjuna surrendered to Him. He said, “Finally! I was waiting for that since so long! Finally, that moment has arisen inside of you. That moment has arrived, I will help you. I see that you are ready to receive everything. I will give you everything. I will reveal everything to you. I will reveal to you who you are.” That’s why He smiled, and said, “Now you are ready. Before, you were not ready.”

So the same, when you are ready, then everything takes its place. Arjuna was in a state of challenge when he saw in between the two armies, he saw his family, friends, what he is letting go, what will happen, what will fall and what will be lost. He became very distressed. Through the guidance of the Lord, he asked that this pain be removed from his mind. He said, “You know better how to deal with that. You are the doctor, You know the medicine. I can take all the medicine and drink, it doesn’t mean that I will get cured, but You know best.”

This is the same; when the \textit{guru} sees that the disciple is ready, the \textit{guru} smiles. Here Bhagavān is very happy that this confusion which arose from the mind of Arjuna has changed, this judgment has changed, this total incapacity of making the right decision has changed. He is on the right track. Now Bhagavān Kṛṣṇa speaks.

\Verse[2.11]
{% --- 1. The Verse ---
  \hspace*{1em}śrī-bhagavān uvāca \\
  aśocyān anvaśocas tvaṁ prajñā-vādāṁś ca bhāṣase \\
  gatāsūn agatāsūṁś ca nānuśocanti paṇḍitāḥ
}
{% --- 2. The Word-for-Word (Clean text, no formatting) ---
  \textit{tvaṁ}---you; 
  \textit{anvaśocas}---grieved [for]; 
  \textit{aśocyān}---those who should not be grieved for; 
  \textit{ca}---yet; 
  \textit{tvaṁ}---you; 
  \textit{bhāṣase}---speak; 
  \textit{prajñā-vādān}---words of wisdom; 
  \textit{paṇḍitāḥ}---the learned; 
  \textit{nānuśocanti}---lament; 
  \textit{ca}---neither; 
  \textit{gatāsūn}---for the dead; 
  \textit{ca}---nor; 
  \textit{agatāsūṁś}---for the living
}
{% --- 3. The Translation ---
  \hspace*{1em}Bhagavān Kṛṣṇa said: \\
  You are grieving for those who should not be grieved for, yet you speak words of wisdom. The truly wise ones lament neither for the dead nor for the living.
}

Now is the beginning of the Gita. The real Gita starts now with the smile of the Lord. All that came before was just a preparation for this.

The Lord says, “ You are grieving for those who should not be grieved for, yet you speak words of wisdom.” You see, Bhagavān Kṛṣṇa has been watching and observing Arjuna all the time. The Great Observer always observes life and sees life. He observes how life is going on. Your true Self, the \textit{ātmā}, is always observing life from within you. Kṛṣṇa says, “You grieve, ‘yet you speak words of wisdom’. You know about life, you have the right feeling inside but when you speak, you say something stupid.” You have the deep feeling inside that you should not do something! But you do it anyway! And then afterwards you say, “Oops, terrible! I did it again!” Learn this great wisdom. You have the wisdom, you have the knowledge, you have the feeling inside you.

“The truly wise ones lament neither for the dead nor for the living.” Kṛṣṇa says, “There is no beginning and there is no end. There is no birth, there is no death.” Kṛṣṇa is going slowly with Arjuna, saying, “Let Me give him the knowledge step by step.” He is gradually taking Arjuna from the state he is in before. The utmost thought in Arjuna’s mind, at that time, has been about the destruction of the family, the destruction of the race and the inevitable consequences of the war. In Chapter 1, he says that the war will bring confusion; it will cause the downfall of the castes; and the ancestors will go to hell. Kṛṣṇa has said, “You are wise. Mourn not.” 

If you want true wisdom, know that there is no birth and no death. You are the \textit{ātmā}, you are eternal. For the \textit{ātmā}, there is no beginning, there is no end. So why do you mourn? You have come here to attain a higher reality; you have come here to do a higher work than what you realize. Kṛṣṇa tells Arjuna that the wise never grieve in this way. The one who has the knowledge of the Self, the realized one, the Jivan Mukta knows that the \textit{ātmā} is eternal and that there is no point in mourning or grieving. The wise person, the paṇḍitāḥ, knows that God is the embodiment of true knowledge, Satchitananda. The wise person knows that it is only Him who is abiding in everything. He is the Self of all. He is the indestructible Absolute. Whatever is created – the body or whatever you see in the outer world – is not permanent. He says, “Do ‘not lament neither for the dead nor for the living.’ Whatever is created, will be destroyed.” Everything in this world keeps changing. The only one who is not changing is God. The body is temporary, it goes through changes and dies. It doesn’t stay.

The association, the link between the body and the soul, is like a dream. Sometimes you dream of something and in that dream, everything appears so real: you are running, you are very active and very busy in that dream; you feel it! But when the dream ends, you wake up, and realize, “Oh, it was a dream.” When you are in that dream, you don’t see reality. But when you are out of that dream, you are back in reality. Kṛṣṇa is talking about birth and death, about living and not living. This is similar to dreaming. When you are in a dream, everything is imaginary; but when you come back to reality, you realize that it was just a dream. So awake! An awakened man, enlightened man knows that a dream is only a dream and he stays out of that dream.

Kṛṣṇa says, “So Arjuna, why are you in these circumstances in grief? You are wise, you should not be in grief. In your words when you speak, you are grieving. But there is much Truth inside you; there is much wisdom inside you. You are not just stupidly saying all these words. You have great compassion for others; you care deeply for others. But do you really want to help others? First help yourself! If you help yourself and become strong, then it will be much easier to help other people.”

\textbf{2016}---The Supreme Lord says, “You are grieving for those who should not be grieved for.” You are crying for those whom you should not cry for. It’s a waste of time you cry for them.

“Yet you speak words of wisdom.” You see Bhagavān was analyzing him. When he was talking nonstop, Bhagavān didn’t say anything. Bhagavān Kṛṣṇa kept quiet. What He was doing, He was analyzing each word of Arjuna. He saw, “Yes, you are deluded, but there is something greater inside of you, and that greatness wants to awaken, and through the words, what you are asking for, it is very – you have lots of wisdom.”

Here Bhagavān said, “The truly wise ones lament neither for the dead nor for the living.” Those who have reached the point of enlightenment, they don’t cry or enjoy for the living, nor for the dead. They are free. Bhagavān Kṛṣṇa was the witness to Arjuna’s drama. He was watching, observing all the time. Arjuna was conscious of his own state. For sure, inside, he was asking himself, “What is wrong with me? Why am I in that state?” But Bhagavān had to act as a \textit{guru}. Not only for Arjuna’s sake, but for everybody’s sake, for the sake of humanity, He has to create that drama.

Do you think that when He went as an ambassador of peace, He could not have changed everything? He could! In a fraction of a second He could have annihilated everything from that court, but He didn’t. He had a different agenda. A long-term agenda.

Here, Bhagavān is saying that the enlightened one doesn’t mourn for either the living or for the dead. He is saying one cannot experience one’s own death. For if one could, then one would be dead. If one would experience death, they would not be dead, no? They would still be alive. You would still have a certain awareness of it. Sometimes you see people, they die and come back, a near-death experience.

How would it be possible? Death can’t experience life, and life can’t experience death. One is just the witness of life. You are the witness of birth and death. “Can I be that witness without experiencing it?” Yes, you can. That’s what your spiritual \textit{sādhana} stands for, no? It’s an inside feeling.

“The truly wise ones lament neither for the dead nor for the living.” There is no beginning for the \textit{ātmā}, and there is no end for the \textit{ātmā}. The \textit{ātmā} is eternal. Have that knowledge of this eternity of the \textit{ātmā}. Then you don’t need to sit and cry. This Self is the witness which doesn’t have any birth or death. All phenomena, like hunger, thirst, emotional drama, sorrow, pain, pleasure, all these qualities are only attached to the body, but not to the Self. The Self observes only. He doesn’t cry. Your \textit{ātmā} doesn’t cry and say, “Oh, my goodness, that \textit{ātmā} has left this body! Terrible!” No. The body cries, your mind cries due to the attachment that you have, but the \textit{ātmā} doesn’t cry. That Self is free. Here Bhagavān is saying this is very unwise; to sit and cry like that, for the Self is ever free. You are here as a witness, the witness of everything. Observe and see.

Everybody has a certain perception. When you are walking, you have a certain perception in the dreaming state. You have a certain perception when you are in a deep sleep state. When this perception dies, everything dies. When you are awake again, this perception becomes alive again. When you are in a deep state, there is difference; each step has a different perception of itself. When you go into deep sleep, it’s a state of – absence of perception, a deep sleeping state. When you awaken in the morning, you are again in an awakened perception. “That Self which is inside of you, dear Arjuna, don’t grieve for anything.” What you do when you are in a deep state, you don’t remember. Only when you are awake and do something, then you remember. When the mind becomes active, then you start remembering.

So, Kṛṣṇa said, “O Arjuna, why are you in these circumstances of grief? You are wise, you should not be grieving, because in the words that you speak, you are very knowledgeable. You have much inside of you. You have much love inside of you. You have great compassion for others. You care deeply for others. But if you truly have this compassion for others and you truly want to help them, first start by helping yourself. Help yourself to become strong, that will be a great help to others. Have this knowledge of - remind yourself constantly, that you are the soul, the \textit{ātmā}. Have this true understanding of the Self.”

This is what Bhagavān Kṛṣṇa is reminding Arjuna of, that you should not – “You grieve for those who should not be grieved for, yet you speak words of wisdom.” A person has eyes, but doesn’t see. You are crying, you are grieving, but you have that wisdom. If you have eyes to see that you see, if you have ears to hear that you hear, no? So, a person who is blind, they don’t perceive the object. But someone who has seen the object, they know how it is. Arjuna had this realization inside of him, yet he is acting as being blind. You know about the Self, yet you choose to be in ignorance. Here Bhagavān is saying, everybody knows that they are not the body, everybody knows that they are not the mind, they are not the intellect, but they still identify themselves. This is due to a lack of true knowledge, a lack of true understanding. You have got this knowledge through books, you have read wonderful books which have given you this insight, you have heard many talking about it, but it has just stayed words which are pleasing to you, but it doesn’t have any effect upon you.

You are still running after the external pleasures. You are still running after things which are perishable. You still make yourself weak, and with time, you are finishing. This is the sorrow of life, to show how life is short here, when you identify yourself with the body and the mind. You grieve, you cry, thinking that you are losing something. One after the other life you keep jumping. Then, in life, you consider yourself very important. Then you spend your time doing futile things. Then you try to find true happiness in that. It’s very difficult, because it doesn’t go together. “Sitting and crying will not lead you anywhere, Arjuna. You have to choose now. You want to cry and pity yourself, or you want to awaken and make yourself ready? Both are given to you on a plate. Choose wisely.”

\Verse[2.12]
{% --- 1. The Verse ---
  na tv evāhaṁ jātu nāsaṁ na tvaṁ neme janādhipāḥ \\
  na caiva na bhaviṣyāmaḥ sarve vayam ataḥ para
}
{% --- 2. The Word-for-Word (Clean text, no formatting) ---
  \textit{tu}---but; 
  \textit{eva}---certainly; 
  \textit{na jātu}---[there was] never [a time when]; 
  \textit{aham}---I; 
  \textit{na āsam}---didn't exist; 
  \textit{na tvam}---nor you; 
  \textit{na ime}---nor these; 
  \textit{janādhipāḥ}---kings; 
  \textit{na ca eva}---and certainly not; 
  \textit{sarve}---[will there be a time where] all; 
  \textit{vayam}---[of] us; 
  \textit{na bhaviṣyāmaḥ}---will not exist
}
{% --- 3. The Translation ---
  Surely there has never been a time when I didn't exist, nor you, nor any of these kings, nor will there ever be a time when we all don't exist.
}

“Surely there has never been a time...” This is so beautiful. Here Bhagavān Kṛṣṇa reveals to Arjuna the quality of the eternity of the soul. He says, “There was never a time when you nor I did not exist. We all existed even before our birth, before we manifested in these bodies. And we will carry on existing even after these bodies disappear. Even though the body dies, the soul is eternal, it will carry on living. So Arjuna, you should not grieve for these relatives or fear their destruction.”

The soul is eternal: one has to come here; one does one’s duty, one’s \textit{dharma}; and then the body dies. But the soul remains eternal. Before this lifetime here, you were somewhere else in another body; but now you are here in this body. You will continue in this cycle of birth and death until you finally attain the Lord Himself. Kṛṣṇa is pacifying Arjuna saying, “They all existed before. Even you. Even Me, I always existed. We are just playing this game, this drama. When this body is finished, we will still carry on living afterwards. So why do you grieve for that? Awake! Realize yourself! The nature of the soul is eternal.”

\textbf{2016}---This is very important to understand, what Bhagavān is saying here right now. He is giving this deep knowledge to Arjuna that, “There was never a time that I didn’t exist. Don’t see Me as what you see on the outside here, and think that, ‘Okay, I am just talking to you to make you happy.’ No, I ever exist when nothing exists. There was never a time that I did not exist, nor you! As I existed before everything, you were also with me.” This is the relationship that we have with God. We were always with Him, and He is always with us, whether you want it or not. You can be the greatest atheist person, you can say God doesn’t exist, it doesn’t bother Him, because deep inside He knows that God is with each one all the time, maybe not in the same concept as everybody understands of Him, but in a different concept.

Here Bhagavān is saying to Arjuna, “Nor any of these kings,” they were always there, these people, and they will always be there. There is no disappearance, there is no merging into something, so that you are finished. No, you don’t. Here Bhagavān is saying, “You have an identity. You will always have it. You ever exist with Him. You are never separate from Him. We all existed before our birth, before we manifested here in these bodies, and we will carry on existing even way after these bodies have disappeared. That \textit{ātmā} will always live. Even if everything gets destroyed, even if the whole universe, when Brahma closes His eyes and opens them again, there is always Pralaya, everything is happening, everything changes, no? The whole world dissolves itself, and a new creation starts. Here Bhagavān is saying that even if that happens, you always exist, and you are not subjected to disease, this \textit{ātmā} is not subjected to decay or death. The cycle which the body goes through, it goes through different shapes, sizes, forms, colors. But these are all changes. One day you are in white skin, tomorrow you are in black skin, tomorrow you are in red skin, you are in yellow skin, brown skin, whatever, these are all changes. As you go through changes, you see, the body itself, you are born, and from being a baby, being small, you grow up to be big, big how you all are. You age, you die, and it's finished for that body, finished for that mind. Another life again. So again, when you die, the body goes back into this element. The soul is ever - is eternal. The soul will not say, “Okay, I will wait when this body comes back again,” no? Will the body come back again? Recycle? Well, it goes through recycling, because it’s made of these elements, no? The five elements. And the five elements go back into nature. But the Self is not touched, is unchanged, immutable, indestructible, and it is eternal. It is not created. If something is not created, it means it can’t die. Everything that is created has a beginning and automatically has an end. Here Bhagavān is saying that, “No, your soul remains, is eternal, was eternal and will always be eternal.” You shall dwell continuously with Him. Those who attain the Lord, who attain the Grace of the Lord, dwell with Him eternally.

So awaken, my dear soul. Realize yourself. You are the \textit{ātmā}. Stop being in the sleeping state. Because in a sleeping state, you will not realize anything. You have to be awake to realize yourself.

\Verse[2.13]
{% --- 1. The Verse ---
  dehino ’smin yathā dehe kaumāraṁ yauvanaṁ jarā \\
  tathā dehāntara-prāptir dhīras tatra na muhyati
}
{% --- 2. The Word-for-Word (Clean text, no formatting) ---
  \textit{yathā}---just as; 
  \textit{kaumāram}---childhood; 
  \textit{yauvanam}---youth; 
  \textit{jarā}---[and] old age; 
  \textit{asmin}---[take place] in this; 
  \textit{dehe}---body; 
  \textit{tathā}---similarly; 
  \textit{deha-antara-prāptiḥ}---[there occurs] the attainment of another body; 
  \textit{dehinaḥ}---for the embodied \textit{ātmā}; 
  \textit{dhīraḥ}---the wise man; 
  \textit{na muhyati}---is not bewildered; 
  \textit{tatra}---about this
}
{% --- 3. The Translation ---
  Just as childhood, youth, and old age take place in the body, the embodied \textit{ātmā} attains another body at death. A wise person is not bewildered about this.
}

Kṛṣṇa says, “It is unwise to suffer, thinking that by changing one’s body, everything is lost. Why are you grieving? Why are you sad about this? Childhood, youth, and old age don’t belong to the soul.” When are you growing up, what is growing? The body. Once you were a baby; maybe now you are in youth; later you will be in middle age; and later on, you will become old…some of you are already old now. But these changes are not for the soul. Your \textit{ātmā} is eternal! The \textit{ātmā} is still the same. For all which is created, change is only on the outside. The body changes. Nature, Prakriti changes. The soul itself never changes.

When the body dies, the soul goes and takes another body. The one who is unaware of this truth, the unwise, grieve over this transition from one body to another. The wise never do so. The one who is wise, “the self-composed man”, the one who is longing for Realisation, the realized one, sees the world completely differently. Kṛṣṇa says, “Arjuna, you are realized. Don’t see the world in a low state of mind. Don’t be blinded by how everybody else sees the world. Let this knowledge which is inside of you awaken. See that the soul is eternal. Even if the soul moves from body to body, this eternity of the soul – that’s what is the Reality.”

\Verse[2.14]
{% --- 1. The Verse ---
  mātrā-sparśās tu kaunteya śītoṣṇa-sukha-duḥkha-dāḥ \\
  āgamāpāyino ’nityās tāṁs titikṣasva bhārata
}
{% --- 2. The Word-for-Word (Clean text, no formatting) ---
  \textit{mātrā-sparśāḥ}---contact of the senses [with their objects]; 
  \textit{tu}---indeed; 
  \textit{kaunteya}---O son of Kunti; 
  \textit{śīta-uṣṇa-sukha-duḥkha-dāḥ}---give rise to cold and heat, pleasure and pain; 
  \textit{āgama-apāyinaḥ}---they come and go; 
  \textit{anityāḥ}---[and are] impermanent; 
  \textit{tāṁs}---them; 
  \textit{titikṣasva}---endure; 
  \textit{bhārata}---O descendant of Bharata
}
{% --- 3. The Translation ---
  The contact of the senses with their objects, O Kaunteya, gives rise to feelings of cold and heat, pleasure and pain. They come and go and are impermanent, so endure them without being disturbed, O descendant of Bharata.
}

Mātrā sparśās means ‘The contact of the senses with their objects’. Whatever comes into contact with the senses, into contact with the mind – heat, sound, touch, colour, taste and smell – these are the products of the senses. Heat and cold, pleasure and pain, are always part of our experience. If one gives attention to the outside, to the sense organs, and the sense objects, one becomes bound by the opposites of duality, like cold and heat, pleasure and pain. Whoever is focused on living only in the external world, will attain only limited things. One will not be aware of the greatness, the eternity of the soul. This duality of experience, also brings joy, grief and so on. This is not for realized people. When you realize that the object is just passing by and changes, you should not allow the mind to dwell on this. It is unhealthy to dwell on this.

Describing the contact between the senses and the object which naturally comes and goes and the limitations of Nature, Prakriti, Kṛṣṇa asks Arjuna to ignore this: “this will only bring pain.” He advises Arjuna to not concentrate on them, to not rejoice in them or grieve over them. Know that these experiences are there, but don’t go into them. If you start feeding them, if you feed the negative qualities, if you go into them – you are trapped. You are trapped in the drama of the mind. And the mind will make you dance to its tune. And once you are dancing to the tune of Maya, all is finished!

Here Lord Kṛṣṇa tells Arjuna, “Overcome pleasure and pain. Go out of this drama of duality and remind yourself that you are eternal.” Again and again He is repeating the same thing: the eternity of the soul. Sometimes you have to repeat something ten times for somebody to get it. The mind is like this. Until it accustoms itself to a certain reality, that reality is easily forgotten and the mind jumps around.

In Chapter 1, again and again Arjuna is reminding Kṛṣṇa of his grief and sadness and so on -- he is finding many excuses not to fight. The Lord is doing the same thing with Arjuna, but in a reverse way, reminding him that, “All that you have been talking about is not real, is not the reality. It is just passing by. You are here to witness this. You are the great witness.” It is only when you realize that you are the Great Observer that you can observe life and not be trapped in the drama of it. When you are apart from it, then you observe. When you are in your \textit{sādhana}, when you are sitting in meditation, who is perceiving everything? When you are doing your Atma Kriya Yoga, of course, the mind is perceiving the physical state; but when you look deeper inside, on the inside there is somebody else watching. Who is watching? That is the great Self. The great Self is always watching; the Great Observer is always watching and observing life, how life is going. Now it’s observing and hearing in this body; before it was observing and hearing in another body. And those who have not yet attained the Grace of the Lord, will later be in another body. Of course you will go through an evolution.

\Verse[2.15]
{% --- 1. The Verse ---
  yaṁ hi na vyathayanty ete puruṣaṁ puruṣarṣabha \\
  sama-duḥkha-sukhaṁ dhīraṁ so ’mṛtatvāya kalpate
}
{% --- 2. The Word-for-Word (Clean text, no formatting) ---
  \textit{puruṣa-ṛṣabha}---O best of men; 
  \textit{dhīram}---[the] wise; 
  \textit{puruṣam}---person; 
  \textit{yam}---whom; 
  \textit{ete}---these; 
  \textit{na vyathayanti}---do not affect; 
  \textit{sama-duḥkha-sukham}---who equal in pain and pleasure; 
  \textit{saḥ}---he; 
  \textit{hi}---indeed; 
  \textit{kalpate}---is worthy; 
  \textit{amṛtatvāya}---of immortality
}
{% --- 3. The Translation ---
  O best of men, the wise person who is unaffected by these, and to whom pain and pleasure are the same, is indeed worthy of immortality.
}

Kṛṣṇa addresses Arjuna as  best of men. He is reminding Arjuna, “You are like a lion! You are strong! You are not weak. You are powerful! Be firm and wise.”

“The man who these do not trouble”, yaṁ hi na vyathayanty ete, means the one who is not troubled by pleasure and pain. The best and most excellent of men are beyond pleasure and pain. It is easy for them to rise above these opposite qualities, to rise above these dual experiences.

The realized soul doesn’t go into the drama of life. He sees pleasure and suffering equally. The one who makes himself ready becomes ripe. He becomes aware to receive the knowledge of the Self.

Kṛṣṇa is pushing Arjuna. He says, “Don’t just talk about what is right and not right. Make an effort to change. See it from another point of view, so you can rise above this limitation. Make an effort to see that these objects of the senses, these objects of pain and pleasure are just limited. Today they are here; tomorrow they are not here. If you can change this feeling in the mind, if you can rise above it, from then on, you will see everything as equal. Only then will the immortality of the soul be revealed to you. But if you are attached to the duality, you are not ready for it.” You see, here it is very important to know that people may talk a lot about Realisation, but if they are not realized, what good is that?

Kṛṣṇa says, “Make yourself ready. You have been given a life here to make yourself ready to receive the Grace of God.” Kṛṣṇa doesn’t only talk about the immortality of the Self: He also talks about attaining the Grace of God. But in this verse, He is not talking about attaining the unlimitedness of God; rather He says, “Make yourself apt for immortality.” However, Arjuna is so focused on the outside reality, that it is difficult for him to think about anything else.

\textbf{2016}---“The wise person who is unaffected by these, and to whom pain and pleasure are the same“ You are wise, you are strong, in your nature, you are equal to pleasure and suffering. You are immortal. You have this immortality. You are eternal.

Everything in life is subject to change. Everything comes and goes, pleasant, unpleasant. One gets things according to one’s karma, but for the \textit{ātmā}, there is no pleasure, there is no discomfort. The Self is eternal. It is not material. Therefore, understand the completeness of the Self. By the mind, by the body or by the intellect, it is impossible. Bear in mind that these pairs of opposites are only in the mind. The soul which is realized, those \textit{ātmā}s, those people who are realized, they don’t go into the drama of life. They observe. They see this equalness everywhere. They are fully absorbed into the knowledge of the Self. They know that this is just a dream. Why mingle in that dream?

Bhagavān is pushing Arjuna, saying, “Don’t just talk about what is right and what is not right. Make an effort to change!” He is addressing him as the best of men, meaning that you are very strong, and due to this - you are very wise, but if you just sit and don’t make any effort, it’s not possible. You will not get to know the \textit{ātmā}. The \textit{ātmā} will not just reveal itself like that. Only when you do your \textit{sādhana}, then you can rise above that. Make yourself ready. You have been given this life, and this life is to make yourself ready to receive the Grace of God.

Here He is not just saying “Realize yourself as being immortal,” that’s it. No, He said, “Realize, make yourself ready.” Which means, make yourself aware that you are the \textit{ātmā} and you are immortal, you are eternal, and that this eternal \textit{ātmā} has a relationship to God. He is not saying, “Okay, mind, understand that you are immortal.” No, aim for the Grace of God so that where you were from the beginning, you can take back the same rightful place. You can be next to Him, serving Him directly. Once you were doing it! But because you have identified yourself so much with the limitation, you have limited yourself. Here Bhagavān is saying, “No, stop that! Aim for the Grace of the Lord. Aim for the Grace of the \textit{guru}. That’s the only thing that will free you.”


\Verse[2.16]
{% --- 1. The Verse ---
  nāsato vidyate bhāvo nābhāvo vidyate sataḥ \\
  ubhayor api dṛṣṭo ’ntas tv anayos tattva-darśibhiḥ
}
{% --- 2. The Word-for-Word (Clean text, no formatting) ---
  \textit{vidyate}---there is; 
  \textit{na}---no; 
  \textit{bhāvaḥ}---existence; 
  \textit{asataḥ}---of the unreal; 
  \textit{vidyate}---there is; 
  \textit{na}---not; 
  \textit{abhāvaḥ}---non-existence; 
  \textit{sataḥ}---of the real; 
  \textit{api}---surely; 
  \textit{antaḥ}---[this is the] conclusion; 
  \textit{anayoḥ ubhayoh}---regarding the these two; 
  \textit{tattva-darśibhiḥ}---by those who have realized the Truth; 
  \textit{dṛṣṭaḥ}---seen; 
  \textit{tu}---indeed
}
{% --- 3. The Translation ---
  There is no existence of the unreal, and there is no non-existence of the real. This is the conclusion regarding these two that is perceived by those who have realized the Truth.
}

This means that everything which is limited, will be destroyed. And what is not limited, the \textit{ātmā}, cannot be destroyed; it is ever-existent. You will realize the eternity of the soul only when you rise above the duality of the real and the unreal. Why would non-existence exist? Non-existence is not real. By non-existence, Kṛṣṇa is referring to the body and referring to all the people on the battlefield. He says, “These people are not real. They appear to exist, but they are not real. The only reality behind all this is the Self. Look beyond this duality. Perceive that which is beyond what your mind is showing you. The soul doesn’t go through these changes; in all circumstances, it never ceases to exist.” It is unchangeable.

Arjuna is feeling much grief that he will have to attack Bhishma Pita. His grief is not so much for the other Kauravas. The most important thing to him is his relationship with his great uncle. He fears that if he kills him, what will happen to his soul? In the first chapter, Arjuna says, “The ancestors will go to Hell.” So he fears that the soul of Bhishma, who is a great being, will go to Hell. Kṛṣṇa tells him, “No. Bhishma Pita is as pure as the river Ganga. The purpose of his existence here is to do what I have commanded him to do. Nothing can exist without Me willing it that it exists. Everything exists because I will that it exists. But you have to see that it is My Will. If you don’t see that it is My Will, you are stuck in the real and the unreal; and you can’t distinguish between the real and the unreal. The saints and the seers, the wise ones have always reached this conclusion: everything else outside changes; but everything comes from only One Reality, which is God. There is only One Reality, My Reality, nothing else. This is the Truth. This is the essential Truth.”

\Verse[2.17]
{% --- 1. The Verse ---
  avināśi tu tad viddhi yena sarvam idaṁ tatam \\
  vināśam avyayasyāsya na kaścit kartum arhati
}
{% --- 2. The Word-for-Word (Clean text, no formatting) ---
  \textit{viddhi}---know; 
  \textit{tat}---that; 
  \textit{yena}---by which; 
  \textit{idam sarvam}---this entire body; 
  \textit{tatam}---is pervaded; 
  \textit{tu}---[to be] indeed; 
  \textit{avināśi}---indestructible; 
  \textit{na kaścit}---no one; 
  \textit{arhati}---can; 
  \textit{kartum}---cause; 
  \textit{vināśam}---the destruction; 
  \textit{avyayasya}---of the imperishable; 
  \textit{asya}---of this [\textit{ātmā}]
}
{% --- 3. The Translation ---
  Know that which pervades this entire body to be indestructible. None can cause the destruction of the imperishable \textit{ātmā}.
}

Whatever is perishable, like the body, the senses, the mind, all these objects of enjoyment, and pleasure – all these are under the control of Prakriti. And Prakriti is under the control of the Spirit, which means under the control of the Consciousness, the life principle.

“Know that which pervades this entire body to be indestructible.” All matter, everything which is created, comes from the non-created: all emerges from only One Reality, God, the Immortal One, the Supreme Spirit. Everything comes from Him. Everything is just an extension of Him. The one who is realized in this immortality, attains the true Self. He sees that the same Nārāyaṇa Kṛṣṇa that is outside, is also inside. And not only in himself, but everywhere, in everything. The seers, the \textit{yogīs}, and the saints know and see this Truth. They know that everything which is created, everything which is on the outside, is an extension of only Him, and it is only Him that is present in everything.

Kṛṣṇa says to Arjuna, “None can cause the destruction of the imperishable \textit{ātmā}.” How can you kill the \textit{ātmā}? You can’t. At that moment, Arjuna is very worried about what will happen to the \textit{ātmā} of Bhishma. Will he go to Hell? Kṛṣṇa says, “No. How can he go to Hell? He can’t.” It’s not because of his actions on the outside that he will go to Hell: his essence is only the Supreme Spirit. The soul, which is under the control of the Immortal Spirit, how can it be destroyed? It is not a material object. The soul is not something material that you can cut with a sword.

\Verse[2.18]
{% --- 1. The Verse ---
  antavanta ime dehā nityasyoktāḥ śarīriṇaḥ \\
  anāśino ’prameyasya tasmād yudhyasva bhārata
}
{% --- 2. The Word-for-Word (Clean text, no formatting) ---
  \textit{ime}---these; 
  \textit{dehāḥ}---bodies; 
  \textit{nityasya}---of the eternal; 
  \textit{anāśinaḥ}---imperishable; 
  \textit{aprameyasya}---incomprehensible; 
  \textit{śarīriṇaḥ}---\textit{ātmā} dwelling in the body; 
  \textit{uktāḥ}---are said to be; 
  \textit{antavantaḥ}---perishable; 
  \textit{tasmāt}---therefore; 
  \textit{yudhyasva}---fight; 
  \textit{bhārata}---O descendant of Bharata (Arjuna)
}
{% --- 3. The Translation ---
  The bodies of the eternal, imperishable, and incomprehensible \textit{ātmā} dwelling in the body are said to be perishable. Therefore fight, O descendant of Bharata.
}

“The bodies of the eternal, imperishable, and incomprehensible \textit{ātmā} dwelling in the body are said to be perishable.” Kṛṣṇa says, “All bodies, which have been created from matter, from the five elements, ‘have an end’.” Everything which is matter, will return to matter. Here the Lord is saying that the mind, the senses, and the physical body are all perishable. Whatever you see outside is perishable, has its limit and will have its end. But you are here to use these limitations to do your duty, to do your \textit{dharma}. This infinite, the \textit{ātmā}, which is present inside of you is the Reality. He says, “Don’t be in ignorance. Know that you are eternal. Know that you are indestructible. So fight!”

Here Kṛṣṇa is not only talking to Arjuna, He is talking to everybody, “Don’t go into your weakness, My dear children. It is Me who is seated within you. Look at life. Face it! Be strong! Know that I am with you. And know that we have this deep, eternal relationship together, that we are always next to each other. You are not alone because I, the Supreme Spirit, am ever-present with you.” Don’t look at the outside and become miserable! Don’t look at the effort that you are putting into your \textit{sādhana}; just do it! Do your maximum! Don’t go into your weakness when you do your \textit{sādhana}, but be in your power! Only then will you realize your Self. Only then will God reveal Himself to you. Don’t look at the outside world. Don’t listen to ‘x, y, z’. Go with your feeling. Open your heart and be powerful! Be strong! So Kṛṣṇa says, “Arjuna, don’t show any hesitation or unwillingness to engage yourself in the fight.”

\textbf{2016}---“The bodies of the eternal, imperishable, and incomprehensible \textit{ātmā} dwelling in the body are said to be perishable.” The body has an end. As much as you can think, “Yes, let me beautify myself by putting so much makeup upon me, by painting my eyes big, big,” it will end. Whether you want it or not, you will die. Today or tomorrow you will die. This body I am talking about. But you, as the \textit{ātmā}, will always be.

“The bodies of the eternal, imperishable, and incomprehensible \textit{ātmā} dwelling in the body are said to be perishable. Therefore, fight, O descendant of Bharata.” Here He said, everything that you see, all this drama that you have just gone through, all this will finish one day. Whether you want it or not, you are getting old, gray hairs are appearing on your head, and death is near. The body goes through this transformation itself, when you look at life: birth, you are growing, aging, aching, diseases, decaying, death, you disappear, the body is gone. Like I said, or you become a saint, then you enter the relic room. There is a chance for that body somewhere also. But even that is not eternal. One day, it will all disappear. But this \textit{ātmā}, which is inside of you, is avinashi, indestructible, nothing can destroy it. It doesn’t depend on any changes. Like the body, which is dependent on changes, like the mind, which is always dependent on changes, the thoughts, which are dependent on changes, all these, which are external, depend on changes. Everything which depends on change disappears after some time. But the \textit{ātmā} doesn’t depend on anything apart from the Grace of the Lord. That’s why the \textit{ātmā} keeps coming to attain that Grace, nothing else. You think that, “Yes, I have incarnated, it’s beautiful.” What happens afterwards? How long would you be in that cycle of birth and death? How long would you enjoy this change happening?

All that is in nature keeps changing. Everything which you perceive with your eyes keeps changing. Transformation happens, and everything is destroyed after some time. But the truth, which is the soul, the Self, doesn’t get destroyed. It just takes its rightful place next to the Lord, serving Him.

Here Bhagavān is saying, “Don’t look at the outside and become miserable. Don’t look at the effort that you are putting in your \textit{sādhana}, “Oh, my God, it’s so far away this realization! This Swami is always talking about realization. He is always talking about attaining the Grace, it’s so far away there. Distance, terrible.” Your soul is worried already about how to reach there, then how would you reach it? Like I said, everything can be given. Look at Arjuna, he is deep in his weakness, yet Bhagavān Himself is giving him that Grace and making it easy for him. Same thing for you all, you see. Bhagavān, He has given this Bhagavad Gita not only for Arjuna, but He has given it to all the souls. He said, “Yes, I am giving it to Arjuna as a drama which I have made, but in reality it is for each soul to have the true knowledge.”

Imagine Bhagavān Kṛṣṇa spoke that five thousand years ago, but still today we are singing the glory of the Gita! This is eternal. He said, “Don’t look at the effort that you are putting in your \textit{sādhana}.” Because when you look at the effort you are putting in your \textit{sādhana}, you will always discourage yourself. Don’t! Look at the goal, you will see that it is nearer than your eyelashes themselves.

God doesn’t reveal Himself until you are ready. The \textit{guru} doesn’t give the Grace if you are not ready. Here Bhagavān Kṛṣṇa is saying to Arjuna, “Don’t concentrate on the limitation. Don’t limit yourself! If you choose to limit yourself, you will be unhappy. You will see that it is very difficult. You will see that it is far away, when in reality it is not that far. Don’t create this illusion of distance when it is the dearest and the nearest to you.”

\Verse[2.19]
{% --- 1. The Verse ---
  ya enaṁ vetti hantāraṁ yaś cainaṁ manyate hatam \\
  ubhau tau na vijānīto nāyaṁ hanti na hanyate
}
{% --- 2. The Word-for-Word (Clean text, no formatting) ---
  \textit{yaḥ}---one who; 
  \textit{vetti}---believes; 
  \textit{enam}---this [\textit{ātmā}]; 
  \textit{hantāram}---to be the killer; 
  \textit{ca}---and; 
  \textit{yaḥ}---one who; 
  \textit{manyate}---thinks; 
  \textit{enam}---this [\textit{ātmā}]; 
  \textit{hatam}---[can be] killed; 
  \textit{ubhau}---both; 
  \textit{na}---[do] not; 
  \textit{vijānītaḥ}---know; 
  \textit{ayam}---this [\textit{ātmā}]; 
  \textit{na hanti}---does not slay; 
  \textit{na hanyate}---nor is slain
}
{% --- 3. The Translation ---
  One who believes the \textit{ātmā} to be the killer, and one who thinks it can be killed — both do not know; for the \textit{ātmā} neither slays nor is it slain.
}

The \textit{ātmā} is in the killer, and the \textit{ātmā} is in the one being killed. But the \textit{ātmā} is not perishable. Those who think that the soul of the slayer or the soul of the slain are limited or perishable, are in deep confusion.

Kṛṣṇa says that an ignorant man attributes the functions of the body to the soul. He says, “Look, you should not mix up the body and the soul. Your soul is the Great Observer, it observes things; it comes here and does its \textit{dharma}, but it remains separate from the mind, the intellect, and the body. It is something apart from who you are. The one who thinks that the slayer is the soul is completely mistaken, and the one who thinks that the soul is killed, is completely mistaken.” This does not mean that now that you can think, that, “We are eternal. It’s okay to go and kill people!” No, it doesn’t mean that one can start killing others. Here Kṛṣṇa is talking from the point of view of a warrior, a yoddha. He was on the battlefield. He says, “Look, how can you kill the soul? The soul is eternal, it’s not perishable. How can you say that the killer is the soul, or the person? Which one is the killer? The soul or the person? Neither. The killer is the body; but the soul has nothing to do with the body. The soul doesn’t kill; it is separate from this reality.”

Here Kṛṣṇa explains that the soul doesn’t go through any modifications. The soul in itself is pure and is not even modified by the effects of karma. Even if it is always moving from one body to another, the soul is ever-free. Last week, someone asked me, “Are we liberated or not?” I said, “The \textit{ātmā} of everybody is realized. Everybody is free, because you are part of God. And that part of God which is present inside of you, is forever realized. It is ever-lasting and is fully realized.” In your \textit{sādhana}, what you do is, you polish the outside to get rid of all the karma. Later on, Kṛṣṇa explains this with an analogy (3.38): “If a mirror is covered by dust, the mirror is not touched by the dust; the mirror always stays behind the dust. But when the dust is cleaned away, then the mirror appears clear.”

This is the beauty of the Gita: here Kṛṣṇa describes the soul as being like a mirror. This shows that people already knew about mirrors at that time. It’s true! In the west, they say that mirrors were invented 2000 years ago. They don’t know that 5000 years ago, Kṛṣṇa had already spoken about mirrors in the Gita. Five years ago, scientists discovered a piece of glass that was 2000 years old. They came to the conclusion that this was the oldest glass that had ever existed. But in India, glass existed 5000 years ago! Imagine how advanced this culture was! In Kṛṣṇa’s time, in Dwapara Yuga, people had a certain knowledge. That was why Kṛṣṇa could talk about the Self. Arjuna knew about the Self; he had the knowledge that he was the \textit{ātmā}, that the \textit{ātmā} and the body were completely separate. Kṛṣṇa was only there to remind him of this truth.

\Verse[2.20]
{% --- 1. The Verse ---
  na jāyate mriyate vā kadācin \\
  nāyaṁ bhūtvā bhavitā vā na bhūyaḥ \\
  ajo nityaḥ śāśvato ’yaṁ purāṇo \\
  na hanyate hanyamāne śarīre
}
{% --- 2. The Word-for-Word (Clean text, no formatting) ---
  \textit{na kadācit}---never; 
  \textit{jāyate}---is [the \textit{ātmā}] born; 
  \textit{vā}---nor; 
  \textit{mriyate}---does it die; 
  \textit{na ayam}---[nor is it the case that] this [\textit{ātmā}]; 
  \textit{bhūtvā}---having come into existence; 
  \textit{na bhavitā}---It will not to exist; 
  \textit{vā}---or; 
  \textit{bhūyaḥ}---in the future; 
  \textit{ayam}---this [\textit{ātmā} is]; 
  \textit{ajaḥ}---unborn; 
  \textit{nityaḥ}---eternal; 
  \textit{śāśvataḥ}---everlasting; 
  \textit{purāṇaḥ}---most ancient; 
  \textit{na hanyate}---[it] is not killed; 
  \textit{hanyamāne}---upon slaying; 
  \textit{śarīre}---of the body
}
{% --- 3. The Translation ---
  The \textit{ātmā} is never born, nor does it ever die, nor is it the case that once having come to existence it will cease to exist in the future. It is unborn, eternal, everlasting and most ancient. It is not killed when the body is slain.
}

This part is so wonderful. Kṛṣṇa says that the soul was never born: we can’t even say that we are born from God. All we can say is, “We come out of Him.” The soul is never born, nor does it die. Kṛṣṇa is saying that the soul is not something which is passing by or has modifications with birth and death. Death doesn’t have any grip on the soul: even birth doesn’t have any grip on the soul. It is everfree. Modifications are only for the body. Kṛṣṇa says that all that is “nor is it the case that once having come to existence it will cease to exist in the future.” meaning that what one has gone through before, it won’t happen again. These modifications, these changes which happen to the body, once they are finished, they are gone! You were a baby, then you are middle age, then you get old: can you reverse the process? No, you can’t. But the \textit{ātmā}, is separate from that.

There are the six modifications which the body goes through: birth, becoming, growth, transformation, decay, and destruction. These modifications have nothing to do with the soul. The soul is ever-free from them, “It is unborn, eternal, everlasting and most ancient.” ‘Unborn’: the soul was never born. We can’t say that the soul is born from God. No, it came out of God. You can say that you are born from someone, like a child is born from the mother. But the \textit{ātmā} is an everlasting image of God, an ever-lasting drop of Himself. It was never born. It is everlastingly present inside. It’s not something which is separate from God; it’s eternally a part of God Himself – whether one thinks about it or not, whether one is aware of it or not. ‘Ancient’: the soul is beyond time. We can’t say, “You are a young soul,” or “You are an old soul.” You hear this very often in the esoteric world. “This person is an old soul, he is very wise.” But this kind of wisdom is based on the outside, or what one hears, or perceives only with the mind. Here Kṛṣṇa is reminding Arjuna that the soul is eternal – it’s not young or old. It doesn’t decay. Usually when the body dies, it decays, is destroyed and is finished! The soul doesn’t go through these kinds of transformations. It is everlasting; it is never destroyed. Even if the body is slain, the soul is not.

Here Kṛṣṇa again explains that even if you cut the body, nothing will happen to the soul. Again and again, He is repeating this to remind Arjuna that the Satchitananda is inside of him. But here, Kṛṣṇa is not using the word ‘God’. He is not using the word ‘Paramātmā’. He will use these words later on. Here He is just trying to reason with Arjuna and explains that within this body, there is something greater, which we call the \textit{ātmā}. That is why the yoga of Chapter 2 is called the yoga of the Self. The only thing that you have to do is to realize the soul, the \textit{ātmā}, the Self. This body has been only given to you to realize this truth.

\Verse[2.21]
{% --- 1. The Verse ---
  vedāvināśinaṁ nityaṁ ya enam ajam avyayam \\
  kathaṁ sa puruṣaḥ pārtha kaṁ ghātayati hanti kam
}
{% --- 2. The Word-for-Word (Clean text, no formatting) ---
  \textit{pārtha}---O son of Pṛthā (Arjuna); 
  \textit{yaḥ}---one who; 
  \textit{veda}---knows; 
  \textit{enam}---this [\textit{ātmā}]; 
  \textit{avināśinam}---[to be] indestructible; 
  \textit{nityam}---eternal; 
  \textit{ajam}---unborn; 
  \textit{avyayam}---[and] unchanging; 
  \textit{katham}---how; 
  \textit{kam}---[and] whom; 
  \textit{saḥ puruṣaḥ}---that person; 
  \textit{hanti}---kills; 
  \textit{kam}---whom; 
  \textit{ghātayati}---causes to be killed
}
{% --- 3. The Translation ---
  O son of Pṛthā, if one knows this \textit{ātmā} to be indestructible, unborn, unchanging, and eternal, how and whom does that person kill or cause to be killed?
}

In this verse, Bhagavān Sri Kṛṣṇa brings to Arjuna’s attention that the one who has the true knowledge of the soul, can never think of harming anyone, or cause anyone to be killed. The thought of harming another person will never arise in the mind of one who has realized Brahma Jyaan or \textit{ātmā} Jyaan -- because he is aware. If one has the knowledge that one is not the body, mind or intellect, rather the soul, one will never identify oneself with what one sees on the outside, and believe that one can kill. Kṛṣṇa is saying that there is only one Consciousness that remains the existence of the Spirit only, everywhere. And this Spirit can neither die nor harm anyone or anything. Kṛṣṇa is reminding Arjuna again and again: look, if you come to the Realisation of the Self; if you come to the Realisation that there is only one Divine Spirit, then who is killing who? If you are ignorant of the \textit{ātmā}, the soul, of course then you will see everything from the point of view of the outside world. But if you have the knowledge of the Self, and if you realize the Self, then you rise above this. Kṛṣṇa says, “Knowing this truth, why are you grieving? Why are you crying? Why are you going into this weakness?” Remember that Kṛṣṇa is talking to a soldier on the battlefield. Know that Kṛṣṇa is also talking to each spiritual seeker. Why do you fear to let go? Why do you fear to transcend the outside reality? Devote yourself completely to your \textit{sādhana}! Embrace life and let the Divine reveal Himself to you.

\Verse[2.22]
{% --- 1. The Verse ---
  vāsāṁsi jīrṇāni yathā vihāya navāni gṛhṇāti naro ’parāṇi \\
  tathā śarīrāṇi vihāya jīrṇāny anyāni saṁyāti navāni dehī
}
{% --- 2. The Word-for-Word (Clean text, no formatting) ---
  \textit{yathā}---just as; 
  \textit{naraḥ}---a person [who]; 
  \textit{vihāya}---has discarded; 
  \textit{jīrṇāni}---worn-out; 
  \textit{vāsāṁsi}---clothes; 
  \textit{gṛhṇāti}---puts on; 
  \textit{aparāṇi}---other; 
  \textit{navāni}---new ones; 
  \textit{tathā}---so; 
  \textit{dehī}---the embodied \textit{ātmā}; 
  \textit{vihāya}---having discarded; 
  \textit{jīrṇāni}---worn-out; 
  \textit{śarīrāṇi}---bodies; 
  \textit{saṁyāti}---enters into; 
  \textit{anyāni}---other; 
  \textit{navāni}---new ones
}
{% --- 3. The Translation ---
  Just as a person who has discarded worn-out clothes puts on new ones, so the embodied \textit{ātmā}, having discarded its worn-out bodies, enters into new ones.
}

Here Lord Kṛṣṇa says, “When one’s clothes get old, one changes them for new ones. Like that, when your body becomes old – but you still have \textit{karma} to work out, and you are not yet fully realized – you will take another body. You will come back again into this material field to continue to advance, until the full Realisation awakens.” But then, one will ask, “What is the full Realisation?” Is it the Realisation of the \textit{ātmā} or of the mind? It is the \textit{ātmā}nivedam: the \textit{ātmā} Itself surrenders to the Supreme. Kṛṣṇa refers to the body as old worn-out clothes. The body will go through the same cycle of birth, life, and death, again, again and again. This kind of suffering is because you are trapped in the game of Maya. But, you are meant for something greater.

When the body gets old and it is time to take a new body, it is like when a mother changes the clothes of her child and the child cries. I have seen this with my niece. She cries crazily when her clothes are dirty; but afterwards, when she has new diapers, she is relaxed. Isn’t it like this? A baby will cry when he has dirty diapers, but when the mother changes them, the baby is again happily quiet. When the body gets old, when you have finished with the changes, this means that you are advancing towards a new body, which will be even better! This is preparing you for a better life in the next body, where you will be free from certain dirt and be fine.

The mother stays indifferent through the cries of her child because the mother is concentrating on the welfare of the child. Even so, God, for the good of the \textit{ātmā}, for the good of the Jiva, cares little for the tears of the body, because He knows what is best for each person. Death can be very terrible sometimes for someone and one will ask, “Oh my goodness! Why this and that?” But the Eternal Father knows that when life has come to an end, when the game of life in this body is finished, one has to move on to the next body – until one fully attains the Divine.

\textbf{2016}---“The embodied ātmā, having discarded its worn-out bodies.” The body is subjected to growth, decay and death as a result of time. Time consumes everything. This decay goes through different phases of transformation. Gradually, you start to deteriorate, the physical -- first you lose your hearing. Then you lose your sight, you lose your ability to do many things, and you can’t walk. Gradually, you see how the body is deteriorating. But still, you try to make it happy. You still claim that you are that body. You sit and cry, “I am becoming old. Why am I becoming old? Get me some plastic, I will infuse it inside of me! Botox. I will keep it! I will make sure it stays fresh.”

The saints, the realized souls, accept themselves. The body, it is important. Even Bhagavān Kṛṣṇa is talking about the body being an illusion, yet it is the vehicle that leads you to that realization. He didn’t say, “Yes, okay, because of that, you have to just throw your body somewhere.” But He said that how you see the body, this is important. It doesn’t mean that now that you have heard that it is decaying and falling down, all this stuff, then who cares about it, and finished. No, then that realization will not come to you because this is the vehicle to attain that realization.

Even the Lord Himself manifests on Earth in a body. But you should not forget about the \textit{ātmā} also. Nowadays, people concentrate so much only on the external reality that they forget that they have a soul, they forget that they are living beings. If sometimes you will tell them, “You are an \textit{ātmā}.” “What is this word?” “You are soul!” “Don’t know, I have never heard this word 'soul.'” “What is the soul?” “I don’t know.” You see? In ancient knowledge, they knew.

Here, Bhagavān is saying that this body goes through different stages itself when it is alive, when the \textit{ātmā} is still dwelling in it. The body goes through it, not the \textit{ātmā}. When the body is done, the \textit{ātmā} finds another place to live. It changes house; it looks for a better house.

You see, when your clothes are dirty, you change your clothes, no? Well, you change, I hope you change every day, not only when it’s dirty. But \textit{ātmā} is the same. When the body is worn out completely, then the \textit{ātmā} has to move on and take a new body. Here Bhagavān Kṛṣṇa is also referring to that, “Just as a person who has discarded worn-out clothes puts on new ones”. The same for the \textit{ātmā}; when the body is worn out completely, it can’t do anything, so the body departs from that vessel and enters another one. The \textit{ātmā} is free. You think that you have control over it, no? The mind has made you think that, “Yes, I am controlling.” No, the \textit{ātmā} is free. It will leave when it’s time to leave. You can put as many machines as you want to keep it, but it will not happen. There is a right time to be born, and there is a right time to depart, and the \textit{ātmā} will not stay one second more in that body when it’s time to leave. And sometimes it will create any excuse to leave. Once that veil of ignorance is removed, that enwraps you and keeps you busy, once this veil of Maya is lifted from your mind, when that delusion disappears, all this spell which Maya has placed upon you, which the Lord has allowed Maya to do, also dissolves. But this happens only when you make an effort. Without effort, not possible.

God has given you the mind to make that effort. God has placed you on the spiritual path so that what your soul truly desires happens in an easy way. Then you can rise and realize how everything was conditioned by the mind, how every confusion and delusion, it was only from the mind.

\Verse[2.23]
{% --- 1. The Verse ---
  nainaṁ chindanti śastrāṇi nainaṁ dahati pāvakaḥ \\
  na cainaṁ kledayanty āpo na śoṣayati mārutaḥ
}
{% --- 2. The Word-for-Word (Clean text, no formatting) ---
  \textit{śastrāṇi}---weapons; 
  \textit{na chindanti}---do not cut; 
  \textit{enam}---this [\textit{ātmā}]; 
  \textit{pāvakaḥ}---fire; 
  \textit{na dahati}---does not burn; 
  \textit{enam}---it; 
  \textit{āpaḥ}---water; 
  \textit{na kledayanti}---does not wet; 
  \textit{enam}---it; 
  \textit{ca}---and; 
  \textit{mārutaḥ}---wind; 
  \textit{na śoṣayati}---does not dry [it]
}
{% --- 3. The Translation ---
  Weapons do not cut the \textit{ātmā}; fire does not burn it, water does not wet it, and wind does not dry it.
}

Overtaken by his feelings for his relatives, Arjuna asks, “How can I kill them?” He starts to worry and grieve, thinking, “How could I do that?” Lord Kṛṣṇa reminds him that weapons cannot cut the soul, nor can fire burn it. Referring in this verse to fire, water, and wind, He explains that the elements don’t have any effect on the soul: neither fire, water, or wind, can do anything to the soul. He says, “Even if the body is cut into pieces by a weapon, the soul is not. Even if the body is burned, the soul will not be burned. Even if water dissolves the body, the soul is not touched by it.”

During the time of the Mahabharat War, the elements were used as weapons. Here Kṛṣṇa says that even if a weapon of water is thrown at the body, it will never dissolve the soul. Even if the weapon of air, the Vayu Astra, is thrown at the body, it will not dry up the soul. The soul will always remain untouched by these elements. Kṛṣṇa explains why the soul is superior to the elements. The body is made up of the five elements. Since these elements are already present inside the body, how can these elements destroy the soul? Even if you throw a fire, a bomb, or a missile at the body of a realized person, this won’t affect him. If you send a water or wind weapon at a realized person, it will have no effect on him. You will often see this in the lives of saints. 

\Verse[2.24]
{% --- 1. The Verse ---
  acchedyo ’yam adāhyo ’yam akledyo ’śoṣya eva ca \\
  nityaḥ sarva-gataḥ sthāṇur acalo ’yaṁ sanātanaḥ
}
{% --- 2. The Word-for-Word (Clean text, no formatting) ---
  \textit{ayam}---this (\textit{ātmā}); 
  \textit{acchedyaḥ}---cannot be cut; 
  \textit{ayam}---this (\textit{ātmā}); 
  \textit{adāhyaḥ}---cannot be burnt; 
  \textit{akledyaḥ}---[it] cannot be wetted; 
  \textit{ca}---and; 
  \textit{aśoṣyaḥ}---[it] cannot be dried; 
  \textit{eva}---indeed; 
  \textit{ayam}---this (\textit{ātmā}) [is]; 
  \textit{nityaḥ}---eternal; 
  \textit{sarva-gataḥ}---all-pervading; 
  \textit{sthāṇuḥ}---changeless; 
  \textit{acalaḥ}---immovable; 
  \textit{sanātanaḥ}---[and] everlasting
}
{% --- 3. The Translation ---
  It cannot be cut; it cannot be burnt; it cannot be wetted, and it cannot be dried. It is eternal, all-pervading, changeless, immovable, and everlasting.
}

Kṛṣṇa continues explaining that the soul cannot be destroyed by any weapon because it is indivisible, unmanifest, constant and immutable. Kṛṣṇa says that weapons don’t have any power to destroy this reality. Here the word ‘weapons’ also refers to karma: Kṛṣṇa means that \textit{karma} has no power to destroy the soul. Kṛṣṇa uses the words ‘changeless’, and ‘immovable’ to describe the soul. The soul is still and motionless. The soul is ever focused in the stillness. Whatever is moving, whatever is bound by vibration, is in constant change, and whatever is in constant change, has a limited existence.

“All pervading.” This means that there is no place where the soul is not present. The cosmic Soul is present everywhere. If you are realized, you don’t need to move from one place to another; you don’t need to dance from left to right. This is why the saints don’t need to go around to ten places: just by sitting in one place, they observe everything within the Self. This is the story of Māhā-Viṣṇu, explaining about the Gita and why He always stays in one place with His eyes closed. While His eyes are closed, He is in His eternal, true aspect, one with everyone, and everything. In that state of oneness, He is constantly energising the word of the Gita. Like that, one who is centred in the \textit{ātmā}, perceives the oneness – not out there, but inside oneself.

\Verse[2.25]
{% --- 1. The Verse ---
  avyakto ’yam acintyo ’yam avikāryo ’yam ucyate \\
  tasmād evaṁ viditvainaṁ nānuśocitum arhasi
}
{% --- 2. The Word-for-Word (Clean text, no formatting) ---
  \textit{ayam}---this (\textit{ātmā}); 
  \textit{ucyate}---is said [to be]; 
  \textit{avyaktaḥ}---unmanifest; 
  \textit{ayam}---this (\textit{ātmā}) [is]; 
  \textit{acintyaḥ}---inconceivable; 
  \textit{ayam}---this (\textit{ātmā}) [is]; 
  \textit{avikāryaḥ}---unchanging; 
  \textit{viditvā}---knowing; 
  \textit{enam}---this (\textit{ātmā}) [in this way]; 
  \textit{arhasi}---you should; 
  \textit{evam}---thus; 
  \textit{tasmāt}---therefore; 
  \textit{na anuśocitum}---not grieve
}
{% --- 3. The Translation ---
  This \textit{ātmā} is said to be unmanifest, inconceivable, and unchanging. Knowing the \textit{ātmā} in this way, you should therefore no longer grieve.
}

Here Kṛṣṇa asks Arjuna, “Why are you grieving, My dear? Why are you crying? The soul ‘is unmanifest, it is inconceivable’”. The soul can’t be conceived by the mind, the mind can’t comprehend it. The mind can’t understand the soul, because the soul is beyond the mind.

If you are looking at the soul only from the point of view of the mind, do you think that the mind can understand the soul? A mind which is always moving around is, ‘cañcala’, “restless,” dancing and jumping like a monkey, from one thought to the other. How can the mind which is in constant movement, be still and perceive the unmovable? The mind will always move from one thing to another. Until one has gone deeply into the \textit{sādhana}, deeply into meditation and calmed the mind, only then will the soul reveal itself. But before that, if the mind is still jumping around, the soul is unknowable.

Kṛṣṇa says that the soul “is unchanging.” Everything which is mutable is bound by the mind, and whatever is from the mind is bound by Maya Prakriti, bound by the laws of Nature. The soul is completely separate from Prakriti. It undergoes no transformation, in any circumstance. When one realizes the soul itself, one knows that it is eternal, omnipresent, immovable, everlasting, unmanifested, unthinkable and immutable. So, Kṛṣṇa asks, “Arjuna, why are you grieving?”

\Verse[2.26]
{% --- 1. The Verse ---
  atha cainaṁ nitya-jātaṁ nityaṁ vā manyase mṛtam \\
  tathāpi tvaṁ mahā-bāho nainaṁ śocitum arhasi
}
{% --- 2. The Word-for-Word (Clean text, no formatting) ---
  \textit{atha ca}---and even if; 
  \textit{manyase}---you consider; 
  \textit{enam}---this [\textit{ātmā}]; 
  \textit{nitya-jātam}---to repeatedly go through birth; 
  \textit{vā}---and; 
  \textit{nityam}---repeated; 
  \textit{mṛtam}---death; 
  \textit{mahā-bāho}---O mighty-armed one (Arjuna); 
  \textit{tathāpi}---even then; 
  \textit{tvam}---you; 
  \textit{arhasi}---should; 
  \textit{na śocitum}---not lament; 
  \textit{enam}---this
}
{% --- 3. The Translation ---
  Even if you consider this \textit{ātmā} to repeatedly go through birth and death, O Arjuna, even then you should not lament.
}

Here the Lord explains that the soul is not subject to birth and death, what dies is only the body. But you as the soul, you are eternal. He says, “Don’t grieve, My dear, because grieving makes you weak.” Learn to accept this as the Will of God. And when you learn to accept life – everything that happens in life as the Will of God, as a present from God – then you will be free from grief.

Kṛṣṇa says, “O Arjuna, even then you should not lament.” Again Kṛṣṇa is addressing Arjuna as the, “mighty-armed,” mahābāho. He means, “You are strong! You are not weak. Don’t go into weakness.”

\Verse[2.27]
{% --- 1. The Verse ---
  jātasya hi dhruvo mṛtyur dhruvaṁ janma mṛtasya ca \\
  tasmād aparihārye ’rthe na tvaṁ śocitum arhasi
}
{% --- 2. The Word-for-Word (Clean text, no formatting) ---
  \textit{mṛtyuḥ}---death; 
  \textit{hi}---[is] indeed; 
  \textit{dhruvaḥ}---certain; 
  \textit{jātasya}---for one who is born; 
  \textit{ca}---and; 
  \textit{janma}---birth; 
  \textit{dhruvam}---[is] certain; 
  \textit{mṛtasya}---for one who has died; 
  \textit{tasmāt}---therefore; 
  \textit{tvam}---you; 
  \textit{arhasi}---should; 
  \textit{na śocitum}---not grieve; 
  \textit{arthe}---the matter; 
  \textit{aparihārye}---inevitable
}
{% --- 3. The Translation ---
  Indeed, death is certain for one who is born, and rebirth is certain for one who has died; therefore do not grieve over what is inevitable.
}

\textbf{2014}---There are many people and schools of thought who don’t believe in liberation, in salvation. They always think that they have to be reborn, that they can’t attain the Lord. They bind themselves to this wheel of life and death. They are in ignorance of liberation, in ignorance of attaining the Grace of God.

Kṛṣṇa asks, “Why are you crying, Arjuna? Birth and death are conceived only in ignorance.” If one is in ignorance about liberation, one will never attain it. People who believe that they can’t run away from birth and death – because of their belief – they are stuck in that reality. If people want to believe that, why should you grieve? The drama of birth and death has nothing to do with the soul. The only aim of the soul is to attain the greatest perfection, to attain the Supreme Lord, Nārāyaṇa Himself. Nothing else matters. Kṛṣṇa says, “Whoever wants to go into this drama of birth and death – it is their choice!” Some of the Kauravas had the knowledge of the Self: Bhishma Pita knew, Dronacharya knew, Kripacharya knew. But they willingly chose to play their roles in the drama and be on the Kauravas side. Kṛṣṇa says, “If they have personally chosen this side, why are you grieving? They are fully aware of their choice to be on the Kauravas side. So don’t grieve for that.”

\textbf{2016}---“Indeed, death is certain for one who is born.” You are born, you will die. Whether you fear death or not, you will have to leave this body. Sad, no? Whether you are sad or you are happy, it doesn’t matter. Even the Lord Himself has come, He has taken birth, He has manifested Himself, but when He leaves, He also leaves this body here and disappears. Because what you took from nature has to go back to nature. This is the law of nature. But the \textit{ātmā} has come from a higher plane. and it will have to go to that higher plane. All the karmic things which are around; the dirt, the dust which you have accumulated through lives, this also has to be removed. Sooner or later you will reach there. But when you have that Grace to be in company, in association with the holy people, blessed people, who will help you to liberate yourself, to awaken yourself, to reach salvation, then that Grace of the Lord makes everything easy to realize yourself.

“Indeed, death is certain for one who is born, and rebirth is certain for one who has died.” It’s a cycle. You are born, the moment you are born, you are looking toward death. The moment you let go of life, as life leaves you, it is already entering another body. Life is carrying on.

“Therefore, do not grieve over what is inevitable.” Here, Bhagavān is saying to Arjuna, “Why are you crying? Birth and death are conceived only in ignorance. It’s only in the mind. If one doesn’t have that knowledge of liberation, if one is ignorant about it, one doesn’t attain it. Here Bhagavān is saying, look at the world, how many people know about liberation? How many people know that they are the \textit{ātmā}? They are content just to work, eat, drink, sleep, enjoy pleasure, and finish. What do they attain in life? Nothing. Here Bhagavān said to Arjuna clearly, “Those who are in ignorance about the fact of the Self, about liberation, they don’t attain Him. But those who are aware of liberation, whether they want it or not, it will be given on a golden plate to them.” So, being aware of this will make you change, transform and grow. Due to that belief, due to that trust, Bhagavān Himself will make everything possible so that you can attain Him. That mind itself will purify so that Sriman Nārāyaṇa Himself reveals Himself to you.

Whoever wants to go into this \textit{dharma} of birth and death, you have a choice. You know that you are eternal. You know that you should not play that game, yet you choose it. Often you know so many things, but you do the opposite of it. Whatever you do opposite, you will have to suffer.

Here Bhagavān is saying, “Look at these, the Kauravas, they have great people with them. They have great knowledge of the Self. Bhishma, the grandsire, Dronacharya, which is the one who knew so much. They are the great teachers. Kripacharya, who is very knowledgeable, yet the willingness to side with the opposite - it’s their choice, when they know they will lose everything.” Like that, the world also goes like this. People always know that, “The world will never give me the full satisfaction,” but we long for that world, we cry for that world, we are attached to that world, we make that world our utmost reality. Yet we long for liberation also, because it’s a pleasing word. We long for Love for God, but we don’t want to let go of this world. We don’t want to let go of anything. Here, in the mind. I am not telling you to let go of the outside, but it is about how much importance you give to it. The more importance you give, the more sorrowful you become. That’s what Bhagavān Kṛṣṇa is saying right now.

“The acharyas, they are fully aware of their choice to be on the Kauravas’ side, so why do you sit and cry, Arjuna, when they know about it?”

\Verse[2.28]
{% --- 1. The Verse ---
  avyaktādīni bhūtāni vyakta-madhyāni bhārata \\
  avyakta-nidhanāny eva tatra kā paridevanā
}
{% --- 2. The Word-for-Word (Clean text, no formatting) ---
  \textit{bhārata}---O descendant of Bharata; 
  \textit{bhūtāni}---all beings; 
  \textit{avyakta-ādīni}---are unmanifest in the beginning; 
  \textit{vyakta-madhyāni}---manifest in the middle; 
  \textit{eva}---[and] indeed; 
  \textit{avyakta-nidhanāni}---unmanifest at death; 
  \textit{tatra}---in that case; 
  \textit{kā}---what [is the reason for]; 
  \textit{paridevanā}---lamentation
}
{% --- 3. The Translation ---
  O descendant of Bharata! All beings are unmanifest in the beginning, manifest in the middle and unmanifest at death. So why lament over this?
}

In the beginning, before birth, the soul has no connection with this body. “In the middle”, when the body gets manifested, the soul still doesn’t have any connection to the body. And after death, the soul will maintain no connection to the gross body. The soul manifests in a body to do its karma, to finish the \textit{karma} which had been created in past lifetimes, yet it remains untouched by anything. Then the soul frees itself upon passing from this body: it doesn’t hang onto the body. The \textit{ātmā} goes through this process. Lord Kṛṣṇa says that this is like a dream world: the dream world is nonexistent before and after the dream. It is only during the dream that the dreamer seems to have a relationship with the body. It is only during the dream that you feel that everything is real. But in reality, it is not! You don’t have any connection to that dream. So why do you lament? A dream is only for a short period of time.

\Verse[2.29]
{% --- 1. The Verse ---
  āścarya-vat paśyati kaścid enam \\
  āścarya-vad vadati tathaiva cānyaḥ \\
  āścarya-vac cainam anyaḥ śṛṇoti \\
  śrutvāpy enaṁ veda na caiva kaścit
}
{% --- 2. The Word-for-Word (Clean text, no formatting) ---
  \textit{kaścit}---one person; 
  \textit{paśyati}---sees; 
  \textit{enam}---it (the \textit{ātmā}); 
  \textit{āścarya-vat}---as a wonder; 
  \textit{ca}---and; 
  \textit{tathā eva}---likewise; 
  \textit{anyaḥ}---another; 
  \textit{vadati}---speaks [of it]; 
  \textit{āścarya-vat}---as a wonder; 
  \textit{ca}---yet; 
  \textit{anyaḥ}---another; 
  \textit{śṛṇoti}---hears; 
  \textit{enam}---[of] it; 
  \textit{āścarya-vat}---as a wonder; 
  \textit{ca}---but; 
  \textit{api}---even; 
  \textit{śrutvā}---after hearing; 
  \textit{enam}---[about] it; 
  \textit{eva}---indeed; 
  \textit{na kaścit}---no one [really]; 
  \textit{veda}---knows [it]
}
{% --- 3. The Translation ---
  One person sees the \textit{ātmā} as a wonder, likewise another speaks of it as a wonder; yet another hears of it as a wonder. But even after hearing about it, no one really knows it.
}

Here Kṛṣṇa says that the soul is something wonderful and that it’s very rare to meet an individual in this world who has perceived the true wonders of the soul. The soul is not a worldly object. He says, “None truly know it,” because the soul cannot be perceived by the senses, the mind or the intellect. To know something, you usually need the mind, no? But how can the mind perceive something that is beyond the mind? The soul is transcendental and unique. It is the only reality that abides eternally within the consciousness of every human being. The only eternal reality is the \textit{ātmā}. How can one even marvel at the wonders of the soul if one is so much focused on the outside world? And to just describe the soul, does not mean that one is truly perceiving the wonders of the soul.

“That Self, which we look on and speak of and hear of as that which is wonderful, beyond our comprehension, no human mind has ever known this Absolute.” Here Kṛṣṇa says that all words are limited. We are using beautiful words to describe the soul, but the soul can’t be comprehended by words. We can only use words to explain. But the one who has truly realized the soul, doesn’t even talk about it, because there are no words to describe it. The moment you start talking about the soul, it becomes unreal. Kṛṣṇa says that to perceive the soul, one has to go deep inside and enter into the quietness. Then one can see the brightness of the soul and one can marvel at the wonders of the soul. Whoever has experienced this knows that you can see it.

You don’t need to talk about it. It’s emanating from deep within you to the outside. This Realisation of the soul emanates from the person who is in this Reality. So it’s not about words. Yet one can feel the soul, one can see it.

There are a few special people in the world who possess such virtue and purity of heart, that when they hear a description of the soul, they spontaneously experience the true wonders of the soul. Because of their faith in God, their reverence for God, just hearing about the wonders of the soul creates a great bhav inside them. This is very rare. Kṛṣṇa says, “no one really knows it” because it is beyond the mind.

\Verse[2.30]
{% --- 1. The Verse ---
  dehī nityam avadhyo ’yaṁ dehe sarvasya bhārata \\
  tasmāt sarvāṇi bhūtāni na tvaṁ śocitum arhasi
}
{% --- 2. The Word-for-Word (Clean text, no formatting) ---
  \textit{dehī}---the embodied \textit{ātmā} [is]; 
  \textit{dehe}---in the body; 
  \textit{sarvasya}---of all; 
  \textit{bhārata}---O descendant of Bharata; 
  \textit{ayam}---this (\textit{ātmā}) [is]; 
  \textit{nityam}---eternal; 
  \textit{avadhyaḥ}---[and] indestructible; 
  \textit{tasmāt}---therefore; 
  \textit{tvam}---you; 
  \textit{arhasi}---should; 
  \textit{na śocitum}---not grieve for; 
  \textit{sarvāṇi}---all; 
  \textit{bhūtāni}---living beings
}
{% --- 3. The Translation ---
  The embodied \textit{ātmā} is in the bodies of all, O descendant of Bharata. It is eternal and indestructible, therefore you should not grieve for any living being.
}

Here Kṛṣṇa explains that the soul, the dweller in the bodies of all living entities, can never be killed. The Lord says that in the bodies of all human beings in this world, there resides only One Supreme Soul. It is only due to ignorance that one sees differences in the diverse bodies. In reality, there is no diversity: there is only One Supreme Soul which we are all connected to, and which we are all part of. What is in me, is in you. It is in Arjuna. It is in Kṛṣṇa. Lord Kṛṣṇa reminds Arjuna, “I am in everyone as the Supreme Soul, Paramātmā.” Every \textit{ātmā} is part of Paramātmā: every \textit{ātmā} is part of this cosmic body which makes Paramātmā, the only One Reality. He uses the word ‘indestructible’: this ultimate reality, the soul, can never be destroyed. It can never be destroyed by anything or by anyone. Again and again Kṛṣṇa says, “Don’t grieve. Arjuna, prepare yourself to fight!”

\Verse[2.31]
{% --- 1. The Verse ---
  sva-dharmam api cāvekṣya na vikampitum arhasi \\
  dharmyād dhi yuddhāc chreyo ’nyat kṣatriyasya na vidyate
}
{% --- 2. The Word-for-Word (Clean text, no formatting) ---
  \textit{api ca}---additionally; 
  \textit{avekṣya}---considering; 
  \textit{sva-dharmam}---your personal duty; 
  \textit{na arhasi}---you should not; 
  \textit{vikampitum}---hesitate; 
  \textit{kṣatriyasya}---[since] for a kṣatriya; 
  \textit{hi}---indeed; 
  \textit{na vidyate}---there does not exist; 
  \textit{anyat}---anything; 
  \textit{śreyaḥ}---more auspicious; 
  \textit{dharmyāt}---than a righteous; 
  \textit{yuddhāt}---war
}
{% --- 3. The Translation ---
  Additionally, considering your personal duty, you should not hesitate, since for a \textit{kṣatriya}, there is nothing more auspicious than a righteous war.
}

Kṛṣṇa again reminds Arjuna of his duty as a Kshatriya, as a soldier. He says, “There is no greater thing for a soldier to do, than to fight. You should not divert yourself from your aim. Knowing now that the soul is immortal and that the body is separate from the soul, why do you lament? You should go and fight! Don’t just stand there and think! Don’t continue to analyse the problem from many different points of view. You are from the warrior caste, so do your duty! This is the natural duty of a Kshatriya – to not run away from a battle. If ‘the Kshatriya’ starts running away from the battle, all is finished! Then all that you have worried about will happen. It will not happen because you are fighting, but because you are running away from fighting.”

In Chapter 1, Arjuna says to Lord Kṛṣṇa that society will be destroyed and so on, and so on, and so on, if there is a war. Here Kṛṣṇa says, “Hey listen! Society will not get destroyed because of this war. It will get destroyed if you run away, if you don’t do your duty as a warrior.” Imagine you are at the doctor’s and the doctor puts you on the operating table. He has already put you to sleep, then he cuts you open and in the middle of the operation he says, “Oh, no, I can’t look at it! I am running away.” You can’t run away from doing your \textit{dharma}, your duty. God has given each person their \textit{dharma} to do in life. You have to do your duty! My duty is to sit here and talk to you, so I am sitting here and I am talking to you. Your duty is to listen, so you are listening. Your duty in the outside world is to do what God has given you to do as your daily duty in life. You have to learn to accept it. That’s why Kṛṣṇa says, “You have to accept your duty as a Kshatriya. There is nothing more wonderful and righteous than to do your duty.”



\textbf{2016}---“Additionally, considering your personal duty, you should not hesitate.”

“Don’t look at the others. Do your duty, what you have been assigned for, and do it fearlessly. You as a Kshatriya, as a soldier, your \textit{dharma} is to fight. You should never divert yourself from your aim. Especially when you have the knowledge of the soul, that the soul is immortal and that this body will one day separate from it, why do you cry?” Here He said, “Don’t just stand and think. Don’t just analyze your worries and problems. All of you are on this battlefield. Don’t run away from it. Don’t become a coward. Finish what you have come here for.” This body is manifested, the \textit{ātmā} is unmanifested. The body is limited, the \textit{ātmā} is unlimited. Here Bhagavān is saying, “Now that you have this knowledge of the Self, remind yourself that that is your \textit{karma} here, remind yourself what your \textit{dharma} as a warrior is. To run away?” You see, He gave him the knowledge of the \textit{ātmā} first, then He reminded him why that \textit{ātmā} has come here. He reminded him of his duty toward that body also, because the \textit{ātmā} doesn’t just incarnate like that. When it takes a body, it’s due to certain \textit{dharma} that it has to fill and finish. It’s not to attach, but it’s to detach.

“Go through that battle, not for the sake of killing, but for the sake of doing your duty, why you have incarnated. You are bound to do that as karma, because you are born into a warrior family, and to fight this fight is the most respectable duty for people who belong to your standard.”

Here Bhagavān is saying, after revealing the eternal soul to Arjuna, “You can’t run away from your duty. It’s important. Because if you run away from your duty, it will not be good. You have to do your duty.”

\Verse[2.32]
{% --- 1. The Verse ---
  yadṛcchayā copapannaṁ svarga-dvāram apāvṛtam \\
  sukhinaḥ kṣatriyāḥ pārtha labhante yuddham īdṛśam
}
{% --- 2. The Word-for-Word (Clean text, no formatting) ---
  \textit{pārtha}---O Pārtha (Arjuna); 
  \textit{kṣatriyāḥ}---[the] kṣatriyas; 
  \textit{sukhinaḥ}---[are] joyous; 
  \textit{labhante}---[when they] get; 
  \textit{īdṛśam}---such an; 
  \textit{yuddham}---[opportunity for] battle; 
  \textit{ca}---and; 
  \textit{upapannam}---that has come; 
  \textit{yadṛcchayā}---by its own accord; 
  \textit{apāvṛtam}---[is an] open; 
  \textit{svarga-dvāram}---gate to Heaven
}
{% --- 3. The Translation ---
  O Pārtha, the \textit{kṣatriyas} rejoice when they get such an opportunity for battle. A war that comes of its own accord opens the gates to Heaven.
}

Here again, Bhagavān Sri Kṛṣṇa tells Arjuna, “ O Pārtha, the \textit{kṣatriyas} rejoice when they get such an opportunity for battle.” He says, “Look! You should be rejoicing! You have not come here to make the war happen, but it came to you! It gives you the opportunity to do your \textit{dharma}. You should welcome this opportunity to do your \textit{dharma}. You should accept it! Welcome it!”

“...it opens the gates to Heaven”: when you do your duty in life, accepting what God is giving you to do, it draws you closer and closer to God-realization. By accepting what God is giving you in your daily life, you are opening the gates to receive more. By showing that you are doing your duty; by accepting what God gives you in life you are showing God that you are ready for a greater purpose in this life, for greater work. Kṛṣṇa says, “If you run away from this, it will not bring you any good. It will pain you. It will be unrighteous. Especially on such an occasion, where doing this duty will help you open all the doors to Heaven.” He says, “Accept life. Accept whatever God gives you. Declaring this war will open the gates of Heaven.”

You know, in war, they often check how you died – whether you were running away from the war or you were facing the war. They check this in a war, no? To see if you got the bullet from the front or from the back. If the bullet came from in front of you and the enemies were in front of you, then you will get lots of praises. But if the bullet came from behind and the enemy was behind you, that’s a problem. Your name will not even be mentioned – because that means that you were running away. When one meets one’s death in a righteous war, one obtains direct access to Heaven. This is Karma Bhumi, \texti{puṇya} Bhumi, so because of all this good merit, one will go to Heaven. Happy are the Kshatriyas who have such an opportunity! Not all Kshatriyas get the opportunity to fight in a righteous war. Only the fortunate ones get to do such a duty.

When I refer to a righteous war, I am not talking about a religion war. Don’t mix these up. Righteous is beyond religion. Nowadays war is due to religion: one religion claims it is better than another. Then they make a war and each side says that it is fighting for the right cause. No! This has nothing to do with a righteous war. A righteous war, a holy war, is the \textit{dharma} that one has to do in life. That’s why Kṛṣṇa says here, “If you run away, you will not do your swadharma, you will not do your duty. Especially when it is coming to you. You are not making the war, you are not the cause of the war, but this war has been given to you to fight.”

\Verse[2.33]
{% --- 1. The Verse ---
  atha cet tvam imaṁ dharmyaṁ saṅgrāmaṁ na kariṣyasi \\
  tataḥ sva-dharmaṁ kīrtiṁ ca hitvā pāpam avāpsyasi
}
{% --- 2. The Word-for-Word (Clean text, no formatting) ---
  \textit{atha}---but; 
  \textit{cet}---if; 
  \textit{tvam}---you; 
  \textit{na kariṣyasi}---do not fight; 
  \textit{imam}---this; 
  \textit{dharmyam}---righteous; 
  \textit{saṅgrāmam}---war; 
  \textit{tataḥ}---thus; 
  \textit{hitvā}---abandoning; 
  \textit{sva-dharmam}---your duty; 
  \textit{ca}---and; 
  \textit{kīrtim}---honor; 
  \textit{avāpsyasi}---you will incur; 
  \textit{pāpam}---sin
}
{% --- 3. The Translation ---
  But if you do not fight this righteous war, thus turning away from personal duty and honor, you will incur sin.
}

Here Lord Kṛṣṇa reminds Arjuna, “If you don’t learn to accept what God is giving you in life and run away from your duty, run away from your \textit{dharma}, then you are committing a big sin.” If you don’t first accept what God is giving you to do in life, and you don’t realize your \textit{dharma}, you will fail in life, and you will ‘incur only sin’.

It is because of the fear of sin that Arjuna wants to run away from fighting. But here Kṛṣṇa turns it around and says, “No, My dear! You will commit sin if you run away.” First Arjuna wants to run away, fearing the sin and the guilt that will come upon him. But Kṛṣṇa says, “No way, don’t do that! It will be not virtuous to do such a thing. If you run away, you will not awaken.”

\textbf{2016}---Here, He says, “Look, Arjuna, you have to fight because this is your \textit{dharma}. Your duty is to fight as a warrior. If you don’t do that, you will incur great sin. You will increase your dismerit, which will bring you to hell for a long time. You will have to do your duty.”

Very often, people ask, “Swamiji, what is my \textit{dharma}?” But, you see, your \textit{dharma} is where God has placed you. What you are doing, you have to learn to accept it. If you run away from that, not accepting what God is giving you directly, you will never do your greater good, you will never do your true \textit{dharma}. First you should accept what is ..?..[39:48] given to you. This is what Bhagavān Kṛṣṇa is saying to Arjuna right now, “You have been born to do that. Changing that will not bring you anywhere. Once you accept that reality, you will be free.” He didn’t say that, yes, you have to change to be free. No, by doing your \textit{dharma} properly in whatever you do, by keeping your mind focused upon Him, no matter what work you are doing, you will realize Him.

Because of fear of sin – before, remember in the first chapter, He said, “If I kill everybody, I will have great sin, I will go to hell.” Because of that, here Bhagavān Kṛṣṇa has reversed it, and said, “Look, this is your perception of sin, but it is not like that.” If you have fear of sin and run away from your \textit{dharma}, then you will have a lot of sin. You will have committed sin toward what God has given you. He said, “No, do your \textit{dharma}. Be wise about it, you will awaken. But if you run away, then nothing good will come of it. If your mind is impure, nothing good can come out of it. If your mind is filled with guilt, you will become the victim of your own guilt. You see, everybody, they have a \textit{dharma}, and when you do that \textit{dharma}, everything goes smoothly. If you run away, you will not be able to bear that sin, you will not be able to bear that guilt. You are running away for not doing your \textit{dharma}, but then you will have another guilt. One guilt is the guilt of the mind, but then you will have the guilt of the soul, which is even worse. Because the guilt of the mind, when you have created certain things, you can forget about it after some time. But the guilt of the soul, when the soul creates that guilt, it’s far more painful than the guilt of that mind, and that will incur more sin.”

\Verse[2.34]
{% --- 1. The Verse ---
  akīrtiṁ cāpi bhūtāni kathayiṣyanti te ’vyayām \\
  sambhāvitasya cākīrtir maraṇād atiricyate
}
{% --- 2. The Word-for-Word (Clean text, no formatting) ---
  \textit{api ca}---then; 
  \textit{bhūtāni}---people; 
  \textit{avyayām}---forever; 
  \textit{kathayiṣyanti}---will speak; 
  \textit{te}---of your; 
  \textit{akīrtim}---disgrace; 
  \textit{ca}---and; 
  \textit{sambhāvitasya}---for an honorable man; 
  \textit{akīrtiḥ}---dishonor; 
  \textit{atiricyate}---is worse; 
  \textit{maraṇāt}---than death
}
{% --- 3. The Translation ---
  Then people will forever speak of your disgrace, and for an honorable man, dishonor is worse than death.
}

Here Kṛṣṇa is using reverse psychology, talking about guilt in another way. He says to Arjuna, “If you don’t do your duty, you will lose your reputation and you will be exposed to sin. People will talk badly of you, but not only that, the rishis and the celestials will all look at you in a questioning way. Your ancestors will also be crying. So, get up and fight! Don’t let people talk ill of you. You are from the warrior clan!”

“...dishonour is worse than death.” It is very important to understand that, during this time, people were very concerned about their honour and about giving their word. When they gave their word, they kept it. They would rather have died than not have kept their word. The Lord reminds Arjuna, “It will be painful to lose people’s respect and be dishonoured: better death than dishonour.”

Kṛṣṇa tells him that if he runs away, he will not be able to walk in this world because everybody in every corner will be saying nasty things about him. Kṛṣṇa doesn’t leave it there. He says that even in Heaven, they will say bad things about him. He says, “Forget about your honour here on Earth: ‘men will recount your perpetual disgrace.’ For time immemorial, people will remember it. And this will not only react here, it will also react in Heaven. Imagine that one day you will go to Heaven, and there they will also nag you! When you are dying here, you think you will go somewhere and be free, but this will continue with you there.”

Know that this is very important to understand. To run away from certain things, people often kill themselves, they commit suicide. They think, “Ah, I have escaped!” Yes, in this life, people have run away. Do you think that their problem has left them? No, it has not left them! It will continue and become even worse! Kṛṣṇa says, “They will also be nagging you in Heaven.” This means that in the next life, it will be worse.

\Verse[2.35]
{% --- 1. The Verse ---
  bhayād raṇād uparataṁ maṁsyante tvāṁ mahā-rathāḥ \\
  yeṣāṁ ca tvaṁ bahu-mato bhūtvā yāsyasi lāghavam
}
{% --- 2. The Word-for-Word (Clean text, no formatting) ---
  \textit{mahā-rathāḥ}---the great warriors; 
  \textit{maṁsyante}---will think; 
  \textit{tvām}---you; 
  \textit{uparatam}---have withdrawn; 
  \textit{raṇāt}---from the battlefield; 
  \textit{bhayāt}---out of fear; 
  \textit{ca}---and; 
  \textit{yāsyasi}---you will earn; 
  \textit{lāghavam}---the disrespect; 
  \textit{yeṣām}---of those [for whom]; 
  \textit{tvam}---you; 
  \textit{bhūtvā}---had become; 
  \textit{bahu-mataḥ}---highly regarded
}
{% --- 3. The Translation ---
  The great warriors will think that you have fled from the battlefield in fear and you will earn the disrespect of those who held you in high esteem.
}

This verse is also about Arjuna’s honour, about how people will look at him. Kṛṣṇa says, “‘The great warriors’: Bhima, your brothers, Bhishma, Duryodhana, Drona and many others – they all have great esteem for you. If you run away, they will have an extremely low image of you. They will be broken! They will be shocked if you run away out of fear!” Raising the question of fear on the battlefield, Lord Kṛṣṇa says, “In this state of mind, Arjuna, you are being disgraceful, because a warrior should not have fear. These people will look at you and say, ‘My goodness! Where did this fear come from?’ They will never believe that your motive in fleeing battle was to protect them. How could they know that? For them, you ran away because of cowardice. They will never know your true motive. When the other Pandavas see you sad and running away out of fear, they will be pitied by the Kauravas and will also develop compassion for them, and then this whole war will be completely different. This war is to be a dharmic war and restore balance. But there will be no dharmic war. Automatically there will be adharma.”

\Verse[2.36]
{% --- 1. The Verse ---
  avācya-vādāṁś ca bahūn vadiṣyanti tavāhitāḥ \\
  nindantas tava sāmarthyaṁ tato duḥkha-taraṁ nu kim
}
{% --- 2. The Word-for-Word (Clean text, no formatting) ---
  \textit{ca}---moreover; 
  \textit{tava}---your; 
  \textit{ahitāḥ}---enemies; 
  \textit{nindantaḥ}---ridiculing; 
  \textit{tava}---your; 
  \textit{sāmarthyam}---ability; 
  \textit{vadiṣyanti}---will speak; 
  \textit{bahūn}---many; 
  \textit{avācya-vādān}---words that should not be spoken; 
  \textit{nu}---indeed; 
  \textit{kim}---what [could be]; 
  \textit{duḥkha-taram}---more painful; 
  \textit{tataḥ}---than that
}
{% --- 3. The Translation ---
  Moreover, your enemies, ridiculing your ability, will use words which should never be uttered. Indeed what could be more painful than that?
}

You see, Kṛṣṇa is using the feeling of guilt, the knife, and putting it inside again. He says, “If you run away from your enemies, the infamy of the public scandal will cause acute distress. It will be terrible. Duryodhana and the other Kauravas will use this; they will enjoy this. They will ridicule you and not only you – they will also ridicule your whole family. They will say to your other brothers, ‘Your brother, Arjuna, ran away.’ They will start calling your family, even your mother, many names. They will speak words of reproach and defame you. If you lose fame, Arjuna, how would you take this?”

You see, the Pandavas also cared very much about their fame. Kṛṣṇa says, “This will bring great suffering to you, Arjuna. Whereas these enemies who you are concerned about will be joyfully happy; they will plot and use this to do more mischief. They will use it for their own purpose. They will say that Arjuna was not a warrior.” Kṛṣṇa reminds Arjuna, “They will even dishonour your holy bow. They will say many disrespectful things about it. How could you bear that? If they talk badly about your Gandiva, could you bear this? You know how powerful your Gandiva is! How strong it is! You will not be able to bear this disrespect to your Gandiva.” Pointing to his bow, He says, “There is nothing worse than running away. So look Arjuna, open your eyes! See that when you accept what God is giving you to do, it will bring you great honour to you. But the opposite of that, what will it be?”

\Verse[2.37]
{% --- 1. The Verse ---
  hato vā prāpsyasi svargaṁ jitvā vā bhokṣyase mahīm \\
  tasmād uttiṣṭha kaunteya yuddhāya kṛta-niścayaḥ
}
{% --- 2. The Word-for-Word (Clean text, no formatting) ---
  \textit{vā}---if; 
  \textit{hataḥ}---slain; 
  \textit{prāpsyasi}---you shall gain; 
  \textit{svargam}---Heaven; 
  \textit{vā}---if; 
  \textit{jitvā}---victorious; 
  \textit{bhokṣyase}---you shall enjoy; 
  \textit{mahīm}---the Earth; 
  \textit{tasmāt}---therefore; 
  \textit{uttiṣṭha}---arise; 
  \textit{kaunteya}---O son of Kuntī (Arjuna); 
  \textit{kṛta-niścayaḥ}---with firm resolve; 
  \textit{yuddhāya}---for battle
}
{% --- 3. The Translation ---
  If slain, you shall gain Heaven; if victorious, you shall enjoy the Earth. Therefore, arise, O son of Kuntī, with firm resolve to fight.
}

As you saw earlier, Arjuna is confused in his mind and even wants to run away from fighting. Here Kṛṣṇa says, “Your state of confusion has been removed, which means that you will now win the battle, you will conquer the Kauravas. You will be victorious.” He says, “You shall enjoy the Earth, – therefore arise, O son of Kunti, and resolve to fight. You will be victorious in this battle and you will also be victorious in Heaven. You will win both ways.” So this removed Arjuna’s worry about whether he would be victorious or not. Here Kṛṣṇa says it, “You will be victorious!” Because He knows that Arjuna will have full victory over the Kauravas. This is a tricky point. Kṛṣṇa says, “You shall gain Heaven.” He doesn’t say, “You have already won Heaven.” He is saying this in a way that shows that He already knows that they will win. He says to Arjuna, “You are the victorious one. Because this war has already begun, this war has already ended.”

Physically, the war has not yet ended; but spiritually, Kṛṣṇa has already revealed the future to him. Kṛṣṇa says, “ If slain, you shall gain Heaven; if victorious, you shall enjoy the Earth.” Here Kṛṣṇa uses the word “if.” He says, “If you die, you will go directly to Heaven.” That’s good! It’s a passport! A visa! He tells Arjuna, “If you meet your death on this battlefield, you will not be a loser. You will get a direct passport to Heaven. Whereas, if you win the battlefield, if you don’t die, you will inherit all the sovereignty of the Earth. You will be the king of this Earth.” Again Kṛṣṇa says, “Shake off your faint-heartedness! Take up your position in this chariot and be determined to fight!”

\Verse[2.38]
{% --- 1. The Verse ---
  sukha-duḥkhe same kṛtvā lābhālābhau jayājayau \\
  tato yuddhāya yujyasva naivaṁ pāpam avāpsyasi
}
{% --- 2. The Word-for-Word (Clean text, no formatting) ---
  \textit{kṛtvā}---considering; 
  \textit{sukha-duḥkhe}---pleasure and pain; 
  \textit{lābha-alābhau}---gain and loss; 
  \textit{jaya-ajayau}---victory and defeat; 
  \textit{same}---[to be the] same; 
  \textit{tataḥ}---then; 
  \textit{yuddhāya}---for battle; 
  \textit{yujyasva}---engage yourself; 
  \textit{na}---not; 
  \textit{evam}---thus; 
  \textit{pāpam}---sin; 
  \textit{avāpsyasi}---you will incur
}
{% --- 3. The Translation ---
  Considering pleasure and pain, gain and loss, victory and defeat to be the same, prepare yourself for battle. In this way, you will not incur sin.
}

Here Kṛṣṇa says to Arjuna, “Awaken, My dear! Realize, that all this is just a drama put on by Me.” Later on, He will say, “I am the One who does everything, not you.” Here He is just telling Arjuna to awaken, “See it from the point of view of your soul. Don’t see only what your mind perceives.” The Lord again tells him, “Fight Arjuna! Turn to battle; this way you shall not incur sin. Treating pleasure and pain alike, O Arjuna, get ready and do your \textit{dharma}!” He says, “Don’t bother about whether you want to gain the kingdom of Heaven. You don’t desire the kingdom, so you want to run away from the kingdom. If you don’t desire all these things, if you don’t want the kingdom, then it is easy for you, no? Just fight! And, cultivate the feeling of equanimity. You shall engage in this fight. I’m not telling you to fight to gain something. I’m not promising you anything, but just fight! Don’t think about how you fight. You already know you don’t desire anything. That’s good! That’s the first thing. Very good! Please just fight! Lift up your weapon and fight. That’s the only thing I am asking you. If you just do what I tell you, it will bring everlasting and supreme peace.”

Here, Kṛṣṇa is saying, “My dear, I am the \textit{guru} and you are the disciple. Your duty is just to listen. As the \textit{guru}, I am ordering you to fight, so don’t even think about it. You have already decided that you don’t want anything. In your mind, you have thought about it and you know you don’t want anything. This is good! Then just do what I am guiding you to do. I am showing you the path to supreme peace.” Bhagavān Sri Kṛṣṇa continues, “If you fight, if you fight in this spirit, no guilt from committing any sin will arise in you.” This is Kṛṣṇa’s reply to Arjuna’s question in the first chapter. He says, “There will be no sin involved if you do your \textit{dharma}. If you participate in the battle, treating pleasure and pain, ‘victory and defeat equally’, with this attitude, you will not be bound to the results of the battle. No \textit{karma} will be created. You will not be responsible. There will be no slightest trace of \textit{karma} arising when you do your \textit{dharma}. You will be completely released from the bondage of action.” Karma arises due to expectation; when one expects in whatever one does, one creates karma. Karma binds and makes one suffer. But if one does one’s duty, if one does anything in the spirit of serving the Lord, one doesn’t create any karma.

\textbf{2016}---Arjuna says, “How can I fight?” Still, that mind of Arjuna is trying to run away from fighting. Bhagavān Kṛṣṇa has given him something to remove his confusion. He is standing now on the battlefield, realizing that this is his duty as a soldier, he must fight.

Bhagavān said to him, “Make grief and happiness, loss and gain, victory and defeat equal to your soul.” Remind yourself that this is all just a game. Take this pair of opposites as being equal. Even if they appear different, they are the same.

A wise man should not be taken over by joy when something is happening, just stay calm. And those who are wise don’t complain when there is some trouble. They analyze the situation. They are in balanced state in both situations; in a happy moment or in an unhappy moment. You see, in the mind there are countless thoughts that keep jumping up and down at all times. When something is happy, the mind is happy. When something is terrible, the mind becomes terrible and amplifies it also. A small trouble like this, you will stretch it out.

Here Bhagavān Kṛṣṇa is saying to Arjuna, “Stay awake. Realize that all this is just a play of Mine. I am the one who is doing everything, not you. I am telling you all this to remind you, remind your mind, remind your soul that it is the Supreme Lord who is doing it, and that you have been incarnated to do one thing. You have incarnated to do your \textit{dharma}. Don’t let yourself be fooled by left and right, by the pairs of opposites. When you have this attitude, then there is no karma. But when there is judgment, this is where \textit{karma} is being created. When \textit{karma} is being created, either in good or bad, one will always be feeling the pressure of not being satisfied enough. Then one can’t truly do one’s duty. How would one serve the Lord? How would one attain Him?

\Verse[2.39]
{% --- 1. The Verse ---
  eṣā te ’bhihitā sāṅkhye buddhir yoge tv imāṁ śṛṇu \\
  buddhyā yukto yayā pārtha karma-bandhaṁ prahāsyasi
}
{% --- 2. The Word-for-Word (Clean text, no formatting) ---
  \textit{eṣā buddhiḥ}---this knowledge [which]; 
  \textit{abhihitā}---has been taught; 
  \textit{te}---to you [so far]; 
  \textit{sāṅkhye}---[is based] on Sāṅkhya; 
  \textit{tu}---but [now]; 
  \textit{śṛṇu}---listen to; 
  \textit{imām}---this [teaching]; 
  \textit{yoge}---concerning [\textit{karma}]-\textit{yoga}; 
  \textit{pārtha}---O Pārtha (Arjuna); 
  \textit{yuktaḥ}---endowed with; 
  \textit{buddhyā}---[this] wisdom; 
  \textit{yayā}---by which; 
  \textit{prahāsyasi}---you will escape; 
  \textit{karma-bandham}---the bondage of \textit{karma}
}
{% --- 3. The Translation ---
  The knowledge which I have taught to you so far is based on Sāṅkhya. Now listen to the teaching concerning \textit{karma-yoga}. O Pārtha, by being endowed with this wisdom, you will escape the bondage of \textit{karma}.
}

In this verse here, you can see that Arjuna is progressing from one step to another. Kṛṣṇa sees that Arjuna is ready for the next step. He says, “I have been teaching you the yoga of the Self. I have just talked about how you will realize yourself: what you have to do, what you have to let go of, how you have to accept. Now I’ll talk about something greater, because I see that you have the attitude of equanimity inside of you. Now, we can move to the next teaching. Due to this attitude which you have cultivated, Arjuna, now we will talk about the yoga of equal attitude.”

Yoga means equanimity. Due to ignorance, one perceives the duality of this world, one perceives difference or diversity. But looking at it from the soul’s point of view, through Realisation, one gains true Knowledge. One will perceive no difference between one’s soul, the \textit{ātmā} and Paramātmā. One will see that one is just a part of Paramātmā, the Supreme Soul. The only difference which the soul will perceive, is in the quantity. But more than that, no. Because all the qualities which are in Paramātmā, are in the \textit{ātmā}. One who is completely absorbed in this Divine Self, in God Consciousness, finds this unity everywhere. There is no trace of ignorance. Kṛṣṇa says, “Hear Me! Now is the time to teach you how to attain this unity. Now I see that you have this interest and that you want to go deeper to possess this knowledge of the soul; and you are willing to do it. You have no conflicting feelings of heat and cold, or pleasure and pain. You don’t desire heaven and you don’t desire the kingdom. You are ready to know that, beyond all this, there is something greater.” Kṛṣṇa advises Arjuna, “This attitude of equanimity which you have, is your spiritual discipline. Now I start telling you about the \textit{jñāna-yoga} and the \textit{karma-yoga}.”

\Verse[2.40]
{% --- 1. The Verse ---
  nehābhikrama-nāśo ’sti pratyavāyo na vidyate \\
  svalpam apy asya dharmasya trāyate mahato bhayāt
}
{% --- 2. The Word-for-Word (Clean text, no formatting) ---
  \textit{iha}---here [on this path]; 
  \textit{na asti}---there is no; 
  \textit{abhikrama-nāśaḥ}---loss of invested effort; 
  \textit{na}---nor; 
  \textit{vidyate}---is there; 
  \textit{pratyavāyaḥ}---[any] failure; 
  \textit{api}---even; 
  \textit{svalpam}---a little; 
  \textit{asya}---of this; 
  \textit{dharmasya}---of this \textit{yoga}; 
  \textit{trāyate}---saves [one]; 
  \textit{mahataḥ}---from great; 
  \textit{bhayāt}---fear
}
{% --- 3. The Translation ---
  On this path, there is no loss of invested efforts nor is there any failure. Even a little practice of this \textit{yoga} saves one from great fear.
}

“On this path, there is no loss of invested efforts” This is \textit{karma-yoga}. Kṛṣṇa says that regarding whatever you do in this world – when you sow a seed in the land, there will be the result of the seed. Even if you fail to protect the seed, it will grow. So on this path of \textit{karma-yoga}, there is no loss of any effort. Nothing is lost.

“nor is there any failure. Even a little practice of this \textit{yoga} saves one from great fear.” Here Bhagavān Sri Kṛṣṇa says that if you do a little effort, if you have first shown interest, even just a little interest, and done a little practice – you have shown the willingness to advance. This will help bring liberation, salvation to your soul. “nor is there any failure,” There are no rules that will stop you. There are no rules for specific times, because salvation is beyond time. Kṛṣṇa says, “a little practice of this \textit{yoga} ” brings a big help.

The effort that you do in life has a great impact on your life. It delivers one from great fear, from not attaining one’s goal. Even if you are doing your duty, without even knowing why you are doing it, it will help in contributing to your advancement towards God Consciousness. It is said that when you start to do a little practice, there will be no problems. Because even if a test comes, you will not perceive it as a problem. You will embrace it and you will go through it. Like infinite waves in the ocean, countless waves of life and death; appear and disappear, you will move without fear; you will prevail. This is the knowledge that I want to give you.

\Verse[2.41]
{% --- 1. The Verse ---
  vyavasāyātmikā buddhir ekeha kuru-nandana \\
  bahu-śākhā hy anantāś ca buddhayo ’vyavasāyinām
}
{% --- 2. The Word-for-Word (Clean text, no formatting) ---
  \textit{iha}---here [on this path]; 
  \textit{vyavasāya-ātmikā}---[the] resolute; 
  \textit{buddhiḥ}---intellect; 
  \textit{ekā}---[is] one-pointed; 
  \textit{kuru-nandana}---O descendant of Kuru (Arjuna); 
  \textit{hi}---indeed; 
  \textit{buddhayaḥ}---the thoughts; 
  \textit{avyavasāyinām}---of the irresolute; 
  \textit{bahu-śākhāḥ}---[are] many-branched; 
  \textit{ca}---and; 
  \textit{anantāḥ}---endless
}
{% --- 3. The Translation ---
  On this path the resolute intellect is one-pointed, O descendant of Kuru; but the thoughts of the irresolute are many-branched and endless.
}

Bhagavān Kṛṣṇa says, “Because you are not attached to your actions, and you are not attached to the results of your actions, the mind becomes one-pointed. The intellect which determines only one thing, remains unshakeably fixed on it.” The mind becomes one-pointed on God and on God alone. Remaining steadfast in God, nothing can shake that person. 

Vyavasāyātmikā buddhir--The resolute mind is single-pointed. When the mind is settled for meditation, when the mind is calm and is focused on one goal, nothing can move its single-pointed devotion. The same is true when one surrenders to one spiritual path and looks only to that path; when one has found his way, his path in life, there’s no need to look elsewhere. Then one becomes strong. But if following one path, one still puts his nose in many places, he won’t move anywhere. Lord Kṛṣṇa says that “the resolute mind is one-pointed.” The one who has found his way, the one who has surrendered to his spiritual path, will find Satchitananda within, will find the Divine in one’s path itself, because the mind has fully given itself to this path.

Avyavasāyinām refers to those who do not possess this resoluteness, whose minds are diluted by feelings of diversity, and ignorance, and are deeply attached to enjoyment of the world. These people will never find peace, because they are always moving, shopping from left to right, jumping from one boat to the other, from one enjoyment to the other. They engage themselves in so many things that they can’t focus themselves anymore. Someone who does hundreds of mantras in their \textit{sādhana}, doesn’t get the benefit of the mantras, because each mantra vibrates differently. Whereas when you take one mantra, and focus on one mantra, it becomes more powerful. Then you have the full benefit of that mantra, which is much more than 100 mantras. Then the \textit{guru} Mantra or the Paramparam Mantra can help you.

Kṛṣṇa says that “ the irresolute are many-branched,” They don’t know this. They have so many things to do: today I’m here, tomorrow I’m there; I’m in this way, I’m in that way; I take this teaching from here, I take that teaching from there. Then there is a big mishmash of everything, a biryani of everything, and they don’t know anything. But the wise one is single-pointed. Even if there is a big mishmash of everything, the wise one will be able to take it and separate everything. This is yoga! Yoga doesn’t mean to sit down and do some exercises. To the western mind, yoga means physical stretching exercises, doesn’t it? No! That is why Atma Kriya Yoga is not based on what you do on the outside, but what you do on the inside.

Atma Kriya Yoga is to rediscover the \textit{ātmā}, focus on the \textit{ātmā}, focus on the Self, focus on the Divine within the Self. If you are focused on this Consciousness like this, then it’s much better. It will bear its fruits. It’s much stronger than having hundreds of paths which will not lead you anywhere. I always give this example – there are two boats next to each other. One boat says, “I am going to England.” The other boat says, “I am going to France.” You are standing on both boats and the boats start going. Sometimes you can continue on both boats together. But when the boats start going their separate ways, what will happen? Either you choose one boat or you fall down! This is what Kṛṣṇa is saying to Arjuna, “O Arjuna, when your mind becomes singlepointed towards God Consciousness, it’s not jumping on this and on that. When the mind is focused, you have control of the mind. You will not be able to destroy the mind, but you can control it. You can focus it! And that is yoga. Learn to focus the mind.”

The mind of Arjuna is not aiming to get something, like an object; it is not aiming to get other peoples’ recognition. Kṛṣṇa is very happy to see this in Arjuna, “Now you are ready for even greater things.” If somebody has this kind of determination, spirituality becomes very easy. If you have expectations on the spiritual path or on any path, it always brings sorrow. Because you expect, “Ah, I will get this,” “I will do that”, “I will do this”, “I will do this meditation and it will bear this fruit.” But when you don’t get what you expect, what happens? You become miserable. Then you don’t have any more faith.

Often people put many other things first in life. Sometimes they put love. Then other times it becomes hate. Sometimes they agree; sometimes they don’t agree. They are always moving from one thing to another. These people have ‘cañcala’ restlessness. Cañcala is a Hindi word describing somebody who is not stable, who is always flickering. You should not trust these people. Sri Kṛṣṇa sees this single minded devotion in Arjuna and He is very happy. He says, “This is good determination! Let’s strengthen this determination.”

\textit{Intro second commentary
Bhagavān is reminding that on this path whenever you are surrendered to the Feet of the Master, whenever you are doing your effort, you don’t lose anything. But you lose something when you are in ignorance of the mind. So, when you are in ignorance of the mind, fear arise, you know. Then you are never at peace.}

\textbf{2016}---So here Bhagavān is saying to Arjuna; those who have a certain determination in life, they are one pointed. They don’t jump looking here, there, have one foot here, another foot somewhere else. No, He said, if you want, if you have such resolution, you stand with both feet in one place. And when you have your both feet in one place, then you can achieve something. But if you have one foot here, another foot somewhere else, you will fall and drown yourself in delusion. So, one must have the will to follow the path of the right mind.

So, this mind must be clear about the value and objective which one wants to follow. That mind must be pure. So, a pure mind is calm like the surface of a lake which is in a calm state. That’s what Atma Kriya helps; brings to you to that state of mind where you become one single-pointed, where your devotion is singlepointed. And those who are not single-pointed, they dwell in the state of confusion. The mind is disturbed. They don’t have true value inside of them.

So, Bhagavān is saying to Arjuna, “Have singlepointed mindedness, you know. Don’t be in confusion of the mind, because the mind will tell you so many things, you know. Because you can’t – you have to learn to control the mind, but as long as the mind is not controlled, you keep jumping from one place to the another, one thought to the other. He said; ignore what comes out from your mind. Hold strongly to what is strong inside of you, you know. Hold strongly to that single-pointedness inside of yourself. Stay in one path which belongs to you. Don’t jump in different paths. When you take a path, when you surrender to one path, honour that path and give all your effort into that path. 10 people will come and tell you, “Oh, this is wrong. Our path – my path is better! Come!” If you are singlepointed, you know what is right. Nothing can move you. You see, you have something great inside of you, but yet, you try to look for it elsewhere through different kind of knowledge. You try to fill your mind with different things which will arise only confusion. Then after that you will say, “Oh, my goodness, why am I so confused?” Because you have not honoured your path. It’s like you have a bunch of keys which is inside of your pocket, but yet, you are searching for it everywhere. You know, you have a light, you hold it, and you go around and say, “Can I have a light?”, when you have that light yourself, when you just need to be aware of it.

So, on this path many will give you different kinds of knowledge, you know, how to free yourself. But know your path. That’s what Bhagavān Kṛṣṇa is saying here. Have this resolution of a single- pointed minded person. Let your mind be single-minded. Somebody will tell you; ‘yes, you are wrong’, but don’t listen to that person. Because the moment somebody tells you you are wrong for following that, that person itself is wrong. Due to his wrongness he can perceive wrong, so how can he guide you? Due to his ignorance he can just talk about liberation, but he can’t give it to you. “Do this, do that.” No, you have to do things which belong to you. So, that knowledge which belongs to you, it is your path which will lead you to your freedom. Doing your \textit{dharma} will lead you to your salvation. Nobody can bring that salvation to you. \textit{guru} will teach you, will push you, give you a kick from behind, but the first effort you have to, because that’s what \textit{guru} has inside of them, you know. They have this love for the Lord and they want to give it. But if you choose to be in deep sleep state, you know, or you want your mind be everywhere else, like Bhagavān is saying, “The thoughts of the irresolute are many-branched and endless.” So, when your mind is so much disturbed, you know, you can’t. Only when it becomes single-pointed, one-pointed, then it is easy for you. Why dig so many holes when you can just dig one hole and get the water?

\Verse[2.42/44]
{% --- 1. The Verse ---
  yām imāṁ puṣpitāṁ vācaṁ pravadanty avipaścitaḥ \\
  veda-vāda-ratāḥ pārtha nānyad astīti vādinaḥ \\
  kāmātmānaḥ svarga-parā janma-karma-phala-pradām \\
  kriyā-viśeṣa-bahulāṁ bhogaiśvarya-gatiṁ prati \\
  bhogaiśvarya-prasaktānāṁ tayāpahṛta-cetasām \\
  vyavasāyātmikā buddhiḥ samādhau na vidhīyate
}
{% --- 2. The Word-for-Word (Clean text, no formatting) ---
  \textit{pārtha}---O son of Pārtha (Arjuna); 
  \textit{avipaścitaḥ}---those who are unwise; 
  \textit{veda-vāda-ratāḥ}---delight in the words of the \textit{Vedas}; 
  \textit{vādinaḥ}---declaring; 
  \textit{iti}---thus; 
  \textit{asti}---there is; 
  \textit{na anyat}---nothing [superior to this]; 
  \textit{kāma-ātmānaḥ}---full of desires; 
  \textit{svarga-parāḥ}---[and] with heaven as the goal; 
  \textit{pravadanti}---[such people] utter; 
  \textit{puṣpitām}---flowery; 
  \textit{vācam}---speech; 
  \textit{kriyā-viśeṣa-bahulām}---full of various specific rites; 
  \textit{bhoga-aiśvarya-gatim}---aimed at enjoyment and power; 
  \textit{janma-karma-phala-pradām}---leading the results of birth and \textit{karma}; 
  \textit{prati}---toward; 
  \textit{yām}---which; 
  \textit{imām}---this
  \textit{buddhiḥ}---the intellect; 
  \textit{bhoga-aiśvarya-prasaktānām}---of those attached to pleasure and power; 
  \textit{apahṛta-cetasām}---[and] whose minds have been carried away; 
  \textit{tayā}---by that (flowery speech); 
  \textit{na vidhīyate}---cannot be established; 
  \textit{vyavasāya-ātmikā}---firmly; 
  \textit{samādhau}---in single-pointed focus
}
{% --- 3. The Translation ---
  O Pārtha, those who are unwise delight in the words of the \textit{Vedas}, declaring \enquote{there is nothing superior to this!} Full of desire and with heaven as their goal, such people utter flowery hymns about special rites to gain power and pleasure, which lead to further birth and \textit{karma}. The intellect of those who cling to pleasure and power and whose minds have been carried away by such flowery speech, cannot be firmly established in single-pointed focus.
}

Many people say that the world is unreal. These people who believe that the only goal of life is to attain Heaven are ignorant of the Truth. They are so attached to the teaching of the four Vedas, that they use this teaching for their own personal gain. Kṛṣṇa says, “People often talk, talk and talk using flowery speech. Don’t be bewildered when somebody just, ‘blah blah blah’. In this world, there are many wise and very knowledgeable people, but their knowledge is only based on books. They think that they know everything, just because they have mastered the books.” Here, Lord Kṛṣṇa says, “O Paartha, the words which these people are saying are unwise, because they don’t really know about the Vedas. They are just taking what they have read in books and then they put it in nice words. But if they don’t have devotion, if they don’t have love and surrender to God, this is nothing. Whatever these people do, they do it for their own selfish gain. They only desire to reach Heaven.”

Do you know the saying, “an empty drum makes lots of noise”? Lord Kṛṣṇa is referring to this. Many people can talk, talk, talk, but they don’t even know what they are talking about. They are very educated and intelligent, but they are not realized. They can give big lectures and use flowery speech. They can go on for hours talking, but their goal is completely different. Lord Kṛṣṇa says, “These people are unwise. Don’t listen to them. They do not have single-pointed devotion to God. They are just here to criticise others. They are here to point their fingers at others. This is due to their own ignorance and their own imbalance. This is because they don’t have full faith in their own religion. That’s why they use words to divert the mind. And people who have similar interests will follow their advice.”

“They teach rebirth as the result of actions and engage in various specific rites for the attainment of pleasure and power.” Here you see that these are very self-centred and egoistic motives. When people are not doing things for God-realization, but only to attain material pleasures, they are not really wise. They are ignorant. These people will always say that there is nothing higher than Heaven. They don’t even talk about God. They don’t even know God; they want only the fruits, the gifts which God will give them. But they don’t want God Himself. True knowledge is when you desire God Himself. You look beyond the pleasures of the world. You look beyond the reality of this world. You do your \textit{dharma} in this world, but you remain untouched by this world. This is what brings salvation and this is what men should strive for – to attain the Divine, attain the supreme goal of life.

When one becomes focused on realizing God Consciousness, one is free. But if one hangs onto these words and desires worldly enjoyments and power, then the mind is far away from God. One can’t attain the Divine. There are so many people in the world -- sometimes they can fill an entire stadium with people. But what are their minds focused on? What are their goals? People who have this kind of mind will always run to the people who speak flowery words and attach themselves to the world of enjoyment; and they will lose themselves. For them it’s very difficult to realize God. Due to their ignorance, they will always be looking at outside things. But one has to look for the Divine within oneself.

“Those who cling to pleasure and power are attracted by these teachings and are unable to develop the resolute will of a concentrated mind.” Those people who are always focused on the outside goals, will only attain the goals of the outside. But the one who is seated within himself through yoga, through \textit{sādhana}, through the love for God, will serve Him, love Him, develop this relationship with Him, first of all inside the heart. This one will be free.

\textbf{2016}---“Flowery speech,” Here Bhagavān starts saying, “O Pārtha, those who are unwise delight in the words of the Vedas, declaring “there is nothing superior to this!” Full of desire and with heaven as their goal.” Here, He is saying there are so many people who have knowledge of the mind, and they can say so many wonderful things to people. Here Bhagavān Kṛṣṇa said, no, they are not that knowledgeable. They are unknowledgeable people and they consider that by giving flowery, beautiful, nice, attractive, impressive speech, they are somebody.

That you see every day in life, no? All the politicians giving flowery speeches, because they tell people what they want, what the people want to hear. You want to hear beautiful things? They will tell you beautiful things. They will utilize the words of the Veda itself, they will utilize the content of the Veda, they will use the wisdom of the Veda, they will use the wisdom and try to reason with it, but does it bring happiness? Does it free and truly give peace to someone? No! It’s flowery words, you see? It’s beautiful, you sit there and listen, “Ah, such beautiful speech!” You go out, “What did he say?” Because it’s not related to you, these words that have been uttered. It doesn’t make any sense for you, because the person who was giving that speech doesn’t mean it. But when somebody who truly means it when giving a certain speech, puts his whole heart into it – it’s easy to go and tell somebody, “Go, do service, you will be free.” But here your mind is wandering around, how can you be free?

That’s why, here Bhagavān Kṛṣṇa is saying that these people’s aim is for heaven. They talk about heaven as if it is a place of permanent residence. He said, “No, it’s not permanent.” As hell is not permanent, heaven also is not permanent. The only one permanent place is to attain the Lord Himself. That’s why the bhakti path taught you to build your relationship with the Supreme Himself. Long for Him, surrender to the \textit{\textit{guru’s}} feet. That somebody else will not tell you. Somebody else will tell you words which you want to hear. You want flowery words? Well, I’ll give it to you. Normally, human beings, in their deep mind, think that they are always unhappy. In order to find this happiness, what do they do? They turn to comfort and pleasure of the outside. Then, to make themselves even happier, they have everything outside, they are attached to it, and they don’t want to let go of it. They don’t want to detach. Again, the word detachment here doesn’t mean that you have to throw everything out of the window. No, God has given you everything. Use it. But use it in a proper way. But many people don’t want to. They are so satisfied, they have built up their own little heaven thinking that this will last forever. And then they find, they take verses from the Veda which is according to their own nature. They interpret the scriptures to make themselves feel good, and say, “Yes, do all those good things, you will go to heaven.”

Here Bhagavān Kṛṣṇa said, “No, this is not permanent happiness. If you build a certain reality and believe in that, one day that reality will disappear. Even though you use many flowery words to prove to yourself that you are right, it will not be. Because you will still jump from one thought to the other, you will still make yourself feel good about something.

You see, you go to – Somebody says, “Yes, I’ve been doing good things in life and for sure when I die I will go to heaven.” It’s true. You do good things, you will have good merit. Like I said, it’s like a bank account, no? When you do bad – you have two bank accounts, but you don’t have the number.

One bank account you have downstairs, which is minus. One bank account you have upstairs, which is plus. I’m not too good in accounting, but let’s see. All the good things that you do in life, you receive merits. And, of course, these merits have to accumulate somewhere. Your bank account in heaven gets accumulated. It’s rising. And when you do bad things here on Earth, your bank account down is decreasing. Which bank account do you want to fill? The one up or the one down? Knowing that nothing is permanent, you will go to these places for a short-term until it is finished, and then you have to come back and work again. \textit{karma-yoga}, no? You are here to work.

That’s why it is said, out of the 7 billion people, how many truly attain the Lord? Very few. Because there are so many people who will tell you so many nice things just to please you, to please your mind, to make your mind happy. And you, of course, say, “Ah, he must be a very good person, he has told me things that I wanted to hear.” Yes or no? If somebody comes and tells you the truth in your face, would you be happy? As much as you say, “Yes, yes,” you don’t want to hear the truth.

You see, the \textit{guru} will not tell you things which will elevate your ego or your pride. He will tell you things which you don’t want to hear. This is how he crushes the pride out of you, he removes that ignorance out of you, not by flowery words, but by putting you into action, by making you work for that salvation. If somebody says, “Yes, don’t do anything, just sit around, I’ll give you heaven. You are baptized, you will go to heaven, don’t worry. Just sit around.” Many people think like that, “The Master will do everything for you.” The Master will not do it for you! He will show you the way. He will polish you, he will help you to control the mind, he will help you to conquer the senses, conquer the mind. He will help you to banish this unhappiness from you, so that automatically, through discrimination, through reflection, through analysis, through true knowledge of the Self, that true happiness awakens inside of you, this permanent happiness awakens inside of you. Otherwise, it would just be putting you in a big stadium, “Here, Hallelujah!”, you all will be saved. It doesn’t mean that you will be saved.

You have to make your effort. You have to purify yourself. You have to long for Him. You have to build that relationship with Him. You have to make your effort. You have to do that. Bhagavān is waiting, the Lord is waiting for you! That first step that you make, He will run to you. You see, in you, inside of you, when you develop this love for Him, there is a certain percentage, no? Imagine what is happening in Him, how much He longs for you. When you start walking toward Him, He doesn’t walk toward you, He runs toward you. That much He is longing for you. You think that you are longing for Him? No, you don’t! He longs for you more than that. So, He said, “Don’t get yourself trapped by these words from people who tell you things which you want to hear.”

Here Bhagavān carries on explaining, “Such people utter flowery hymns about special rites to gain power and pleasure, which lead to further birth and karma.”

Here many people center themselves for the attainment of personal gain, material things. You see, people are like that. They go to a saint, “Maharaj, I have this business, it’s not working properly, can you do something?”

One thing is very important to know; What is the nature of a saint? The nature of a saint is one with the Lord. Let’s analyze the saint now. A saint has left everything. He has left heaven itself to come and help people here, no? He is fully detached from everything. His aim is to raise you toward God, to raise you so that He can present God to you, “Here, I know Him, I present you, this is God, this is Kṛṣṇa, this is you. Here, build your relationship with Him. I’m giving Him to you.” This is the nature of a saint, where he can present God, take God and say, “Here, love Him. He is the Ultimate.” This is not the Ultimate, but He is the Supreme, and He is the eternal. It is with Him that you have to build your relationship, not with this world of illusion outside.

Here, people go and say, “Can you do something for my business?” Of course, in this world, you have all kinds of people, no? That’s why Bhagavān Kṛṣṇa in this verse said, they have flowery words which will - they are so engrossed in their own personal gain. They always want things for themselves. They are so selfish. Due to their selfish motive, when somebody comes, “What can I gain from that person?” That mind of certain people works in that way, so that they are wrapped up in their own imagination. They are wrapped up in their own dreams about a certain heaven. They create a fantasy of heaven, and they talk about it. They desire to be born in the next life in a very rich family. They are already planning for their next life. “How will my next life be?” Because of their knowledge only of the mind, they think that they are very superior. Because they can talk a lot, they think that they are somebody of a high pedigree. Here, Bhagavān is saying these are just words. Nothing can come out of them, because their motive is wrong. They don’t care for you, they care only about themselves. They will tell you, “Yes, yes, you are right,” for your sake, because you want to hear that, you want your desires to be fulfilled, desires which are only external, material, perishable.

They are rooted, their motive is selfishness, they are rooted in that motive of selfishness. They forget about the law of karma. They utilize the scriptures for their own benefit. And there are many of those, “Come ,come, I will give you a very long mantra and a big exercise you should perform. Next life, you will be well off. Forget about this life.” They are always thinking what will be next. They don’t think that this life can be the end of this

If you attain the Supreme, you attain Giridhariji, there is no need to come back here. But this, nobody can teach you apart from the \textit{guru}. He said, when you surrender to the Master’s Feet, the Master has detached from this world. This doesn’t mean anything. The Master brings the Lord to you. The Master reminds you of your relationship with the Supreme beyond this nature. He says, “Here, you want something permanent? He is permanent. I am giving Him to you. Why should I give you things which are not permanent?” Everything has an end here. Every creation, whatever is created, ends. Even if in your fantasy you think that this will last forever, it will not! Today is here, tomorrow it will not be here. And at the time of leaving this body, everything will be gone. So, what is your aim is very important. The Self doesn’t reveal Himself just like that. Mahāvatāra Bābājī has given the kriya for you to do. For what? For His concern? No! Look at Him, He is wearing only a dhoti, nothing else. What He wants is that you elevate, that you remind yourself of that Self, let the Self remind you of who you are. That’s your life purpose. The rest will be taken care of. Everything in nature is taken care of. It doesn’t mean that you should stop your duty, you should stop your work and say, “Okay, Swamiji has said, we have to go to heaven, now we don’t do anything. He will look after, everything will fall.” EFood will not fall down from heaven on your head, eh? For the saints and the sadhus, the Lord takes care. He does provide it. Also, sometimes it falls down.

But He said, “Do your duty, earn it!”, and He is ever present with you. As you are doing your duty, He is standing next to you. You know, for the devotees it is like that. They think they are alone doing their work, and say “I don’t have time to do the japa. I don’t have time to do this and that.” Once somebody asked, “I don’t have much time. I do one time japa in the morning, but during the day I’m so busy in my work. My mind is so focused on the work that I don’t have time for it.” But, you see, the moment you have become a devotee, the Lord is your permanent companion. Whether you feel Him or not, He is next to you. Whether you think of Him or not, He is next to you. This one time, when you think of Him truly from your heart, it builds this connection from your heart to Him.

“Such people utter flowery hymns about special rites to gain power and pleasure, which lead to further birth and karma. The intellect of those who cling to pleasure and power and whose minds have been carried away by such flowery speech,” here Bhagavān is saying, teaching, those who are egoistic, they will have egoistic motives. They will not teach you that you have to build a relationship with God. They will teach you how to attach yourself more to this world, because they are already attached. “I am damning myself, I want you also to damn yourself. I want to go to hell, and I bring you also together with me to hell.” That’s what the beautiful words are doing. They don’t teach you to be free, because they themselves, even if they are talking, their knowledge is only based on books. Through books, they can give big, big speeches. They can stand in front of millions of people, “Waya, waya, waya.” It’s just blah, blah, blah, nothing else. You see, if you speak, but don’t mean it, it doesn’t bring you anywhere.

“Those who cling to pleasure and power are attracted by these teachings and are unable to develop the resolute will of a concentrated mind.” He said, these people who are attached to these beautiful words, the minds of such persons, of people who think of high position, fame, glory, wealth, pleasures, comfort, they are so much preoccupied with these things: how to get more. They are so greedy. Their minds are so focused on the external things only, they can’t achieve the Lord. That’s why in spirituality, the first thing it tells you, “Calm the mind. Don’t let this mind control you. Show the mind who is the master.” If the mind is the master, then you will run outside. But if you show your mind that you are the master of the mind, you control the mind, it is you who allows the mind to think the way the mind wants to think – the way you want to think, and not run where the mind wants you. Because the mind is always attached only to the outside, and this attachment will bring you always – of course, when you have attachment, you know there will be another rebirth. And these people know, and say, “Okay, listen, you have a good life in this life, next life also you will have a wonderful life.” But not knowing what wonderful life is. Wonderful life, even if you don’t succeed in attaining the Lord in this life, next life to have a wonderful life is to still be a bhakta, still being a devotee of the Lord.

Here what is Lord Kṛṣṇa saying, “We have this eternal relationship. Once you are my devotee, I will look for you always. I will care for you and I will send you to the right place.” But if you start running toward the outside, by listening to some people telling you what you want to hear, such a mind cannot be trained. This mind cannot be controlled. That mind will always create a certain fantasy of what it is. They create a certain fantasy in their mind and say, “Oh, it is like that. It’s not me!” This is where mental problems arise.

Such beings, such people, they always point fingers toward the others. Here Bhagavān is saying that these people with such minds, who don’t have any control over their minds, they don’t have harmony in them. They can’t learn to control the mind. The mere mention of meditation will freak them out. The mere mention of yoga will freak them out. Here yoga doesn’t mean gymnastics. I’m not telling you how to put your feet on your head. Yoga is the union between you and the Lord. This is the greatest yoga; to realize this union between you and Bhagavān.

So, a mind, a free mind, a mind which is calm, it’s free from attaching itself to the outside. It is concentrating only on attaining this love and building this love for the Supreme.

\Verse[2.45]
{% --- 1. The Verse ---
  trai-guṇya-viṣayā vedā nistrai-guṇyo bhavārjuna \\
  nirdvandvo nitya-sattva-stho niryoga-kṣema ātmavān
}
{% --- 2. The Word-for-Word (Clean text, no formatting) ---
  \textit{vedāḥ}---the Vedas; 
  \textit{trai-guṇya-viṣayāḥ}---related to the three \textit{guṇas}; 
  \textit{arjuna}---O Arjuna; 
  \textit{bhava}---[you must] become; 
  \textit{nistrai-guṇyaḥ}---free from the three \textit{guṇas}; 
  \textit{nirdvandvaḥ}---[and] the pairs of opposites; 
  \textit{nitya-sattva-sthaḥ}---[be] ever established in \textit{sattva-guṇa}; 
  \textit{niryoga-kṣemaḥ}---free from concerns of gain and protection; 
  \textit{ātmavān}---[instead be] established in the \textit{ātmā}
}
{% --- 3. The Translation ---
  The \textit{Vedas} deal with the three \textit{guṇas}, O Arjuna. You must become free from these \textit{guṇas} and the pairs of opposites. Be ever established in \textit{sattva-guṇa}, free from concerns of gain and protection—instead, be established in the \textit{ātmā}.
}

“The Vedas deal with the three guṇas,” which are sattva, rajas, and tamas. “You must become free from these \textit{guṇas} and the pairs of opposites.” Here Lord Kṛṣṇa says that everything in this world functions because of the three \textit{guṇas}: sattva, rajas and tamas. Even the Vedas talk about these three \textit{guṇas}. But Kṛṣṇa says, “If you want to rise, you have to go above the three \textit{guṇas}.” If you are trapped in the game of the three \textit{guṇas}, you will always have judgment. The tamas is always judging the sattva and the sattva is always judging the rajas, and so on like that. Each one is thinking, “Oh, I am sattvic.” The tamas always hate the sattvic and rajasic.

Recently somebody was talking with a friend about some food and asked, “Can I eat this?” The other one said, “That’s very rajasic!” There is always a fight between rajas, tamas, and sattva. These three \textit{guṇas}, these three qualities, are the three qualities of the mind. It’s not the outside thing that we are eating or drinking. No! It’s the attitude of the mind, how one perceives things and how much one is hanging onto judgment. One is continuously judging, “This is like this and this is like that.”

“Be ever established in \textit{sattva-guṇa}.” Sattva is the state of the Self. Whoever abides in the Self, “never cares to acquire things or to protect what has been acquired, but is established in the Self.” The one who is free from this game of the three \textit{guṇas}, the one who has risen above these three \textit{guṇas}, is fully centred in the \textit{ātmā}. Here the word, \textit{ātmā}, means the Divine, God, Nārāyaṇa Kṛṣṇa Himself. As long as one has not controlled the mind, one is in the body consciousness. Here the Lord says, “Move from the body consciousness: focus, centre yourself, establish yourself in the Self.” If these three \textit{guṇas}, these three qualities are under control, one can rise above them. Know that God is above these qualities. And God is in all, everything: sattva, rajas and tamas. Slowly, slowly, Kṛṣṇa is bringing Arjuna out of his depression. He is revealing to him the \textit{ātmā}. Later on, He will reveal to him God Himself.

\textbf{2016}---“The Vedas deal with the three \textit{guṇas}.” Here He is teaching him. What are the three \textit{guṇas}? Sattva, rajas and tamas. He said that these three \textit{guṇas} have the characteristic to condition the mind of a person.

“O Arjuna. You must become free from these guṇas.” Here Bhagavān is not saying that you have to be sattvic, you have to be tamasic or you have to be rajasic. He said everybody is controlled by these three \textit{guṇas}. Everybody is born, this is their nature itself. This is the very nature of you being born in this world. You are born with certain characteristics, certain qualities. And due to those qualities you are brought where you have to be. But here He is saying to him, “No, you should not attach to this.” It’s not to put a placard on your face and say, “I am sattvic, don’t touch me!” “I am tamasic, keep away from me.” “I am rajasic, I get involved with everybody.”

The sattvic quality; “Be ever established in sattva-guṇa, free from concerns of gain and protection -- instead, be established in the ātmā.” Here He said everything. Those who have sattvic qualities, they aim for true knowledge, true wisdom, and they have a certain quality in life.

Rajas leads one toward ambition and desires. And tamas causes mental disorders, stress, laziness, hate and so on. You see, in life, when you are at peace with yourself, you can truly enjoy what the Lord is giving you. Then you really appreciate what you have. Here, appreciation doesn’t mean you attach to it. You appreciate it. You are grateful for what God has given you. When singing the Name of the Lord, you are participating in it. You are in deep – this inspires you. 

In your mind, you are saying, “It is difficult!” But it gives you a certain inspiration, no? The Name of the Lord, when you are singing, when you hear it, even, it’s different.

Rajas makes you always think of yourself. You are always running towards= name, fame, glory, “How can I be known? How can I make myself more important?” I, I, I, mine, mine, mine, this is rajasic. Not knowing who this “I” is. People of the rajas – everybody talks about “I,” also people who have sattvic qualities, tamasic qualities, rajasic qualities, they all talk about “I,” but not knowing who that “I” is. But we use words. It doesn’t mean that when you mean “I,” it’s me. It depends on the quality.

Sattvic people don’t think of themselves, but they think of others. “How can I make somebody else happy? Not for my sake, but for the sake of others.”

The rajasic quality will say, “How can I make myself happy by utilizing others?”

And the tamasic quality, mental stress, mental disorder. People of such mental illnesses, they are always thinking, “How can I harm others? I am not happy, so you should not be happy.” These are people which are born with demonic qualities. They build a certain understanding in their mind through drugs, through narcotics. They are full of this heaviness, laziness. They don’t have any aim in life, and because they don’t have any aim in life, they don’t want anybody else to have an aim in life. They project their depression all around them. These people, they don’t have any - they are far away. They still have to purify themselves so much.

That’s why Bhagavān is saying, different qualities, in which quality do you abide? But still, you have to go above these qualities. But certain qualities help you to be yourself. Tamasic qualities are far away. You see, if you say, “Okay, I have tamasic qualities, I will abide in that, I will realize myself in that.” No, you will not! You’ll have to climb that ladder step by step. That’s why we say that in God we have the divine qualities. Through the divine qualities, through the good qualities, we achieve Him. Through bad qualities, we are far away from Him. It doesn’t mean that God – God is everything. All the qualities are inside of Him. But to attain His Love, to realize His love, it’s only when you are pure, no? Your \textit{sādhana}, your prayer, you are longing for Him. These people who are in the outside world, they don’t. People who are taking drugs, creating a certain fantasy, and saying, “Yes, I have seen this.” When they take drugs, they say, “Aaaahh!” I don’t know, I never took it, but I wonder how it is, because I heard many people talk about it.

I asked somebody once, “Why do you take drugs?” For sure, many of you know about it. He said, “I feel good, and I see things.” That’s what you did, no? Doesn’t matter. You changed now, that’s a good thing. So, the person was telling me, “Yes, you are talking about God, but I am experiencing Him. I experience Him.” I said, “When did you experience Him?” He said, “When I take the drugs.” I said, “Now are you still experiencing Him?” He said, “No. Now it seems that He is far away.” Mental disorder! You create a fantasy, and you believe in that fantasy. But that fantasy stays a fantasy. Once that effect is gone, you are back to reality. Then you are back into depression. Where has that experience gone? Because once you experience the Love of God, it is ever-growing, it never leaves you. From time to time, Bhagavān will remind you of His Love, He is pulling the cord, that you come nearer to Him. But this is not through something. This happens through true Love. Something which is outside is limited. Whatever imagination or whatever illusion you experience through that, doesn’t last forever.

That’s why Bhagavān Kṛṣṇa is saying, center yourself in sattvic qualities because that will bring you closer to Him. Because if you center yourself into the duality – here He talks also about the pairs of opposites, no? You must free yourself from the three \textit{guṇas} and from all duality. Be above the \textit{guṇas}, but be above the pairs of opposites: hot and cold, failure and success, pleasure and pain, good and evil.

How to do that? Abide in pure sattvic qualities, abide in what is more important inside of you. He is exhorting Arjuna to know and to understand this jungle which we call life. So, learn to live with this yearning for pure qualities, learn to attain self-control. That’s why He said, go above all duality and have – Bhagavān is showing to Arjuna, that one cannot just let go of things like that if you don’t have the true knowledge of how to do it, if you are not truly guided. Otherwise, you will leave something and attach to something else. You see, when you leave something, you have to attach to something which is higher. Once I was talking with one person, she said, “Ah, I want to become a brahamcharini.” I said, “Okay, give me a good reason why you want to become brahamcharini?” It’s easy just to say: yes, it’s such a nice trend, let me put a title upon myself and take one title and stamp on myself, “I am brahmachari, I am renunciant.” “Give me a good example, why do you want to become brahmacharini.” “Ah, because I want to become a Swamini.” I said, “Go away from here!” To become a renunciant, you have to detach yourself from this world. You have to let go of that world that you have left behind, and you have to attach to Him and the Master. Not for the sake of something. Why should you leave something and attach again to something else? Just to praise your pride and ego? It will not lead you anywhere. To have a title, I am so and so, it will not bring you anywhere. If you are detaching, it must be for the right reason. And that right reason is to attain the Love, Supreme Love for God, nothing else. Whatever God will give you on the way, it is an extra bonus for you. But if you are renouncing something, but for something else again, it is wrong. And that’s what people do very often. “I want, I want, I want, I, I, I, mine, mine, mine, only me!” You will never be free, you see. Even if you have orange clothes upon yourself, it doesn’t mean anything if you don’t have this love and if you don’t understand what this life of a renunciant is, what the life of a brahmachari or brahmacharini is.

\Verse[2.46]
{% --- 1. The Verse ---
  yāvān artha udapāne sarvataḥ samplutodake \\
  tāvān sarveṣu vedeṣu brāhmaṇasya vijānataḥ
}
{% --- 2. The Word-for-Word (Clean text, no formatting) ---
  \textit{yāvān}---as much as; 
  \textit{arthaḥ}---utility; 
  \textit{udapāne}---[there is] in a well; 
  \textit{sampluta-udake}---flooded by water; 
  \textit{sarvataḥ}---from all sides; 
  \textit{tāvān}---that much [use]; 
  \textit{sarveṣu}---[there is] in all; 
  \textit{vedeṣu}---the \textit{Vedas}; 
  \textit{brāhmaṇasya}---for the \textit{brāhmaṇa}; 
  \textit{vijānataḥ}---who is realized
}
{% --- 3. The Translation ---
  As much use there is for a well when there is water everywhere, that much use there is in all the \textit{Vedas} for a realized \textit{brāhmaṇa}.
}

In this verse ‘brāhmaṇa’ doesn’t mean Brahmin as a caste: a brahmin here means a realized soul. A realized soul who is standing at the brink of a lake overflowing with sweet, life-giving and wholesome water, has no use for a small well. If you are standing by a whole ocean, why would you even want a small bottle of water? You already have the whole source. You don’t need the bottle since you have the source. What use would be a well, when there’s a flood all around it? What use would be all the water of a well, if you have a whole lake? Even so, a brahmin renounces every attachment to the objects of enjoyment and realizes God. That brahmin, that realized soul, when he has the Realisation of God Consciousness, when he has the Realisation of Nārāyaṇa Kṛṣṇa Himself, then even that ocean, which is called the Vedas, becomes of no use. Then one attains the supreme bliss. The Vedas are only talking about the Divine; but when you have the Divine Himself, what would you do with all this talking? You don’t need it, you have Him.

Kṛṣṇa says, “Those who attain this ocean of supreme bliss, in the form of Nārāyaṇa, don’t depend on any happiness or gratification, or on objects of enjoyment. One doesn’t even need any ritual which is recommended in the Vedas, because the desire for whatever is needed, is not there anymore.” There is no desire to get fulfilled. Everything has been fulfilled in the deep ocean of supreme bliss. To attain this state, one must let go, one must renounce all the desires, one must abandon all the desires for worldly enjoyment. All the other earthly things are the fruits of the Vedas, the rituals, and so on. When one is attached to any kind of thing, one is bound by the three \textit{guṇas}. You can take whatever you need from the Vedas which will benefit your advancement towards God-realization. But if it doesn’t help your advancement towards God-realization or help you attain the Grace of God, then it’s of no use.

\Verse[2.47]
{% --- 1. The Verse ---
  karmaṇy evādhikāras te mā phaleṣu kadācana \\
  mā karma-phala-hetur bhūr mā te saṅgo ’stv akarmaṇi
}
{% --- 2. The Word-for-Word (Clean text, no formatting) ---
  \textit{te}---yours; 
  \textit{eva}---[is] only; 
  \textit{adhikāraḥ}---the right; 
  \textit{karmaṇi}---to action; 
  \textit{mā kadācana}---not at any time; 
  \textit{phaleṣu}---to the fruits; 
  \textit{mā}---do not; 
  \textit{bhūḥ}---become (consider yourself); 
  \textit{karma-phala-hetuḥ}---the cause of actions' outcomes; 
  \textit{mā}---nor; 
  \textit{astu}---let there be [any]; 
  \textit{te}---your; 
  \textit{saṅgaḥ}---attachment; 
  \textit{akarmaṇi}---avoiding action
}
{% --- 3. The Translation ---
  You have only the right to act, but never to the fruits of action. Do not consider yourself the cause of actions' outcomes, nor let there be any attachment to avoiding action.
}

As I said earlier, if you work, if you do your duty, but your mind is only focused on what you want to gain, you will never be free. And that work will not be fully effective, because there will always be doubt about whether you will do it properly or not. There will always be doubt about whether you will achieve your goal or not. Because of this deep fear of unknowing, this fear of not doing it properly enough, you will not do it effectively. Here Kṛṣṇa says that you have the right to action, meaning you have to do your \textit{dharma}. But you only have the right to do your \textit{dharma} – the action, the work.

“Do not consider yourself the cause of actions’ outcomes.” Don’t be attached to the results, to the fruits of your actions. The result doesn’t depend on you, it depends on Him. In life it’s like this. You have a big plan to do something. And you plan it, maybe for days and for months, but when it happens, it happens the opposite way! What is your reaction? Even if the result is good, you are sad, because it’s not how you planned it. But if it is not good, you’re also sad. There’s no freedom in this.

In work, Kṛṣṇa says, “Do not consider yourself the cause of actions’ outcomes.” “You have the right to action, but only to action.” You have the right to action only. Kṛṣṇa gives a certain freedom here: the freedom of choice – a freedom that human beings alone have. When you are attached, you become ineffective, because when you are only focused on trying to do things properly, you are not free! You don’t have any freedom inside of you, because you are tense. Kṛṣṇa says, “If you do your action with the aim of serving the Lord, if you renounce this attachment to the result of the action, you will attain God-realization. You can’t renounce your actions, but do your actions with a selfless motive. Do them without egoism.”

If you try to renounce your actions by force, you will not succeed. Everything which is done by forcing, will not lead one to freedom. One will not be free, but will be under pressure, and always be thinking, “Oh my goodness, am I doing it right or wrong?” When you sit in meditation, you say, “When will I finish?” When you’re doing your japa, you say, “How many malas have I done? Where am I? Oops I missed one!” You’re sitting and doing your Atma Kriya Yoga and instead of enjoying the practice, you are thinking, “Which chakra am I on? Am I on the fifth chakra? No, I guess I was on the fourth one.” Isn’t it like this? You’re doing your pranayama, and thinking, “How many did I do? One more or one less?” You’re not free. But if you do your duty, if you do your \textit{sādhana} – just do it and be happy that you’re doing it! It doesn’t matter whether you breathe five or ten times. You will gain the fruit of the \textit{sādhana}. Your intention is very important. That is why Kṛṣṇa says, “Do not consider yourself the cause of actions’ outcomes, nor let there be any attachment to avoiding action.” By this He means that if you do your action being free, everything good will come to you. But if you do it with tension, nothing good will come to you. Then you will be disappointed and you will be miserable afterwards.

\textbf{2016}---Here Bhagavān is saying that you should do your work, but you should not attach to the fruit of the work. You have your duty. You have to focus the mind, you have to live, no? You have to work. But your work will truly bear its fruits when you do it out of love, when you love what God has given you, and happily do it knowing that you are serving Him. This is where work becomes worship. When you think of Him while you are working, then everything will be perfect, because it doesn’t come from you. It comes from Him. It is Him who is working through you. That’s why He said, “Don’t attach to these outside things what you are doing because you are limited, your power is limited, you have a limitation, but I don’t have any limitation. Hold tight upon me. I can give you more than that what you think. You are aiming only for something, but I can give you much more than that. Don’t attach to only these, because if you attach to only these, you are creating a certain limitation to it. Then you will get only that. Then you become a slave only to that, when there is much more you can do. What do you want? That's why He said to him, “Don’t attach yourself to the fruit of that action. Don’t look at this limitation. Don’t make that your motive. Know that it is I who is doing everything.”

Like Tukārām was always saying, “Lord it is not me, it is You who is doing through me. It is only Vittala Panduranga which is doing everything.”

In life, you have certain choices. For example, if you have to make a certain decision and you have to achieve a particular objective - let’s say you want to learn something, you have different ways, no? You start searching, you choose. You choose – you have three ways, you have the tamasic, rajasic and sattvic way. You want to choose a course because you have a certain motive, what is the motive for it? Then you have these choices in front of you that you have to follow. Here Bhagavān is telling to Arjuna; Look, you have the tamasic way, rajasic way and sattvic way, and you have your motive. If you follow this, it will lead there, if you follow that, it will lead somewhere else. If you are following the sattvic, you will reach me. What is your choice? You have this intellect, no? Through your intellect, you have to choose wisely. Which course do you want to do?

You see, when you do the course, do you know the result, what it will be? When you are sitting for an exam, do you know what the result will be? You don’t know, no? Is it in your hands? No, it is not. What will be, it will be. That’s why Bhagavān is saying, don’t attach yourself to that fruit of that gain. Because that fruit of your action will not come to you per your demand.

He says that it is not you which controls everything. Let not the fruit of your action be your motive, because it’s not you which have control over that. It is the Divine law which holds good under all conditions of action, inaction or non-action. It is this divine law that gives you these things.

To consider that – “Since the result of my \textit{karma} is not in my hands, why should I do something then?” By thinking like that, it is already a negative attitude. That’s why Bhagavān said, inactivity. You are born as a human being, you are always in action. If you say, “No, I don’t want to do something,” what will happen? You will attach yourself to laziness, incapacity. You will be incapable of doing things. You will become lazy. And this is tamasic. So, a person who is lazy, do you think they are happy? Do you think they are happy to be lazy? But they don’t have any choice. They have taken that path.

Very often, people come to me and say, “My son is lazy, not doing anything, just closing themselves in their room.” Do you close your child in the room? No. You see, what you have taught them, that’s what they are doing. You have allowed them, you have not given them the right knowledge from the beginning, and then you sit and cry, “Oh, my son is like this, my son is like that. My child is like this, not doing anything.” Who is responsible for that? You, the parents! If from the beginning you would have given them the right conduct of living, respect, and if you yourself would have that knowledge, then you could give it, no? But it’s never too late. Here Bhagavān said, “Look, one is given opportunities, and you should never think that you are left to damn yourself, you are finished, that you can’t be saved. The moment you make your first step toward Him, He makes everything become active. Even when you are inactive, you are lazy, even from your laziness, He will bring you out. What do you have to do? Not sit and cry, “Oh my child is lazy, he is closing himself!” No. Offer this laziness to the Feet of the Lord. Offer it, offer your child to the care of Thakurji. He will look after, and He has His way of handling everybody, even the lazy people.

\Verse[2.48]
{% --- 1. The Verse ---
  yoga-sthaḥ kuru karmāṇi saṅgaṁ tyaktvā dhanañjaya \\
  siddhy-asiddhyoḥ samo bhūtvā samatvaṁ yoga ucyate
}
{% --- 2. The Word-for-Word (Clean text, no formatting) ---
  \textit{dhanañjaya}---O Dhanañjaya (Arjuna); 
  \textit{yoga-sthaḥ}---established in \textit{yoga}; 
  \textit{tyaktvā}---[and] abandoning; 
  \textit{saṅgam}---attachment; 
  \textit{kuru}---perform; 
  \textit{karmāṇi}---actions; 
  \textit{bhūtvā}---having become; 
  \textit{samaḥ}---equal; 
  \textit{siddhi-asiddhyoḥ}---in success and failure; 
  \textit{samatvam}---[this] equanimity; 
  \textit{ucyate}---is called; 
  \textit{yogaḥ}---\textit{yoga}
}
{% --- 3. The Translation ---
  O Dhanañjaya, established in \textit{yoga} and abandoning attachment, perform action, remaining equal in success and failure. This equanimity is called \textit{yoga}.
}

Here Kṛṣṇa is saying, “O Dhanañjaya, established in \textit{yoga} and abandoning attachment, perform action.”In this verse Kṛṣṇa says that when you do yoga, it will help you to rise above expectations. Yoga will unite you with the Divine. Whatever you do is yoga, when you’re not attached to the fruits of the action. But whatever you do with expectation, will never free you. You will damn yourself, you will get yourself attached to the world. When a sadhak who is practising \textit{karma-yoga} lets go of attachment to his own actions and the fruits of those actions, when he ceases to have likes and dislikes, then he is free from joy and sorrow. The essence of equality awakens and then one attains the Divine. As long as there is judgment, this is not possible.

Through the constant practice of indifference to success and failure in action, a human being reaches the ultimate state of unshakeable stability in equality. This is the culmination of \textit{karma-yoga}. Therefore, Bhagavān Kṛṣṇa tells Arjuna to perform his duty established in yoga. Kṛṣṇa is trying to remind Arjuna that the practice of equality in failure or success is important. You have to be equal in both failure and success. You have to have equanimity in the performance of your daily duties in your every act, by being free from like and dislike, with reference to every object, every being, every fruit of that action itself. Kṛṣṇa says that this is the essence of yoga: to free yourself, to rise above attachment to the fruits of your actions. Yoga is this equanimity. He shows that one can become a \textit{yogī} and attain this equanimity, through any discipline, or through any \textit{sādhana} whatsoever. That’s why He says to Arjuna, “Establish yourself in this equality. Establish yourself in that yoga. Become a \textit{yogī}.”

\textbf{2016}---Neutrality. Bhagavān has brought Arjuna in between the two armies because of this equality, neutrality. He is neither on the good side, nor on the bad side. He dwells in the middle. When Arjuna told Him, “Move the chariot to the middle so that I can see”, Bhagavān already said, “Ah, that’s a clever move. It’s not a bad move to do, because you are not on either side. But first, I have to give you the right knowledge.” So, you see how He planned everything.

The \textit{yogīs} do their own action without expecting anything in return, and they are not attached to anything. He is saying to him right now, this is the essence of \textit{karma-yoga} itself, “Fixed in Yoga, do your actions, make work a prayer. Pray and work.” You see, Arjuna was so disturbed in the mind, “How can I do that?” Here Bhagavān is saying, “Yes, you can!” Concentration, focus. Yoga means reminding yourself of your union with the Divine, your connection with Him. Whatever you do with that connection with the Divine, everything is blessed.

“Abandoning attachment, perform action, remaining equal in success and failure. This equanimity is called yoga.” Take success and failure in the same way. Purify the mind. Only when the mind is purified, only when the mind is calm, only when the mind is focused, is centered on the Supreme Lord, then you are free. If the mind is wandering in all directions, nothing good can come out of it. You are never sure, “Am I right or not?” The mind makes you say, “Oh, I want to achieve this. I can do it, you can do it. Why? How?” You have hundreds of questions that arise inside of your mind. This is an uncontrolled mind.

That’s why when you do your \textit{sādhana}, I said, firstly try to calm your mind. What awakens in the first instinct, first feeling, first wave which awakens inside of you, that’s what you have to learn to listen to, because it comes from your consciousness, it comes from your higher Self. If you can distinguish between these two, you will always make the right decision.

The mind always runs toward the pairs of opposites. It always wants to justify itself. So, control that mind. By a controlled mind which is not running left-right in all directions, it will give you balance and equilibrium in life. This is what yoga stands for, when the mind is controlled and centered, when the mind is calm, not only when you are doing your \textit{sādhana}, but at all times.

You see, very often you see people, you talk with people that say, “Yes, I have been doing yoga for twenty years.” Once, I met one lady, we were sitting at the table, she said, “I’ve been doing yoga for twenty years.” I looked at that person, I was wondering, “Do you know what yoga means?” Yoga doesn’t mean what you have been doing on the outside. That’s not the true yoga. You see, every action is a yoga itself. Everything which is active is yoga, whatever you are doing. But true yoga is this union between the \textit{ātmā} and Paramātmā, this relationship, if you realize that connection. I see all of you doing more yoga than that person who has been doing twenty years of yoga, because that person doesn't have any understanding of what yoga is. She has been doing exercise to keep herself fit, which is good, it’s important also. But often you see that, the people that do Hatha yoga, they think that this is the ultimate, that’s everything. No, they are still empty. That’s not everything. You see, Patanjali has created these exercises for you to keep your body healthy, to keep yourself proper and to keep the mind centered.

What’s the use if your yoga is only in the class, you are doing the exercises in the class, and then the moment you go out, again you are the same? That’s what Bhagavān is saying, that equanimity, this discipline, this control, should always be with you at all times. It’s not that you say, “I am doing yoga for so many years,” and then when you look at yourself, you are not at peace, you are stressed. No, what are you doing? Have you understood what yoga is? No, you have not! So, the right knowledge is important.


\Verse[2.49]
{% --- 1. The Verse ---
  dūreṇa hy avaraṁ \textit{karma} buddhi-yogād dhanañjaya \\
  buddhau śaraṇam anviccha kṛpaṇāḥ phala-hetavaḥ
}
{% --- 2. The Word-for-Word (Clean text, no formatting) ---
  \textit{dhanañjaya}---O Dhanañjaya (Arjuna); 
  \textit{karma}---action (with attachment); 
  \textit{hi}---[is] indeed; 
  \textit{dūreṇa}---by far; 
  \textit{avaram}---inferior; 
  \textit{buddhi-yogāt}---to [this] \textit{yoga} of wisdom; 
  \textit{anviccha}---[so] seek; 
  \textit{śaraṇam}---refuge; 
  \textit{buddhau}---in wisdom; 
  \textit{kṛpaṇāḥ}---those who are the petty-minded; 
  \textit{phala-hetavaḥ}---are motivated by the fruits of action
}
{% --- 3. The Translation ---
  O Dhanañjaya, action done with attachment is greatly inferior to this \textit{yoga} of wisdom, so seek refuge in that wisdom. The petty-minded are motivated by the fruits of action.
}

“Action done with attachment is greatly inferior to this \textit{yoga} of wisdom.” Here Kṛṣṇa says that when one has selfish motives, when one is doing things in a very selfish way, one is devoid of this knowledge of yoga or equality. But when one doesn’t concentrate on the fruit of the action, one attains the Divine. In this context, Kṛṣṇa is saying that whatever you do, do it with a selfless motive -- because if you do it with a motive of attaining something, “The petty-minded are motivated by the fruits of action.” This is the pride. This is the ego. People will always claim that they have done this or they have done that and brag about it saying, “Oh, we have done this”, “We have done that.” Even if this action appears good, it is not good. Then they start comparing, “If I do that much, I will get that much.” They work like normal people who are doing \textit{karma-yoga} in the outside world. It will not free them.

Kṛṣṇa says we should pity people with such minds because they will damn themselves. He says, “O Arjuna, instead of focusing on that, gain the true knowledge of yoga and surrender.” Because when you have the true knowledge of yoga, God reveals Himself; this Love awakens, because you go through this purification. Then you let go of what is ‘mine’ and ‘I’. The ego is not present, saying, “I am doing this”, “I am doing that.” These evil ways of thinking will disappear. He is advising Arjuna, “If you go into this wretched state of mind and hang onto these objects of action, it’s a pity.” He advises Arjuna not to be poor and wretched like these people and says, “You are wise.”

\Verse[2.50]
{% --- 1. The Verse ---
  buddhi-yukto jahātīha ubhe sukṛta-duṣkṛte \\
  tasmād yogāya yujyasva yogaḥ karmasu kauśalam\\
  karma-jaṁ buddhi-yuktā hi phalaṁ tyaktvā manīṣiṇaḥ \\
  janma-bandha-vinirmuktāḥ padaṁ gacchanty anāmayam
}
{% --- 2. The Word-for-Word (Clean text, no formatting) ---
  \textit{buddhi-yuktaḥ}---united with wisdom; 
  \textit{jahāti}---one rises above; 
  \textit{ubhe}---both; 
  \textit{sukṛta-duṣkṛte}---righteous and unrighteous deeds; 
  \textit{iha}---here in this life; 
  \textit{tasmāt}---therefore; 
  \textit{yujyasva}---engage yourself in; 
  \textit{yogāya}---\textit{yoga}; 
  \textit{yogaḥ}---[this] \textit{yoga}; 
  \textit{kauśalam}---[is] expertise; 
  \textit{karmasu}---in actions
\textit{manīṣiṇaḥ}---the sages; 
  \textit{buddhi-yuktāḥ}---[who are] established in this wisdom; 
  \textit{hi}---indeed; 
  \textit{tyaktvā}---renounce; 
  \textit{phalam}---the fruits; 
  \textit{karma-jam}---born of actions; 
  \textit{gacchanti}---attain; 
  \textit{padam}---[the] state; 
  \textit{anāmayam}---beyond all suffering; 
  \textit{vinirmuktāḥ}---[since they are] freed from; 
  \textit{janma-bandha}---bondage of rebirth
}
{% --- 3. The Translation ---
  United with such wisdom, one rises above both righteous and unrighteous deeds while living. Therefore, engage yourself in this \textit{yoga}. This \textit{yoga} is expertise in action. The sages who are established in this wisdom renounce the fruits of action and attain the state beyond all suffering, and are freed from the bondage of rebirth.
}

Here, in verse 50, Kṛṣṇa says that, “United with such wisdom, one rises above both righteous and unrighteous deeds while living.” A true \textit{yogī} is not attached to anything. What he receives, he can give away, without having any remorse. He is free. He does this because he is established in this equanimity. The residue of all virtues and sinful deeds from one’s past lives, persisting in the form of tendencies, are stored in the Brahma Nadi. A true \textit{yogī} sees this and lets go of it. He removes all these tendencies; he renounces them from the mind itself; he breaks the connection with these actions which he has performed in this life, or in a past life. From the mind itself, he removes them. He doesn’t just sit and enjoy this. When he does something good, he doesn’t sit there and feel good about it and keep praising himself. He removes even that feeling.

It doesn’t make any sense to just love the people you already love. You should learn to go beyond that and love everyone. You have to come to the point of this equalmindedness. Like the \textit{yogī} who sees all in the mind, who sees the connection with his action in the mind itself, from the beginning. Whatever arises inside the mind, how much attention do you give to it? Kṛṣṇa says, from that moment on, you remove it from your mind so it cannot bear fruits in the form of rebirth. If you renounce all fruits of action from the mind itself, the mind is purified. There is no creation of any new \textit{karma} for later on. You start to burn it! That’s what Atma Kriya Yoga does: it burns the karma, releasing you from past tendencies which had been stored inside you. Through the performance of this disinterested action for the good of the world, all the actions of the \textit{yogīs} are neutralized. With this attitude, there is no attachment to personal gain. 

When you help people, there is always an expectation – for a ‘thank you’, for a gratification, for praises. You see this in daily life, you see it everywhere; wherever you go, the world runs on expectation. When you don’t expect anything from the mind itself, all is purified, all is transformed; you stop this creation of karma. Similarly, the virtues and sinful deeds of this current life will also stop having any effect. Lord Kṛṣṇa is advising Arjuna to commit to exerting himself through the yoga of equanimity. The Lord reminds the mind of Arjuna, that such a \textit{yogī} gets liberated in this very life itself. He says that you can liberate yourself, now and forever, in this very life. How? You have to transform; you have to do all your actions without any expectation – and if that is done, you are free! Otherwise you will become a slave of the mind, and you will become as crazy as the mind.

In verse 48, Kṛṣṇa tells Arjuna, “Fixed in yoga, do your actions.” Every action which one does in life, whether it’s big or small, always creates karma, because the very nature of action leads to bondage, to attachment. People can’t really stay inactive, even for a moment. They are always busy. The mind is always busy engaging itself, jumping from one thought to another, going from one action another. It’s not possible to not do action. Even when you are sleeping, you are doing action: you are breathing in and out, you are moving left and right. So there is action in every aspect of your life. Kṛṣṇa says in such circumstances, “Practice the yoga of equanimity.” If you train your mind to see this equality, it is very easy. And it’s the best practice to relieve one from bad karma, to relieve one from the bondage of karma. Therefore, practice this and become equal in failure and success. For it is equality that is the essence of yoga. It’s only through this equalmindedness that one will attain the state of perfection.

“The sages who are established in this wisdom renounce the fruits of action and attain the state beyond all suffering, and are freed from the bondage of rebirth.” So Kṛṣṇa says that the sages “who are established in this wisdom ” They are not attached to the mind which is thinking always of, ‘I’ and ‘mine’. Through their wisdom, the truly wise and learned are established in the state of equalmindedness, and they have attained the goal of their human existence: they have become sages, they have become \textit{yogīs}.

A human birth is the doorway to salvation. Here Lord Kṛṣṇa says, “One who attains a human birth should realize how important it is.” This human birth is not here to create karma. You have not been given this human body to dance to the tune of Maya. You have not been given this human body to enjoy an attachment to this world. You are not here to be bound yourself to this reality. You are here to elevate yourself. You are here to liberate yourself, to detach yourself, to rise to the divine’s virtues, to God Consciousness. You are not here to rebuild karma, but to be free from karma, to be free from misery. Because it is through the creation of \textit{karma} -- firstly in the mind – that misery arises. But you make yourself miserable! Nobody is responsible for your misery! You yourself are responsible for your misery, because in your mind you choose to be miserable, you choose to be poor, you choose to be sad, depressed, terrible, and so on.

Kṛṣṇa is saying to Arjuna, “Rise, My dear! Be wise! If you are wise, you will not attach yourself to all these human qualities and miss the opportunity in this very life. If you miss this opportunity and remain immersed in worldly enjoyment, then you are not wise, you are like the outside world.” Through the stability of the yoga of equal-mindedness, one ceases to have any connection with the fruit of action performed in this current life, as well as with the fruit of action performed in countless of past lives. This automatically liberates one from the cycle of birth and death. One is released from this bondage when one renounces the fruit occurring from action.

Kṛṣṇa says that the sage who is liberated from misery, who is liberated from all the pain of Maya, will automatically obtain the Supreme Abode of God. If one is freed from the mind, liberation is automatically attained. One attains the Lotus Feet of the Lord and one is totally surrendered. One rises, one goes beyond the power of Nature itself and the divine qualities start to manifest through that person. The one who is totally above all the negative qualities, becomes identical with the Supreme Brahman, with Nārāyaṇa. Then the qualities of Nārāyaṇa start to manifest through that person. With the attainment of the Supreme Abode of Lord Nārāyaṇa, with the Realisation of the supreme state of Brahman, one awakens and Satchitananda manifests through that person. All the attributes of the Divine start to shine through this person. And one goes even beyond the reality where Heaven is present: this state is beyond Heaven itself. One goes beyond the dualistic state in the mind: and this breaks all the bondage to birth and death. When the bondage to birth and death is broken, one rises above even the heavens. One becomes a true \textit{yogī}.

\textbf{2016}---“United with such wisdom, one rises above both righteous and unrighteous deeds while living.“ A \textit{yogī}, those who practice yoga, the man of wisdom, somebody who has true knowledge, is always harmonious.

“Therefore, engage yourself in this yoga.” Devote yourself to this realization. “This yoga is expertise in action.” Here Bhagavān Kṛṣṇa is saying it; the performance of action with a balanced mind, whatever you do with a balanced mind, you are not disturbed.

“The sages who are established in this wisdom renounce the fruits of action and attain the state beyond all suffering, and are freed from the bondage of rebirth.” He said, a \textit{yogī} who has everything in balance is not prone to disturbance caused by emotional upset. A true \textit{yogī} is not attached to anything. A true \textit{yogī} is – you see, \textit{yogīs} are people who perform yoga. Because the sage, the \textit{yogī}, has the knowledge and the understanding of what is happening in his daily life. This understanding is deep enough and can’t be shaken - they are well-rooted. This flame of knowledge, the flame of true knowledge, has enlightened his whole inner Self. All the tendencies which the mind creates to attach itself, he has broken all the connections with that. He goes, but he is not attached to. He does what he has to do, but he is not touched by that.

The sages who are detached don’t long for wealth and comfort on the outside. For them, they long only for the Divine Grace. There is only one thing they desire: only the Grace of God. The true \textit{yogī} longs only for that Grace of God. That’s what matters. All actions, all deeds performed, are just objectives to attain that Grace.

“This yoga is expertise in action.“ You see, he cares for nothing; he has become indifferent to the very nature of action itself. He is not attached to that action, whatever he is doing. He doesn’t have any attachment or tendencies; he is not attached to that, he doesn’t have any expectation. In a \textit{yogī}, somebody who performs yoga, the knowledge and action have blended together, and this true knowledge about the Self starts to reflect upon him.

You see, when you are born on Earth, you know you have one thing that you want. Why you have come here is to free yourself, to have salvation. But due to the association with the outside, very often you are fed with the wrong things.

Birth is not just happening like that. Each soul is born to attain the Divine. When wisdom arises, when the ignorance is dissolved, when Bhagavān calls you on the spiritual path, this is where the Light which is inside of you starts to shine. Then you are truly living. Then He makes everything possible so that you are free from Maya, so that you can come and abide with Him. When the mind is free, liberation runs toward you. You don’t need to run toward liberation. Everybody is doing things for something, no? You are doing something to attain something. Here Bhagavān is saying, no, if you long for God, you don’t need to run behind something. Everything will run towards you. People who lack wisdom, they long for liberation without even knowing what liberation is. Then they have a fantasy of that liberation. Then you will create that fantasy. Here Bhagavān said, “No, you should not. The only thing is to attain that Grace, and reside with Me in My Supreme abode eternally. This is where you come from, and this is where you will go.” Before creation, you lie inside of Him. When you were created, you became external to Him. When you leave, when you attain Him, you will still be near to Him. You take certain aspects. Those who have attained the Lord, they take the same aspect of the Lord; they resemble Him. It doesn’t mean that they become Him, but they start to look like Him.

What does the fruit of action matter for a wise sage who has realized himself? Those who have realized themselves, they know what is transient and what is eternal. They know wisely that bonding, that pulls oneself here, what keeps oneself here, what delusion is, what confusion is. When you know what delusion and confusion are, do you like to go into it? So, here I’m telling you.

Long for the Grace of the Lord, long for Him, don’t let yourself be spellbound by Maya. Because those who are centered into the Lord, Maya doesn’t have any grip upon them. Those who are fully centered, the \textit{yogīs}, they are free. They go everywhere. They are eternal. They radiate their peace, their freedom, their true happiness. Without any expectation, they radiate this freely to all, and they become a lantern to those who are seated in darkness. They become a guide to those who sit in darkness so that they can also start shining their light. And this is the work of a devotee, the work of a bhakta; once they have tasted that sweetness of the Love of God, once they are strong enough, they should - this is your duty, to help others find that Light.

\Verse[2.52]
{% --- 1. The Verse ---
  yadā te moha-kalilaṁ buddhir vyatitariṣyati \\
  tadā gantāsi nirvedaṁ śrotavyasya śrutasya ca
}
{% --- 2. The Word-for-Word (Clean text, no formatting) ---
  \textit{yadā}---when; 
  \textit{te}---your; 
  \textit{buddhiḥ}---intellect; 
  \textit{vyatitariṣyati}---crosses beyond; 
  \textit{moha-kalilam}---[the] forest of delusion; 
  \textit{tadā}---then; 
  \textit{gantāsi}---you will attain; 
  \textit{nirvedam}---indifference; 
  \textit{śrutasya}---to what has been heard; 
  \textit{ca}---and; 
  \textit{śrotavyasya}---to what is yet to be heard
}
{% --- 3. The Translation ---
  When your intellect crosses beyond the forest of delusion, you will be indifferent to what you have heard and what you will hear in the future.
}

Here Lord Kṛṣṇa says that when all this has been eradicated from your mind, when all the negative qualities have been discharged, when your mind has been purified of all impurities, and when your intellect is focused on the Divine, then you realize that there is no ‘good’ or ‘bad’. Even when you hear somebody flatter you or not flatter you, it doesn’t make any difference.

In this verse, Lord Kṛṣṇa refers to what the scriptures say you should do or should not do. In the scriptures there are always many things that tell you – you have to do this and you have not to do that.

Kṛṣṇa says that when your mind is completely focused on the Divine, what you read in the scriptures or what you hear from the scriptures, doesn’t make any difference to you. You perceive the Divine everywhere, not only in saintly people, but in everyone – wherever you go and in whatever duty you do. And you do your duty with such devotion, you do it with such eagerness, solely to serve and please God. You renounce any merit that the mind is thinking it may gain by serving the Lord. Some people say, “We pray, we pray, we love,” yet behind this there is always a certain expectation.

In this verse, the Lord tells Arjuna that after the disappearance of attachment, his mind would regain its natural state of transparency. He would then develop genuine indifference towards all, towards the ephemeral objects of this world and the next.

\Verse[2.53]
{% --- 1. The Verse ---
  śruti-vipratipannā te yadā sthāsyati niścalā \\
  samādhāv acalā buddhis tadā yogam avāpsyasi
}
{% --- 2. The Word-for-Word (Clean text, no formatting) ---
  \textit{yadā}---when; 
  \textit{te}---your; 
  \textit{buddhiḥ}---intellect; 
  \textit{śruti-vipratipannā}---[that has been] perplexed by the words of the \textit{Vedas}; 
  \textit{sthāsyati}---remains; 
  \textit{acalā}---motionless; 
  \textit{niścalā}---[and] unshaken; 
  \textit{samādhau}---in single-mindedness; 
  \textit{tadā}---then; 
  \textit{avāpsyasi}---you will attain; 
  \textit{yogam}---\textit{yoga}
}
{% --- 3. The Translation ---
  When your intellect, confused by the words of the \textit{Vedas} remains unshaken in single-mindedness, then you will have attained the state of \textit{yoga}.
}

In this state of yoga, the intellect has crossed the mirror of delusion and fully recoiled from the enjoyments of this world and the next. You are not bothered about this life, or the next. You have wholly freed yourself from the fault of distraction, you are fully concentrated, and you take to the practice of meditation on God.

Kṛṣṇa says, “When your mind is calm, when your mind is not attached to the outside, let your mind be focused on God.” He doesn’t say, “When your mind is controlled, just leave it like that.” He says that the mind must be focused on the Divine form of God, the Divine form of Nārāyaṇa Himself. So practice meditation on God and be unshakably and firmly concentrated on God alone.

When your mind has changed its state and its way of viewing life, it must be focused on something new. Otherwise the mind will jump back into the past activities. Bhagavān Kṛṣṇa says, “Focus your mind on God.” It’s like a child: when a child is naughty, you have to give the child a toy to play with so he can entertain himself. This mind is like a child who is naughty. You have to make the child be very good; you talk nicely with the child, but you also have to give the child a certain instrument. And the instrument is this deep concentration, this meditation on God alone. This is where the mind must rest, steady and undistracted: the mind “firmly placed, remains unshaken in a concentrated mind, then you will attain the vision of the Self, and attain yoga, tadā yogam avāpsyasi.”

It is through that yoga that you will attain this perpetual and complete union with God. This union will only happen when the mind is free from the three \textit{guṇas}, which are impurity, distraction, and obscurity. When you are endowed with the power, and energy of discrimination and dispassion, then you are fully concentrating on God.

These are the seven types of yoga related to the yoga of the mind: yoga of action, yoga of meditation, yoga of even-mindedness, yoga of the Divine glory, yoga of devotion, Ashtanga Yoga consisting of the eight limbs, and Sankhya Yoga.

The first yoga is the yoga of action, \textit{karma-yoga}. Here the action is performed by one who seeks to climb to the highest level of yoga. It is the action of making the effort to attain the union between the \textit{ātmā} and the Paramātmā and realizing the oneness in this equalmindedness.

The second yoga is the yoga of meditation, Dhyaana Yoga. The mind of the \textit{yogī} is as steady as the flame of a lamp kept in a windless place (Ch. 6, v. 19). When the mind enters this state of steadiness and is not flickering, not jumping left and right, one has attained the goal of Dhyaana yoga.

The third yoga is the yoga of even-mindedness, Samatva Yoga. Practicing this form of yoga, one performs his duty while established in yoga (v. 48). Do your duty, but be centred! Establish yourself in yoga, renouncing attachment and viewing success and failure alike.

The fourth yoga is the yoga of the Divine glory. In this verse, Bhagavān Sri Kṛṣṇa asks Arjuna to behold His wonderful power. He says, “Behold, see God’s glory in everything.”

The fifth yoga is the yoga of devotion, \textit{bhakti-yoga}. In this verse, Kṛṣṇa says that the yoga of bhakti is when one worships God constantly and exclusively. In this form of yoga, one places God first and does everything for the sake of God’s Love.

The sixth yoga is of the eight limbs, Ashtanga Yoga. This yoga is mentioned as a sacrifice. It is a mixture of Sankhya Yoga and \textit{karma-yoga}. This is the yoga of physical expression. Lord Kṛṣṇa says that one should not be attached to doing physical yoga. One should do it in a state of non-attachment. This will keep you in a healthy state. It will develop a certain energy and vibration that will help you advance in your \textit{sādhana}.

The seventh yoga is Sankhya Yoga, which is explained by Kṛṣṇa throughout this Chapter.

Later on, Arjuna asks, “What is a \textit{yogī}?” We have been talking about the \textit{yogī} here. To become a \textit{yogī}, you have to be centred in all the forms of yoga. When one has mastered these seven steps of yoga, one automatically becomes a \textit{yogī}. The first kind of \textit{yogī} is called a Maha Yogi: a Maha Yogi is God Himself. The second kind of \textit{yogī} is a knower of the Self. When Kṛṣṇa speaks about this \textit{yogī}, He is talking about the one who knows one’s Self. The third kind of \textit{yogī} is one who has surrendered the body, mind, and intellect to God: this is a bhakta. A bhakta, who has surrendered the body, mind and intellect to God, is a true \textit{yogī}. The fourth kind of \textit{yogī} is a Karma \textit{yogī}. A Karma \textit{yogī} is one who performs all actions without attachment, only for the sake of selfpurification. A Sankhya \textit{yogī} is a \textit{yogī} who realizes and identifies with Brahman. A Bhakti \textit{yogī} is completely and constantly in meditation on God, with an undivided mind. He practices in depth the forms of Bhakti, devotion and action, without any personal interest or any personal aim. A Dhyaana \textit{yogī} is a hermit who retires to caves to seclude himself. Hermits are in constant meditation. Focusing on the Divine is the most important thing to them.

So you can see that \textit{yogīs} have different qualities. Later on Kṛṣṇa explains more to Arjuna about the different aspects of yoga and \textit{yogīs}.

\Verse[2.54]
{% --- 1. The Verse ---
  \hspace*{1em}arjuna uvāca \\
  sthita-prajñasya kā bhāṣā samādhi-sthasya keśava \\
  sthita-dhīḥ kiṁ prabhāṣeta kim āsīta vrajeta kim
}
{% --- 2. The Word-for-Word (Clean text, no formatting) ---
  \textit{arjuna uvāca}---Arjuna said; 
  \textit{keśava}---O Keśava; 
  \textit{kā}---what; 
  \textit{bhāṣā}---[is the] description; 
  \textit{sthita-prajñasya}---of one whose intellect is fixed; 
  \textit{samādhi-sthasya}---[and] who is seated in transcendence; 
  \textit{kim}---how; 
  \textit{prabhāṣeta}---does he speak; 
  \textit{sthita-dhīḥ}---one of steady intellect; 
  \textit{kim}---how; 
  \textit{āsīta}---does he sit; 
  \textit{kim}---[and] how; 
  \textit{vrajeta}---does he move
}
{% --- 3. The Translation ---
  \hspace*{1em}Arjuna said: \\
  O Keśava, what is the description of one whose intellect is fixed and who is seated in transcendence? How does such a person speak? How does he sit and how does he move?
}

Arjuna has by now gained profound knowledge. He says, “My friend, here in front of me, is not just a mere mortal. He is Bhagavān Himself, He is God Himself incarnated.” Arjuna addresses Kṛṣṇa as Keśava, Ke-sha-va: ‘Ke’ as Brahma, the creator; ‘Sha’ as Viṣṇu, the preserver; ‘Va’ as Shiva. He says, “You are the Creator, the Protector and the Destroyer of this entire universe: thus I address you as Keśava, my Lord.” The tune has changed: Arjuna is not talking to Kṛṣṇa the same way as before. He is addressing Him with great humility, because he has realized that He is the controller of all three aspects of the Trimurti -- Brahma, Viṣṇu and Shiva. He addresses Him as the Almighty God, saying, “It is only You, my Lord, who can give me the true answer.”

“What is the mode of speech, O Kṛṣṇa, of one of steady wisdom, who is established in the control of the mind? What does one of steady wisdom say?” Arjuna asks, “You are the Lord of lords, you are the God of creation, preservation and destruction. What will you say? How does a \textit{yogī} who is realized sit? How does he move? Please tell me all these qualities of a true, realized soul.” Arjuna now desires to know the signs of one who has attained Realisation, of one who is perfectly and unshakably stable in God Consciousness. Kṛṣṇa has said that He is not attached to anything. So how does one who is not attached to anything walk, how does he talk? How does he eat?

Bhagavān explains, “The perfect souls who have realized God, exhibit special characteristics in all their ways of life.” In whatever they do, they have their own way. They are Siddha Purushas. They are not bound by their actions; they are not bound by their words. By nothing! Today they can say one thing; tomorrow they can say something else. Today they can say, “This is blue!” Tomorrow they can say, “This is red!” Kṛṣṇa says that realized souls are not attached to anything. Even if they do an action which is inappropriate in the minds of society, they are free from every \textit{karma} related to it. Their ordinary speech, their modes of sitting and walking have a special characteristic: it is all through the vibration, and it is all for the sake of others, it is not for their own sake.

Realized souls don’t do anything for themselves, but for the sake of others. The intention of Arjuna in asking about the modes of speech of a self-realized man, is to know how he should speak to the Lord. He wants to know, “How is my action with the Lord? Am I correct or not?” This is selfintrospection. By asking this question, he starts to look at himself and to ask himself, “How is my action? How is my speech? Did I ask properly? How do I engage in any form of action? How am I?” He starts to think about how he walks and how he acts: he becomes very aware of this. Kṛṣṇa has been talking about becoming a great \textit{yogī}. “Kṛṣṇa says to rise to Self-realization, but how will I become that?” Only by self-introspection can one attain Self-realization. That is why in Atma Kriya Yoga we say, “Do daily introspection so that you can transcend and transform yourself.”

\textbf{2016}---Here he is asking, “What is the quality, how will I recognize?” Bhagavān Kṛṣṇa has talked about the wisdom of the Self, which a \textit{yogī} has. This has inspired Arjuna to know, in an inquisitive way, to know in a deep way,” How can I fully grasp in respect of such an individual?”

Because you can’t see what is happening in somebody's mind. You see only what is on the outside. You can’t perceive what knowledge one has inside. You can’t perceive what understanding another person has. So, how to perceive that? Here Bhagavān is telling to Arjuna, “To recognize such a person, it is only through one’s intuition. Only the intuition, this intuitive awareness, this intuitive assessment, can make one recognize these qualities. You may know your own knowledge, you may have this knowledge of the Self, you have this knowledge, as a warrior. Here, my friend, here in front of me, is not just a mere mortal.” Here, Arjuna knows that in front of him is Bhagavān Himself. Step by step, he starts realizing that, “He is not my friend.” The words that He is speaking are something deep, are something which is very powerful, which only God incarnated can speak about. He addresses Kṛṣṇa, “Please, tell me how I can recognize this great humility?”

Here Bhagavān is saying that those who have realized God, they reflect a different Light. These great souls, whatever they do, is not for themselves. They are free; they are not bound by any words. Today they can say something, tomorrow they can say the opposite. It will not make any difference to them. Today, they can do something wonderful, tomorrow they can do something very destructive. It’s equal. But whatever they do, they do in accordance with the will of God. They don’t do it for their own pleasure, they don’t do it for their own attachment or their own want. Even when they walk, even if it appears that they are not doing anything, even if they are piling a mountain of sand, they are always doing something for the benefit of the world. And also here Bhagavān is saying that the \textit{yogī} doesn’t speak much. One word is enough. Everything that they do, even when they walk or sit, is always absorbed into the Divine. By being near them, you feel different. You feel your heart crying for the Lord. The same thing that they have inside of them, they give it to others.

Here, Arjuna is asking; he said he wants to know about his own actions. He wants to know why he is there, how he should behave toward Kṛṣṇa. By asking the question, he is questioning himself, he is asking Bhagavān, “Please, is my action toward You right? Is it proper for me to ask You something? I am not so sure. You see, before I was in a state of confusion, now I am out of that confusion. Now you are showing me. I start to think, have I acted improperly toward You in any way, in any action?”

Here Kṛṣṇa has been telling to Arjuna, “Become a \textit{yogī}. Let this Self-Realization awaken inside.” Arjuna starts to look at himself and starts to wonder, “For Him to speak such words to me, that means there is something inside of me. Why would Kṛṣṇa speak and tell me that? That means there is a \textit{yogī} inside which is dormant, which wants to awaken.

\Verse[2.55]
{% --- 1. The Verse ---
  \hspace*{1em}śrī-bhagavān uvāca \\
  prajahāti yadā kāmān sarvān pārtha mano-gatān \\
  ātmany evātmanā tuṣṭaḥ sthita-prajñas tadocyate
}
{% --- 2. The Word-for-Word (Clean text, no formatting) ---
  \textit{śrī-bhagavān uvāca}---Bhagavān Kṛṣṇa said; 
  \textit{yadā}---when; 
  \textit{prajahāti}---one gives up; 
  \textit{sarvān}---all; 
  \textit{kāmān}---desires; 
  \textit{manaḥ-gatān}---arising in the mind; 
  \textit{pārtha}---O Pārtha; 
  \textit{tuṣṭaḥ}---[is] satisfied; 
  \textit{ātmani}---within oneself; 
  \textit{ātmanā}---by the Self; 
  \textit{eva}---alone; 
  \textit{sthita-prajñaḥ}---of steady wisdom; 
  \textit{tadā}---then; 
  \textit{ucyate}---[one] is said
}
{% --- 3. The Translation ---
  \hspace*{1em}The Lord said: \\
  When one gives up all desires arising in the mind, O Pārtha, and is satisfied within oneself by the Self alone, then one is said to be of steady wisdom.
}

After the total cessation of all desires in the mind, then the sadhak, the bhakta, perceives the true nature of the Eternal, the true nature of God within the Self: within himself, the Supreme Self awakens. The supreme intelligence awakens. The intelligence which Kṛṣṇa talks about before is the intelligence which the mind understands. Here Kṛṣṇa is saying that in the one who is satisfied in the Self Itself, what arises is wisdom, Brahma Jyaan, the knowledge of the Divine, through the continued and devoted practice of attaining the Divine. When the mind is unshakably fixed on God, then one becomes stable, of stable mind. When one finds God within oneself, this is true intelligence, true wisdom.

\textbf{2016---}Bhagavān replies; a common man has all kinds of desires. He is always looking forward to satisfying them for his own selfish motive and sake, and for the sake of the people around him, but in a selfish way.

One who is selfish goes in life looking for short-term happiness, short-term gain. A man who is realized sees this as something transient, something not permanent, and he raises this way of seeing one’s action. One learns to abandon this urge from the desire, the desire for petty things. When one starts to abandon through true knowledge, when one has this understanding of true wisdom, divine knowledge, automatically, these desires which the mind has created start to fade away. This happens automatically. How much importance you give to the outside, the more or less important the outside becomes. But when you have your mind focused on something which is permanent, if you have your mind focused on the Self itself, if you have your mind focused on Lord Kṛṣṇa, then the mind becomes fixed and stable. When one realizes, when the Divine reveals Himself within the Self, then one receives true wisdom, true knowledge. It awakens by itself automatically from deep within. But here, you see, this doesn’t awaken just like that.

Very often, people imagine things, and they say, “Yes, I hear this.” Or, “I see that.” But when you look at the person, you see how unstable the person is.

That’s why Arjuna is asking Bhagavān Kṛṣṇa, “What kind of quality do these people have?” Those who have truly had a glimpse of the Divine, they are calm. The mind is calm, the mind is focused only on Giridhariji. They are not running after a perception of happiness or a thought of happiness, because they have the supreme happiness with them. And that supreme happiness, they want to share it with others, they want to give it to others. They want to help others to center themselves in the \textit{ātmā}. Because only the \textit{ātmā} can give and reveal this true happiness. They will not help you to fulfil your wishes outside. He said, everything will come to you automatically. Everything that has to be for you outside will come to you. You don’t need to do anything for that. You only have to build your relationship with Bhagavān. Nature looks after.

You know, God created this universe, created this whole world itself, and, of course, He has infused Himself in everything, but everything goes by itself, automatically. Everything runs automatically. God doesn’t take a stick and say, “Here has to change. You have to change.” No.

Those who have centered into the \textit{ātmā}, into the Self by the Self, they have found true happiness. He said, “Don’t look for this happiness outside when this happiness, you have it inside of you.”

You see, when you talk about realization, somebody asked me a question once, “When will I be realized?” But you are always realized, it’s just that you don’t know it yet. The \textit{ātmā} inside of you is part of the Supreme. It’s part of God. It’s a parcel of the Divine. Can that parcel not be realized? That is not realized. You think that you are not realized because of that mind. But in essence, that’s why Bhagavān is saying there, all desires from the mind, expel them. When the mind is calm and is satisfied in the Self by the Self, only the Self can give you that true happiness, this peace, this perpetual happiness. When that happens, your mind comes to a conclusion that there is another, it’s a new harmony that the mind perceives, a harmony which the mind has never experienced before. The mind has always run thinking that, “That’s it. This is my work. My duty is to run around left-right, getting busy. He has always shown that he is the master. But once you enter into the Self, you become stable in intelligence. A new harmony starts to dwell in that mind. One gains control over it. One saves that mind from treachery and wandering around. The mind becomes fixated on the Supreme Himself.

\Verse[2.56]
{% --- 1. The Verse ---
  duḥkheṣv anudvigna-manāḥ sukheṣu vigata-spṛhaḥ \\
  vīta-rāga-bhaya-krodhaḥ sthita-dhīr munir ucyate
}
{% --- 2. The Word-for-Word (Clean text, no formatting) ---
  \textit{anudvigna-manāḥ}---one who is unperturbed in mind; 
  \textit{duḥkheṣu}---by suffering; 
  \textit{vigata-spṛhaḥ}---one who is free from craving; 
  \textit{sukheṣu}---after happiness; 
  \textit{vīta-rāga-bhaya-krodhaḥ}---one who is free from attachment, fear, and anger; 
  \textit{ucyate}---is called; 
  \textit{muniḥ}---a sage; 
  \textit{sthita-dhīḥ}---of steady intellect
}
{% --- 3. The Translation ---
  One whose mind is not perturbed by suffering, who does not crave after happiness, who is free from desire, fear and anger—such a person is called a sage of steady intellect.
}

Here Kṛṣṇa says that when the mind is steady and calm, when the mind of a God-realized soul has no thirst for pleasures of the world, when he takes pleasure and pain alike and remains ever-balanced, in whatever kind of experience, he is a Muni. Even more, this intelligence makes him a Muniraj. He is worthy. Because he is settled in understanding, he is allowed to give knowledge, he is allowed to share his experiences of God. Because of his self-controlled and balanced mind, he is allowed to help others. Because of the stability in his state of mind, because he is pure in his actions, he is not bound by anything. He has full control of his speech, like a Muni. When the mind is fully absorbed in the Divine and when one’s speech is for the greater good of the people, one becomes like a Muni, a great sage.

\Verse[2.57]
{% --- 1. The Verse ---
  yaḥ sarvatrānabhisnehas tat tat prāpya śubhāśubham \\
  nābhinandati na dveṣṭi tasya prajñā pratiṣṭhitā
}
{% --- 2. The Word-for-Word (Clean text, no formatting) ---
  \textit{yaḥ}---he who; 
  \textit{anabhisnehaḥ}---has no attachment; 
  \textit{sarvatra}---anywhere; 
  \textit{prāpya}---when encountering; 
  \textit{tat tat}---that; 
  \textit{śubha-aśubham}---the agreeable and the disagreeable; 
  \textit{na abhinandati}---does not feel attraction; 
  \textit{na dveṣṭi}---does not feel aversion; 
  \textit{tasya}---his; 
  \textit{prajñā}---wisdom; 
  \textit{pratiṣṭhitā}---is steady
}
{% --- 3. The Translation ---
  He who has no attachment anywhere, who feels neither attraction nor aversion when encountering the agreeable or the disagreeable—his wisdom is steady.
}

Earlier we saw that Arjuna asks Kṛṣṇa, “How does a \textit{yogī} speak?” Here Kṛṣṇa says that when a realized soul encounters what is “agreeable or the disagreeable, he feels neither attraction nor aversion.” This is true wisdom. Many people have lots of wisdom in many things, but when you encounter them and you say something they don’t want to hear, what do they do? They fight back, no? They rebel!

In the mind of a true \textit{yogī}, there is no rebelling. Someone who is fully dedicated to his path, should disregard all negativity and fully focus on his path: if one does this, nothing can move one – neither praises or insults. When wisdom or knowledge is given by such a great saint and \textit{yogī}, one should drink from it, one should dive into that wisdom.

\textbf{2016}--- Arjuna is asking Bhagavān, “How does a \textit{yogī} speak? When he speaks does he agree or disagree?” Bhagavān said no, because the mind has changed, the mind has moved away, the mind dwells in that inner harmony, they don’t run around, back and forth. They don’t always run to the extremity of the ends. They don’t run into pleasure or pain, failure or success, gain or loss.

If these pairs of opposites are upsetting the mind, it will get diverted. If there is any characteristic of the mind that is still active in a \textit{yogī}, they will still get disturbed. But those who are beyond that, they don’t get disturbed by anything. There will be no expectation inside of them. For them, all experiences are only the Divine. They have this mental calmness in all situations, whether it is agreeable or disagreeable, whether it is attraction or aversion. They look at things and people around them in an unattached, indifferent manner without any hate or anger. They live in this world without being disturbed by how people are around them. They don’t judge people.

Very often, you see people – here coming back to this verse, you can see that Bhagavān is saying, very often, people have great knowledge, they talk about things, yet in their knowledge, there is still judgment, criticism. They are still very rebellious. But those who have attained the Lord, who are centered in the calmness of the mind, they don’t rebel. They don’t concentrate on the negativity of others, and they don’t concentrate on the positivity of others apart from the Lord Himself. You see the mind.

Often, you encounter people who are very wise, very knowledgeable, very intelligent, but when you talk to them, they are very critical about their path. Only they are good. All the others are bad. It’s everywhere like that, no? “My path is best. Your path is wrong. My God is good, your God is not good.” Because they don’t have any true experience of God. Because when you have a true experience of God, this duality, this judgment disappears. It makes place only for Love, nothing else. That’s what Bhagavān Kṛṣṇa is saying here. “Those who have this true wisdom and knowledge, like the saints, the real saints, the \textit{yogīs}, they perceive everything as God perceives.

\Verse[2.58]
{% --- 1. The Verse ---
  yadā saṁharate cāyaṁ kūrmo ’ṅgānīva sarvaśaḥ \\
  indriyāṇīndriyārthebhyas tasya prajñā pratiṣṭhitā
}
{% --- 2. The Word-for-Word (Clean text, no formatting) ---
  \textit{yadā ca}---and when; 
  \textit{ayam}---this (one); 
  \textit{sarvaśaḥ}---completely; 
  \textit{saṁharate}---withdraws; 
  \textit{indriyāṇi}---the senses; 
  \textit{indriya-arthebhyaḥ}---from the sense objects; 
  \textit{iva}---just as; 
  \textit{kūrmaḥ}---a tortoise; 
  \textit{aṅgāni}---[retracts its] limbs; 
  \textit{tasya}---[then] his; 
  \textit{prajñā}---wisdom; 
  \textit{pratiṣṭhitā}---is firmly established
}
{% --- 3. The Translation ---
  When one completely withdraws the senses from the sense objects, just as a tortoise retracts its limbs, then one's wisdom is firmly established.
}

This verse speaks of withdrawing the senses from all objects of self-enjoyment at the time of meditation, thus depriving them of their power to distract the mind and the intellect. Like a tortoise pulls all his limbs inside, leaving only the shell outside, like that, one has to dive within oneself. You can’t pull your hands inside you, but by entering into the cave of the heart, being aware, and sitting in deep, full concentration within oneself, this is true wisdom. But if one is sitting in meditation, yet one is not there, one is somewhere else, that is not true meditation. How many people sit to meditate, but the moment they close their eyes, they ask themselves, “Where am I? I don’t even know where I am!” They are somewhere else.

Once there was a \textit{yogī} who invited all his students to sit for meditation. As they were all sitting, he told them, “Okay, now go into the Self.” Everybody closed their eyes, started breathing and went within. When the meditation was over, the \textit{yogī} asked, “Who was really here during this meditation?” Everybody was thinking, “Oh, I was in America”; “Ah, my parents!” “I was here,” and “I was there.” They were not fully focused. They were all over the place. Even here, how many of you are really listening to me? If I scream, of course, this will bring you back here. But some of you are elsewhere, even though you appear to be listening and watching. You are watching, but you are not seeing anything, because your mind is somewhere else. When the mind is focused in the Self, that is true wisdom. Then it can’t be elsewhere. Nothing can distract it; nothing can move it.

\Verse[2.59]
{% --- 1. The Verse ---
  viṣayā vinivartante nirāhārasya dehinaḥ \\
  rasa-varjaṁ raso ’py asya paraṁ dṛṣṭvā nivartate
}
{% --- 2. The Word-for-Word (Clean text, no formatting) ---
  \textit{viṣayāḥ}---the sense objects; 
  \textit{vinivartante}---slacken their pull; 
  \textit{dehinaḥ}---for someone (the embodied person); 
  \textit{nirāhārasya}---who abstains [from them]; 
  \textit{rasa-varjam}---except the taste; 
  \textit{dṛṣṭvā}---having seen; 
  \textit{param}---the Supreme; 
  \textit{rasaḥ}---taste; 
  \textit{api}---even; 
  \textit{asya}---of this (person); 
  \textit{nivartate}---ceases
}
{% --- 3. The Translation ---
  The sense objects slacken their pull for someone who abstains from them, but their taste remains. Once having seen the Supreme, even that taste vanishes.
}

Each of the senses has a different food. Each sense has its own taste. The food of the ears is sound; the food of the nose is smell. If all of these are renounced, automatically all our concentration will be focused on attaining one food – the divine nectar of the Supreme God Himself. And once the Supreme has revealed Himself, the aspect of the supreme reality which is inside oneself awakens – Amrit. Then one doesn’t need anything outside, then one doesn’t need anything secondary. One won’t need any secondary object to live, because God Himself will become the all-pervading reality for that person. The food – God will be one’s food. All one’s actions will be pervaded by only this One Reality.

Kṛṣṇa says that the only thing that one has to have is God-realization, the vision of God. Without the vision of God, nothing is possible. Otherwise one will always run and hide behind worldly objects. Yogis who possess a stable mind obtain the direct vision of God, the ocean of the supreme bliss. In this state, the \textit{yogī} ceases to have any trace of attachment to worldly objects or to the senses. All the attachment to worldly things totally disappears when there is a direct vision of God. It is through complete surrender and longing for the Divine, that God will reveal Himself. If one is centred in the object on the outside or controlled by the senses, the Divine will never reveal Himself. If one is limited to the duality of ‘good’ and ‘bad’, the Divine will not reveal Himself. It is only when one becomes like a \textit{yogī} who is above this, that the Divine vision is revealed within the core of your own Self, not outside. The true \textit{yogī} firstly perceives this divinity within himself. Then, through that Divine vision, he sees the same divinity equally everywhere in all God’s creation.

\Verse[2.60]
{% --- 1. The Verse ---
  yatato hy api kaunteya puruṣasya vipaścitaḥ \\
  indriyāṇi pramāthīni haranti prasabhaṁ manaḥ
}
{% --- 2. The Word-for-Word (Clean text, no formatting) ---
  \textit{kaunteya}---O son of Kunti (Arjuna); 
  \textit{pramāthīni}---[the] turbulent; 
  \textit{indriyāṇi}---senses; 
  \textit{hi}---indeed; 
  \textit{prasabham}---forcibly; 
  \textit{haranti}---carry away; 
  \textit{manaḥ}---the mind; 
  \textit{api}---[of] even; 
  \textit{vipaścitaḥ}---[a] wise; 
  \textit{puruṣasya}---person; 
  \textit{yatataḥ}---who is striving [to control them]
}
{% --- 3. The Translation ---
  O Kaunteya! The turbulent senses forcefully carry away the mind of even a wise person who is striving to control them.
}

Here Lord Kṛṣṇa speaks about wise men, the wise people, who have the full knowledge of what is ‘good’ and what is ‘not good’. These people “striving to control them (the senses)” are doing their daily \textit{sādhana} and are very focused and intense, but they can’t let go of the senses of sight, smell, taste, sound, and touch. Even the wise ones are bound by these objects of the senses. The sense objects may withdraw, but not the taste for the enjoyment of them. As long as an attachment to the enjoyment of sense objects remains in the heart of man, his senses will lead him to the enjoyment of sense objects. His mind, his intellect, cannot rest stable in God, because of the senses which are carrying away the mind. God, Kṛṣṇa says that one must control all the senses. He says, “Devote the mind to the service of God alone. Practice meditation!” Focus on His Supreme form within the core of one’s heart. Lose oneself in this deepness; then the turbulence, the insistence of the senses will not tempt or bewilder the mind. Thus the mind will enter into this stability. The mind may still churn, but there will be only one aim: to attain the vision of God.

Kṛṣṇa says to Arjuna, “O son of Kunti, become like the \textit{yogī} who has a stable mind and makes a special effort to control his senses by totally renouncing attachment to the objects of the world.” Bhagavān is also saying that when you serve the Lord, you should not expect anything in return – not even the blessing of God, because you can’t make a deal with Him. You can’t say to God, “I love you, I’m doing this for you, but in return I want you to bless me or love me back.”

\Verse[2.61]
{% --- 1. The Verse ---
  tāni sarvāṇi saṁyamya yukta āsīta mat-paraḥ \\
  vaśe hi yasyendriyāṇi tasya prajñā pratiṣṭhitā
}
{% --- 2. The Word-for-Word (Clean text, no formatting) ---
  \textit{saṁyamya}---having controlled; 
  \textit{sarvāṇi}---all; 
  \textit{tāni}---those [senses]; 
  \textit{āsīta}---one should sit (remain); 
  \textit{yuktaḥ}---in the state of meditation; 
  \textit{mat-paraḥ}---focused on Me; 
  \textit{hi}---indeed; 
  \textit{yasya}---of one whose; 
  \textit{indriyāṇi}---senses; 
  \textit{vaśe}---are under control; 
  \textit{tasya}---his; 
  \textit{prajñā}---wisdom; 
  \textit{pratiṣṭhitā}---is firmly established
}
{% --- 3. The Translation ---
  Having controlled all the senses, one should abide in the state of meditation, focusing on Me; for one who has controlled his senses, wisdom is firmly established.
}

“ Having controlled all the senses, one should abide in the state of meditation, focusing on Me.” Here Lord Kṛṣṇa says, “Fully concentrating on Me, ‘for one who has controlled his senses, wisdom is firmly established.’” One has to sit in yoga, in meditation. Giving oneself fully unto Him, with a controlled mind, with a mind which is not distracted, the sadhak has to concentrate on his spiritual growth. He has to go deeper, seeking God-realization, seeking the Grace of God, seeking to attain Him. This must be his aim. Here Lord Kṛṣṇa says that one has to first completely surrender and control the senses. After subduing the senses, if the mind is also brought under control, one must sit in silence and dive within one’s own Self to find Him. Kṛṣṇa says, “Then one is ‘given up to Me!’ The Self of all the selves is only Me. I am the One who is seated in all the bhaktas, in every human heart.” A sadhak, a devotee, should find the Divine, should find the \textit{guru} within his own Self. That will make one fully realized. But if the mind is scattered all around during your daily activities, then even while you sit for meditation, it will be dancing around. That is why whenever you tell people to meditate, they find it very difficult. And this is why japa is important.

When you start chanting the Name of God, you train yourself to control your mind, your heart, and your senses. “Having controlled all the senses,” having renounced all attachment, all-desires, and controlling the mind, the intellect should focus on Kṛṣṇa, on God. For a bhakta, for someone who is really longing for God, the mind should not be on anything else; it should dwell on the Master, and God.

\textbf{2016}---“Having controlled all the senses.” Here Bhagavān is explaining, “One should abide in the state of meditation, focusing on Me; for one who has controlled his senses, wisdom is firmly established.” Having brought all the senses, the five senses under control, when you have started to control what you are looking at, what you are hearing, what you are listening to, what you are eating, what you are doing, by controlling these senses, you will be able to control the mind. By controlling the senses, sit quietly and concentrate upon Me. Here, Bhagavān is saying, concentrate upon Him. Bring your focus upon Him. He is exhorting Arjuna to learn to restrain the senses. He suggested to him a method which will help him to do so. “Sit in a relaxed way, be comfortable, do not move or swing around. Concentrate your mind upon Me.” Don’t become like a pendulum.

Sometimes you see people sitting in meditation, they are shaking, all kinds of movement. Maybe some animal is under them, but here Bhagavān is saying no, you have to be calm, and you have to be steady, because this reflects on the mind also. If you are swinging around and you are shaking around, feeling that the kundalini is rising, it’s not. Because when the kundalini arises, this is not the right – maybe you are possessed, but when kundalini arises, it’s just a vibration that goes through the spine. It’s a deep feeling of joy and happiness that flows inside of you. It’s like a tingling sensation which moves inside of your spine, and that doesn’t make your whole body shake. That’s why those who have truly experienced that, they look at these people who shake themselves like that, “There is some mental problem here.” What to do? That’s the drama of life, sometimes you have to entertain it. When they realize their stupidity, they can move away from that. Like Bhagavān Kṛṣṇa is saying, “Control those senses, don’t let that mind run away, run in all directions.”

Concentrate your mind. How to make the mind stable? He is saying, “Fix your mind upon Me!” He didn’t say fix it in the void or in the emptiness. He said, “Look, I am here in front of you. I have a form. I am not just a concept. I am visible. Focus upon Me. Bring that mind to my Feet.” This is how a sadhak and bhakta, a devotee, should have their mind focused on the Feet of the Master. They should see the \textit{guru} sitting inside of them. They should not let the mind run around. Whatever they are doing, they should do, but with the awareness that they are never alone. \textit{guru} and God is always with you. And when you sit in your meditation, you should perceive that \textit{guru} and God is seated inside of you. Whenever you sit in meditation, and you see your mind is running, who is seeing that mind running? The Self, no? Who are you? Are you the running mind, or are you the Self? So, tell the mind, “Hey, come back! Sit here! Behave!” When you have the mind focused on the Lord, automatically the control over the senses and the mind – if the mind is there, the mind is centering itself upon the Lord, it will not let the senses move left-right. If the mind is not stable, of course, then the senses will also not be stable. But if the mind is focused upon Lord Kṛṣṇa, then there will be no disaster. This controlled mind will help one to attain the purification of the intellect. This prepares one – this is a preparatory stage for the devotee to ascend to the plane of Self-Realization. It makes one ready.

You see, when you start your spiritual path, what do you do? Japa, no? Do japa, chant the Name of God continuously. Prepare yourself, control the mind. Controlling of the mind will give control over all your desires. Focus the mind at the Feet of Giridhari. Chant His Name. Let the mind dwell continuously on the Name. Let it chant at every moment, not only when you have your japa mala with you, but at every moment, it should, it shall chant the Name. Make it so that every breath is a remembrance of His Name. Every action that you do is for His sake only. This is how, when every breath is transformed into the Divine, every action that you do will also be transformed, it will always be a reminder that \textit{guru} and God are with you, and that they are always pushing you to make yourself better.

\Verse[2.62]
{% --- 1. The Verse ---
  dhyāyato viṣayān puṁsaḥ saṅgas teṣūpajāyate \\
  saṅgāt sañjāyate kāmaḥ kāmāt krodho ’bhijāyate
}
{% --- 2. The Word-for-Word (Clean text, no formatting) ---
  \textit{dhyāyataḥ}---[when] deliberating upon; 
  \textit{viṣayān}---sense-objects; 
  \textit{saṅgaḥ}---attachment; 
  \textit{puṁsaḥ}---of a person; 
  \textit{teṣu}---to them; 
  \textit{upajāyate}---arises; 
  \textit{saṅgāt}---from attachment; 
  \textit{sañjāyate}---arises; 
  \textit{kāmaḥ}---desire; 
  \textit{kāmāt}---from desire; 
  \textit{krodhaḥ}---anger; 
  \textit{abhijāyate}---arises
}
{% --- 3. The Translation ---
  When deliberating upon sense-objects, a person develops an attachment to them; from attachment ensues desire, from desire arises anger.
}

Here Kṛṣṇa doesn’t say that you must let go of everything: he is talking about how to let go. This is very important. He says that by constantly dwelling on the objects of enjoyment, the human being develops an intense form of attachment to them. This is true. When the senses dwell on things, like lust and drugs and other things, when the mind and the senses dwell on these many things, you become attached to them and you are not free. This creates more desire to gain even more things.

When one dwells on even one enjoyment in the mind, one gets attached to it and finds it very difficult to let it go. And then you create not one, but many more desires. Kṛṣṇa says, “All these people who are in front of you on the battlefield, Arjuna, these are all the qualities which arise due to desires. And it takes only one desire to create many more attachments.” Duryodhana doesn’t want to let go of his attachment to the kingdom. From his attachment to the seat of Hastinapur have arisen so many desires to gain more and more. That is why the Kauravas also took Indraprastha: they wanted to conquer everything. This is also true in our daily life, when we become obsessed or controlled by greed, pride, and the ego. From one desire, more desires arise. It is never enough. One becomes greedy. This is how attachment gives rise to desires, and when a hindrance appears in the fulfilment of these desires, one becomes angry. “Why didn’t I get this? Why did my neighbour get that? Why not me?” This anger then becomes hatred.

\textbf{2016}---So, Lord Kṛṣṇa is saying that those whose mind, the senses pick up perception from the external object choicelessly. So, it appears, appeal to the senses, you know, it appears pleasant to the senses. But they are very short living. They don’t last for long. Even the most beautiful, you know, external appeal that you have, the most - what you have – anything which makes you extremely happy in the outside doesn’t last forever. Only short-term. You go through it like a big waves, “Paaaf, I am happy!” Because you make an illusion that you are happy. You have made yourself believe that this is true happiness. You will be happy, but does it last forever? A month? Once honeymoon’s time is finished, you are back to normal. What have you achieved? Nothing.

But they keep repeating themselves, you know because due to that short-term happiness – you see, it’s like, have you seen somebody sitting at the slot machine in a casino putting one coin, “The next one I will win!” Put one more coin, no, not happening. Next one, next one, next one and the machine gives 10. They have already spent hundred, they get 10 and they are so happy for the 10 they got. That excitement which makes them more excited to play more. Then they carry on again and again. This becomes an addiction. Any kind of addiction will make one unhappy. But they are so much engrossed into this addiction that they don’t see it. They can’t reason anymore, because their mind is so focused on the outside. If they could take the machine with them home, they would do it. At night for sure when they are sleeping imagine, you know, your partner is next to you, he has this addiction, he will hold your hand thinking it’s a machine. [The audience laughs] And when you scream he will think that the money is coming out. [laughing]

So, where your mind is focused, that becomes your reality. So, this short-term happiness appears the most appealing, you know, but it keeps repeating themselves again and again. It is due to that that attachments are created; due to the wrong perception of things, due to where your concentration is. One is running after, one is hankering after one desire and the other and keep repeating itself more and more. They are never free from that.

“from attachment ensues desire, from desire arises anger.” Because when your desire is not fulfilled, you try, try, try, try, try, what happens? You become frustrated. Your frustration starts arising inside of you. Then you lose everything. You say, “I am doing so much, I’m struggling so hard, I’m keeping struggling myself more and more for vain hope, you know, that the things will get better”, that the next one will be the jackpot. Each coin you are putting, expecting, each desire you are having, you know, it’s arising, next one will be better, next one will be better.

So, this struggle is endless because they are attached to it. This struggle ends up in complete disappointment, in complete anger. So, the root of anger is desire, which brings frustration in your choices. When anger arises, you are already unhappy with yourself, no? Then you will be unhappy with all around you. When you are unhappy with yourself, you can’t be happy with anybody.

So, this is what is happening in the Kauravas, due to greed, you know, want to have more and more, their desire to gain more and more, the Karuavas are not happy. And because of their own unhappiness [49:18] they don’t want anybody to be happy. From one desire many more desires arise. One desire is enough to create lots of hindrances around you, because that one desire for the external reality will make you unhappy? So, how many these desires do you have that bring anger to you? “Why don’t I get this? Why my neighbour gets that? Why not me?”

Last week, no few days ago I got a message from somebody, somebody won a lottery after having a big accident, nearly died, you know, so he won a lottery, very happy. So, he nearly died, two days later he won a lottery of how many millions. Somebody sent me a message, “I am so unlucky! Why this doesn’t happen to me?” What can you say for such a person? Be happy for somebody else, you know. So, this is the reason, you see, because you are not happy for somebody else. You know, you are always thinking about yourself, you know, in a frustrated way that you can’t truly be happy. So, in reality, when you are rebelling against somebody else, you are not rebelling against anybody apart from yourself. It’s just because you can’t see and accept your own fault, you are projecting it on people around you. That’s what Bhagavān is saying; from desire anger comes forth and explodes, paw! But when your mind is controlled, then you don’t disturb anybody. You bring peace for the people around you also.

\Verse[2.63]
{% --- 1. The Verse ---
  krodhād bhavati sammohaḥ sammohāt smṛti-vibhramaḥ \\
  smṛti-bhraṁśād buddhi-nāśo buddhi-nāśāt praṇaśyati
}
{% --- 2. The Word-for-Word (Clean text, no formatting) ---
  \textit{krodhāt}---from anger; 
  \textit{bhavati}---arises; 
  \textit{sammohaḥ}---delusion; 
  \textit{sammohāt}---from delusion; 
  \textit{smṛti-vibhramaḥ}---loss of memory; 
  \textit{smṛti-bhraṁśāt}---from loss of memory; 
  \textit{buddhi-nāśaḥ}---destruction of discrimination; 
  \textit{buddhi-nāśāt}---from the destruction of discrimination; 
  \textit{praṇaśyati}---one is lost
}
{% --- 3. The Translation ---
  From anger arises delusion; from delusion, there is loss of memory; from loss of memory the destruction of discrimination occurs; and with the destruction of discrimination, one is lost.
}

When anger arises in the heart of man, it deprives him of his power of discrimination. He is unable to think, so he will not heed the consequences of whatever he does in a fit of anger. He doesn’t realize what he is doing. When the anger grows, one’s memory starts to get confused, and when the brain starts to get confused, one loses all control. One forgets about relationships; one forgets about everybody around. One also forgets what one has to do, or not to do. In this state, one is unable to plan; one doesn’t have any determination. One loses all reasoning. When one doesn’t control the senses, one is reduced to this state. And when a man is reduced to this state, it becomes easy for him to abandon the path of duty, the path of \textit{dharma}. It becomes easy for that person to lose himself in different paths.

A person who has anger also has qualities such as harshness, violence, greed, and stupidity. His mind is always focused on doing evil behaviour: he is always wanting to harm people. This is wickedness. When a person is angry and can’t control himself, he reduces himself to the state of existence where he has no chance of being reborn as a human. He lowers himself to the level of the animal species or he is born in a very low, infernal degree of the creation.

Here Lord Kṛṣṇa says to Arjuna that if someone doesn’t take advantage of this human life to realize and attain God-realization, attain the Grace of God, or receive His Vision, this life is wasted. Kṛṣṇa says that if this animallike quality arises in the mind of a man, he will become like this animal. This includes how you think, all the food that you eat, and the food which we were talking about earlier, the food of the senses. People are so attached to what they eat! People are so attached to how they think! People are so attached to what they do in their actions! If they have not profited in this lifetime to attain God’s Grace, in the next life, they will be born as an animal.

In this life, God has given them a human body to attain Him; but they have refused the Grace of God, they have refused the gift of God and this is a great sin. People who eat meat, who can’t let go of meat, in their next life, they will surely be re-born as the animals which they have eaten, because it is a curse to eat animals. If there is one thing that can degrade a person, it is a curse which one puts on oneself. If one doesn’t have the chance to renounce all these things, one will be reduced to the state of animal existence.

It is asked in the Upanishads, “regarding the meat eaters, with what attitude do they eat meat?” This is very important. When an animal is being eaten, you are taking in all the sadness, the aggression and other emotions of this animal. When you are putting the meat inside of you, you are also taking in this \textit{karma} of the animal. You are putting this inferior state inside of you. So in your next life, you will be born again to burn this karma. You will be born in the group of animals that you have eaten the most. But, of course, God is merciful: He will always allow you to remember – because animals also can have the perception of realizing that they should change. Even if you have reduced yourself to this degraded state, God is with you. He will give you an experience to remind you. He will say, “You don’t need to do this. You don’t need to go to the infernal regions. You don’t need to reduce yourself and ruin yourself like this.”

In this lifetime itself, you can attain God. This is mercy. God loves you! God cares for you! God incarnates! God traces your life and makes everything possible for you to realize Him! But He can’t force you to do it. He shows you the way to come to Him, through His Love. He says, “I am waiting for you! Please come, realize, awaken! First do your effort, control yourself! I will help you.”

Lord Nārāyaṇa told a story to Maha Lakshmi. “There is a story I will tell you, Devī, about how the sacredness of the Self, the knowing of the Self is very important. I will tell you the story of Deva Sharma. Deva Sharma was born in Shampur and had a great desire to have the knowledge of God. He had this deep desire, so he prayed to Me, ‘Please send me to the right person who will help me get out of this mundane life, so that I can realize myself and attain the perfection of life.’

“One day, Deva Sharma had a visit of a brahmin, a sage. He offered some prasad to the sage and said, ‘Please, reverend sir, tell me, who can guide me so that I can fully surrender to God? Who can guide me so that I can control my senses and my mind can fully focus on the Divine?’ The sage said to him, ‘Go to a certain village and there you will meet your \textit{guru}. On the outskirts of the village, there is a \textit{guru} by the name of sage Nithyananda. Go to him! He will give you Enlightenment and will free you from the illusion of this life.’ So Deva Sharma started on his way to the village.

“Meanwhile in this village, there was a shepherd who had a goat. Every day the shepherd would take the goat to graze in the field. One day as he was going to the field, a tiger jumped out of the forest and roared terribly! But when the tiger jumped in front of the goat, instead of eating the goat, he just looked at it, got scared and ran away. Then the goat was walking very proudly, ‘Din din dun dun.’ The shepherd was shocked and thought, ‘What happened? Why?’ He went to his Gurudev, Nithyananda Maharaj, and revealed to him all that had happened and asked, ‘Why was this? Normally a tiger would eat the goat or would eat me, but it didn’t do anything. It ran away, frightened.’ Of course, with his inner vision, the sage could see the past lives of the goat, and tiger, and the shepherd. Looking into this vision, he said, ‘Ah! The tiger got frightened of the goat because in its last life, this tiger was a handsome man and this goat was a terrible witch. In that last life, the witch had eaten him! The tiger was going to eat the goat, but then saw that it was the witch who had eaten him before. Fearing that also in this life the goat would eat him, he ran away. Due to the witch’s bad karma, she had been reduced to being a goat.’

“Some days later, Deva Sharma arrived at the ashram looking for his \textit{guru}. He sat in the ashram waiting to speak to his \textit{guru} who was there reading the Bhagavad Gita. Then, the shepherd came to the ashram bringing his goat with him. The shepherd also sat there listening the Bhagavad Gita and it was like this every day. The shepherd would bring the goat there, because he knew that his goat had been a witch and thought that it would be good for the goat to hear the Bhagavad Gita. Hearing the Bhagavad Gita would purify the goat. So he brought the goat to listen to the Bhagavad Gita to free it from its pitiful state in this sub-human species – which was due to the goat’s past karma.

“One day the tiger also came to the ashram, scaring everybody. The \textit{guru} called to it and said, ‘Hey, come here!’ He patted it saying, ‘Sit here!’ He started reading the glory of the Gita to the tiger. So all three of them, the shepherd, the tiger, and the goat were sitting there, listening to the Bhagavad Gita. Then the \textit{guru} revealed to them who they had been in their last lives. He said, ‘You shepherd, were a brahmin, you had the knowledge. But you just practiced it mechanically. As you didn’t apply this knowledge inside of you, you have been reduced to being a shepherd. You used to be very wise and learned, but it was only for your personal gain. Since you were always running after personal gain, you were born here, poor and miserable.’ And the other two, you know who they were in their past lives.

“So the \textit{guru} advised the three of them, ‘Concentrate and listen with deep faith to the Gita. And try your best to change.’ The eyes of the tiger and the goat started to shine and whenever they listened to the Gita, there was no enmity between them: they had reached that state of perfection. Even if in their last lives they had not reached it, by the Grace of the Lord, by hearing the Gita, they came to this point of Realisation. So finally they attained salvation, all three of them.”

Here Lord Kṛṣṇa says to Arjuna, “Don’t waste life! Take every opportunity! Don’t be stupid! Don’t be silly! Don’t let this negative quality take control of you. You are the \textit{ātmā}, you are the Great Observer. You are far better than all this! Don’t go into your weakness. Arise!”

\textbf{2016}---Frustration destroys the sense of right and wrong. When the people around you are frustrated, your dear ones are frustrated, when a wife is frustrated – you see, when the husband is frustrated with his wife or children, he forgets that he is the father or the husband. You forget what your duty is. In frustration, you forget where you are standing. When you forget yourself, you forget what your duty is, what your purpose is, due to frustration, you start to get angry. And this anger, this wrath, toward the people around you is so harsh. It doesn’t have any limit to it. Then people do things which are out of control. The good sense of reasoning is totally lost. When one reaches that stage, it is as if one is already dead. You are standing in self-destruction. You are destroying yourself.

Bhagavān Kṛṣṇa said to Arjuna, “If someone doesn’t take the advantage of this human life to realize and attain God-Realization in one life itself, to attain the Grace of God, to control that mind, to control the senses, life becomes wasted.”

If one keeps the animal quality arising in the mind, if you don’t tame the animal inside of you, you will become like the animals. You see, throughout life you have gone through many different species and in your subconscious, these animal qualities are still active. When you start thinking in an animalistic way, then you become like an animal. Animals are doing their work according to their own instinct, but you, you are a human being. You don’t bark like dogs. You don’t kill like a lion. You have certain intelligence, certain knowledge, you have the power of discrimination. But if you have lost this by the arising of anger in your heart, in your mind, then how would you think? How would you act? Anger will consume you in doing something bad, and anger will make you feel as if you are right. You can’t even think about it. And also, it will destroy all your relationships around you, because anger is such a destructive force where you forget yourself, you forget all around. And the determination of anger is destroyed. One loses all reason. So, if one doesn’t control the senses, all reason is lost.

\Verse[2.64]
{% --- 1. The Verse ---
  rāga-dveṣa-vimuktais tu viṣayān indriyaiś caran \\
  ātma-vaśyair vidheyātmā prasādam adhigacchati \\
  prasāde sarva-duḥkhānāṁ hānir asyopajāyate \\
  prasanna-cetaso hy āśu buddhiḥ paryavatiṣṭhate
}
{% --- 2. The Word-for-Word (Clean text, no formatting) ---
  \textit{tu}---but; 
  \textit{vidheya-ātmā}---one who is self-controlled; 
  \textit{caran}---interacting; 
  \textit{viṣayān}---sense objects; 
  \textit{indriyaiḥ}---with senses; 
  \textit{rāga-dveṣa-vimuktaiḥ}---[that are] free from attraction and aversion; 
  \textit{ātma-vaśyaiḥ}---[and] under the control of the Self; 
  \textit{adhigacchati}---attains; 
  \textit{prasādam}---tranquility
  \textit{prasāde}---in [this state of] tranquility; 
  \textit{upajāyate}---[there] arises; 
  \textit{hāniḥ}---the destruction; 
  \textit{sarva-duḥkhānām}---of all sorrows; 
  \textit{asya}---for him; 
  \textit{hi}---certainly; 
  \textit{buddhiḥ}---[the] intellect; 
  \textit{prasanna-cetasaḥ}---of a person with a serene mind; 
  \textit{āśu}---soon; 
  \textit{paryavatiṣṭhate}---becomes firmly established
}
{% --- 3. The Translation ---
  But one who is self-controlled, who interacts with sense objects with senses that are free from attraction and aversion and under the control of the Self, attains tranquility. In this state of tranquility, all his sorrows are overcome. Certainly, the intellect of a person with a serene mind soon becomes firmly established in Me.
}

\textbf{2014}--- “But one who is self-controlled, who interacts with sense objects with senses that are free from attraction and aversion and under the control of the Self, attains tranquility.” Here Bhagavān is not saying that one should not enjoy what one has. Everything has been given by God. One should enjoy it! But one should not become a slave to it. One should not be attached to it! Kṛṣṇa says that if the God realized soul comes into contact with the senses and sense objects, or with like and dislike, these are instruments of advancement. He is not subject to the negativity of these instruments. He can utilise these instruments because he has no attachment to these instruments and he seeks no personal gain with them. Because of the process of controlling the senses, he is free. Even if these qualities of like and dislike awaken inside of him, even if they appear for him, it doesn’t mean anything.

Since the realized soul is centred in the Self, since he is centred in God, nothing bothers him. Whatever he does, he is like a child. A child doesn’t bother about things. A child can ‘like’, and then, the next moment, he ‘doesn’t like’. Children don’t keep grudges; they don’t take revenge. Kṛṣṇa says, “First you have to discipline the senses.” If the senses are not disciplined, you will be reduced to an animal state. But a realized soul, with a disciplined mind, overcomes everything. By chanting the Divine Names, practising japa and meditation, one is ever remembering the Divine in the daily activities.

Like this, one remains free, because whatever one is doing, one is not attached to it. It will not create more attachment. Even if realized souls move among men, even if they are near people who are negative, it doesn’t affect them. No matter what happens, they are fully centred in the Divine: they are free. They are free firstly from the sense of like and dislike. They are renounced from within.

“In this state of tranquility, all his sorrows are overcome. Certainly, the intellect of a person with a serene mind soon becomes firmly established in Me.” Here Lord Kṛṣṇa says, “When the senses are disciplined and when there is no way of liking or disliking, then the heart of a devotee, the heart of a sadhak, becomes transparent, becomes innocent, becomes pure.” In that state, one enjoys spiritual bliss. One becomes peaceful in every situation: one stays in the calmness of the Self. The purity of the heart is a direct cause of joy and calmness. The prasad that we offer to God, is the purity of the heart. When one works on oneself and becomes calm, this becomes a prasad which is offered to the Divine. Whatever one does, it is all prasad. It is all a prayer, and it is constantly bringing one to higher and higher states.


\Verse[2.66]
{% --- 1. The Verse ---
  nāsti buddhir ayuktasya na cāyuktasya bhāvanā \\
  na cābhāvayataḥ śāntir aśāntasya kutaḥ sukham
}
{% --- 2. The Word-for-Word (Clean text, no formatting) ---
  \textit{na asti}---there is no; 
  \textit{buddhiḥ}---discernment; 
  \textit{ayuktasya}---for one who is not self-controlled; 
  \textit{ca}---and; 
  \textit{na}---[there is] no; 
  \textit{bhāvanā}---[proper] awareness; 
  \textit{ayuktasya}---of one without self-control; 
  \textit{ca}---and; 
  \textit{abhāvayataḥ}---for one devoid of meditation; 
  \textit{na}---[there is] no; 
  \textit{śāntiḥ}---peace; 
  \textit{aśāntasya}---for the unpeaceful; 
  \textit{kutaḥ}---how; 
  \textit{sukham}---happiness
}
{% --- 3. The Translation ---
  For one who is not self-controlled, there is no discernment nor proper awareness. For one devoid of such practice, there can be no peace, and without peace, how can there be happiness?
}

\textbf{2014}---The man lacking self-control can have no peace of mind and can’t think of God. He’s always running to the outside world; so how can he think and concentrate on the ocean of the supreme bliss? How can he become calm and tranquil? He will not be able to concentrate, because he doesn’t even know what peace is. The mind of the worldly man remains ever-distracted and his heart remains constantly burning and agitated under the impulse of love and hate. He can’t concentrate, because he is busy either loving or hating. He is always concentrating always on lust and anger, greed and jealousy.

Kṛṣṇa says to Arjuna, “Centre yourself in equalmindedness. Try your best to become equalminded, so that you can free yourself.” You can control your emotions. An uncontrolled mind has no peace and is always aggressive. Without peace of mind, there can be no true happiness. One can create a certain illusion of happiness and think that one is happy. How many people do that? Many! You ask people, “Why are you happy?” They say, “I don’t know why I’m happy. I’m happy because of this. I’m happy because of that.” Anyone who doesn’t have spiritual knowledge, the knowledge of the Self, who is not on the way to loving God, doesn’t know real happiness.

When people talk about love, they have no clue what love is. What they call love is actually just infatuation. They long for caring or being looked after. When you look at how people love in the outside world, you see that all these relationships are fake! All this happiness which people create in the outside world is not true! It’s only true for the mind, because it’s feeding the mind. It makes one feel happy: “Oh I am happy, if someone is pleasing me to be happy. But if I am alone, I am depressed.” It’s true! “Oh, I am terrible! Oh oh…” It’s always about oneself, always about craving for attention! This is what people call happiness. These self-pitying people crave attention. The world revolves around this.

Kṛṣṇa says to Arjuna, “Don’t feel sorry for yourself! What happiness can come of this? What true joy can come of it? This will reduce you to a very pitiful state, to an animal state! So don’t go into that. Be free from that!” In that state, you will never have peace. You will be restless and unhappy! Even if you create an illusion of being happy, the moment this illusion is removed, you will again be unhappy. Wake up! God is the most important thing! He is the main object of enjoyment. He is the Ultimate – not this illusion that you call happiness. When your heart itself is truly happy, then you will always be happy. But when your mind is happy, you should be careful!

Nāsti buddhir-ayuktasya: here Kṛṣṇa says “The stupid mind! It is not intelligent!” People call this state of non-intelligence, intelligence.

\textbf{2016}--- For him who doesn’t have any concentration, he doesn’t have any peace.

The man who lacks self-control doesn’t have peace of mind and can’t concentrate on the Lord. That mind is always running in the outside world, always trying to think and concentrate in a very greedy way. Whereas a harmonized person keeps himself peaceful within. He doesn’t allow himself to be disturbed by anything outside. So, you have self-control, nothing can move you. Nothing that people will say will divert you from your path when you know where you stand, what your path is, who you are following. You will be reasonable, you will have this self-control, self-discipline, over the strong pillar upon which an intelligent person stands. Those who are centered in yoga, they are in self-discipline, and they have a certain personality. They are never diverted. For the one who is in yoga, the one who is concentrating, who is absorbed in the Self. Here Lord Kṛṣṇa is telling to Arjuna, “Center yourself in equal-mindedness. Center yourself again and again. Make your mind equal. Make your mind – control that mind. Try your best to free yourself. Don’t sit and cry, “No, I can’t, I can’t, I can’t! I tried, I tried, I tried.” Keep trying harder! Because an uncontrolled mind will never give you peace. It will always be aggressive, rebellious. It will never give people around you that peace, because you can’t create an illusion of peace.

You see, people always say, “Yes, be peace, we want peace, we want peace!” But nobody is doing anything for that peace. To have peace, you have to become peaceful. And to be peaceful, you have to have a controlled mind, a mind which doesn’t judge. Then you can be peaceful. Otherwise, it’s always left and right. And when you are truly peaceful inside of you, it reflects upon yourself.

Sometimes I sit and observe my devotees. That’s what I do. I can tell that I have some good devotees. Not all, but some. Enough to be happy about. Because, you see, you have a certain glow, you have a certain calmness which many spiritual people don’t have. It’s not to boost your ego,  that I am saying that, but to encourage you to get yourself better. I didn’t say who has it, who doesn’t have it, but often I see that, and it makes me really happy to perceive that.

This is what Bhagavān is saying to Arjuna also: Have this calmness, and this will reflect. Have this true happiness, and it will reflect on the others. Have self-discipline. Control one’s mind, learn to control one’s reaction toward others. When somebody makes a mistake, with the emotions, thoughts, how is your reaction? Because everybody likes to do what they want, no? Everybody is free to do what they want. It doesn’t mean that everybody has to be how you want. You always want people to be how you want. Learn to accept them. Learn to see them from a different view. If you change your specs and start to see the people in a different way, you will not judge them.

When you sit in meditation, when you do your \textit{sādhana}, when you do your concentration, it helps you to cross over this duality. It removes all judgment from you. It gives you peace of mind. And when you have peace of mind, you are happy, and you can make others happy. But if you don’t have peace of mind, you become restless. Then you make everybody unhappy, including yourself.

\Verse[2.67]
{% --- 1. The Verse ---
  indriyāṇāṁ hi caratāṁ yan mano ’nuvidhīyate \\
  tad asya harati prajñāṁ vāyur nāvam ivāmbhasi
}
{% --- 2. The Word-for-Word (Clean text, no formatting) ---
  \textit{hi}---indeed; 
  \textit{manaḥ}---[when] the mind; 
  \textit{yat}---which; 
  \textit{anuvidhīyate}---conforms to; 
  \textit{caratām indriyāṇām}---(of) the wandering senses; 
  \textit{tat}---that (mind); 
  \textit{harati}---carries away; 
  \textit{asya}---his; 
  \textit{prajñām}---intelligence; 
  \textit{iva}---just as; 
  \textit{vāyuḥ}---the wind; 
  \textit{nāvam}---[drives] a ship; 
  \textit{ambhasi}---[across] the water
}
{% --- 3. The Translation ---
  Indeed, when the mind follows the wandering senses, it carries away one’s intelligence, just as the wind drives a ship across the water.
}
\textbf{2014}---When the mind is happy, you should be careful! It is not true happiness, but short-term happiness! “it carries away one’s intelligence, just as the wind drives a ship across the water.” In the Gita, Kṛṣṇa is already talking about ships. It’s amazing, you know? Much later, in the dictionaries of the west, you will find references to ships. But the ship was already mentioned five thousand years ago in Kṛṣṇa’s time! “When the mind follows the wandering senses”: if a man’s mind and senses are not disciplined, are not focused on the Divine, his senses will drag his mind along with them; it may affect the intellect. Drawing away from God, one will pursue the outside and automatically one will degrade oneself.

“ Indeed, when the mind follows the wandering senses, it carries away one’s intelligence” This is true, especially when one doesn’t get attention. Take a wife who is fully in love with her husband and she doesn’t get the attention she wants. How does she become? This is not only true of women – it is true of everyone! The mind is carried away, the understanding is carried away. Deep inside, the heart is saying, “Yes, this is wrong.” You have the feeling that you are doing something wrong, but the reason is not clearly functioning. You can’t even think, you become stupid. If somebody tells you, “You are stupid,” you will not even see it, because there’s no reason. There’s no understanding! This will automatically lead to one’s downfall.

“Just as the wind drives a ship across the water.” A strong wind may affect a boat proceeding towards its destination. The boat may be diverted from its proper course and be tossed about in the high seas, with huge breakers washing the deck and then sink. In such a state, you will drown yourself. But if an expert sailor knows how to sail, he can use the wind to sail his boat and keep it on course. He can use the wind to reach his goal. It is the same for an elevated soul. An elevated soul is above good and bad. He is still in the world, he is a Jivan Mukta, but he uses everything around him as God’s instruments to attain God-realization. He doesn’t sit and cry, but uses the experiences of life to attain God. He sees God’s hand in every aspect of life.

\Verse[2.68]
{% --- 1. The Verse ---
  tasmād yasya mahā-bāho nigṛhītāni sarvaśaḥ \\
  indriyāṇīndriyārthebhyas tasya prajñā pratiṣṭhitā
}
{% --- 2. The Word-for-Word (Clean text, no formatting) ---
  \textit{tasmāt}---therefore; 
  \textit{mahā-bāho}---O mighty-armed Arjuna; 
  \textit{yasya}---whose; 
  \textit{indriyāṇi}---senses; 
  \textit{sarvaśaḥ}---[are] completely; 
  \textit{nigṛhītāni}---restrained; 
  \textit{indriya-arthebhyaḥ}---from sense objects; 
  \textit{tasya}---his; 
  \textit{prajñā}---wisdom; 
  \textit{pratiṣṭhitā}---is firmly established
}
{% --- 3. The Translation ---
  Therefore, O mighty-armed Arjuna, one whose senses are entirely restrained from their objects, his wisdom is firmly established.
}
\textbf{2014}---Again Kṛṣṇa says, “One who is fully focused is centred in the Divine within oneself, and sees the Divine within oneself.”

Kṛṣṇa addresses Arjuna as ‘mahā-bāho’: “He who possesses long, stout, and powerful arms.” He says, “You are not weak, you are a great hero and a great fighter! It should not be difficult for you to subdue your mind, to have your mind under control. It should not be difficult for you to have your senses under control and tame them. So control the senses, Arjuna! Control the senses – sound, touch, taste, sight and smell without any weaknesses. So move! Grow out of this excitement of the senses!”

The soul has enjoyed the senses throughout many life times. The mind has kept the soul trapped due to this long uninterrupted enjoyment of the objects of the senses. Since the beginning of time, the senses have developed a natural attraction for their objects. So get the senses under control! Change them so that you can free yourself! Don’t let your mind be distracted by these impulses. They have always been there trying to stop you. Now move on! You are not alone. God is with you. When you focus your mind and your mind becomes stable, you will see that the Divine is there with you.

\Verse[2.69]
{% --- 1. The Verse ---
  yā niśā sarva-bhūtānāṁ tasyāṁ jāgarti saṁyamī \\
  yasyāṁ jāgrati bhūtāni sā niśā paśyato muneḥ
}
{% --- 2. The Word-for-Word (Clean text, no formatting) ---
  \textit{yā}---[that] which; 
  \textit{niśā}---[is] night; 
  \textit{sarva-bhūtānām}---for all beings; 
  \textit{tasyām}---in that; 
  \textit{saṁyamī}---the self-controlled; 
  \textit{jāgarti}---is awake; 
  \textit{yasyām}---that to which; 
  \textit{bhūtāni}---beings; 
  \textit{jāgrati}---are awake; 
  \textit{sā}---that; 
  \textit{niśā}---[that is] night; 
  \textit{paśyataḥ}---for the seeing; 
  \textit{muneḥ}---sage
}
{% --- 3. The Translation ---
What is night to all beings, to that the self-controlled is awake, and that to which all beings are awake, is night for the enlightened one.
}

  
\textbf{2014}---The experiences of the ignorant man of the world and the man of knowledge are very different. It’s like ‘day’ and ‘night’; but here it’s has nothing to do with actual day and night. Here, the ‘day’ refers to the awakened, the enlightened one, the one who is out of darkness. “What is night to all beings, to that the self-controlled is awake.”  Here, the ‘night’ means all the excitement of the outside world. The self-controlled one is awake, which means he is vigilant, he is guarding: the Self is observing and controlling, so that one doesn’t go into temptation. Such a person is wise. Whereas when one is ignorant, it is like one is in darkness: it is ‘night’ for those beings.

For the enlightened, the one who is centred in the Self, the one who is finding the Divine in the Self, ‘day’ and ‘night’ are not seen as different from each other. During the ‘night’ one doesn’t see, whereas a \textit{yogī} sees, even in the ‘night’.

When one’s mind has been darkened by the cloud of the senses, one doesn’t see, one doesn’t perceive anything. Only when the Light of God-realization, the Light of the \textit{yogī} awakens, does one perceive the nature of God inside oneself and everywhere else. Although every enjoyment in this world and the next is perishable, momentary, transient and full of suffering, enveloped in the darkness of ignorance since the beginning of time, the worldly-minded man regards these enjoyments as everlasting and full of joy.

People who are in the darkness of the mind and are so attached to self-enjoyment, create the illusion of being happy. They are still in darkness; they are still in the ‘night’. But the selfcontrolled, the one who is awake, finds joy: they find God-realization, they find the wisdom of the Divine, embodied within the core of the heart. Then within them is the awakening. They’re not in the state of sleep, they’re not in the state of unawareness: they have the true knowledge of Brahma Jyaan. They realize that it is God who is working through them, that they have Satchitananda inside of them; they have the Truth of God, the knowledge of God, and the bliss of God within themselves. They are in the world, they can use everything in the world, but they aren’t attached to the world. They utilise every action, every thought and everything which comes from the senses, yet they remain separate from the world. They are in the Light, not in the darkness of ignorance.

\textbf{2016}---This is a very nice verse. The sage, the one who is absorbed in the Supreme, the one whose concentration is on the Lord, they are always in an awakened state. An awakened state is in everything that people perceive. You see, somebody says, “This is wrong. This is right.” No? But for those who are awake, they will find that there is a balance between everything. In the wrong, there is right, in the right, there is wrong. It depends on how you are looking at it.

The men of knowledge, those who have experienced the Divine, they see the opposites differently. You see, day and night here don’t mean the day and night which you understand by the mind. The day and night reflect upon the opposites, the difference, how one who is self-controlled sees and how one who is in darkness sees. One who is self-controlled is always in the light. They perceive things in a different way. And those who are seated in darkness, they don’t see anything, and they perceive certain things differently.

Those who are vigilant, who are absorbed in the Self and have a controlled mind, they don’t go into depression. They are not tempted by the outside. These are the wise.

Whereas the ignorant ones, they are always in darkness, they always pity themselves, they are always sitting and crying about themselves. They don’t want to do anything to change anything, but sitting and crying, that’s what they are good at.

The sages, the \textit{yogīs}, they are awakened while the ordinary people sleep, the ordinary men, human beings, are sleeping. The eyes of the sage are always open toward the truth. Whereas the eyes of normal people are closed. Those who are - bhaktas who know about the Love of Giridhari, the Love for God, they see the world differently. Whereas somebody who is in the world, who is very materialistic, they will tell you you are crazy, they will tell you all this is fantasy, because they see things differently. Because they are so concerned about their fear. You, the devotees, are concerned about God, how to attain Him, how to build that love toward Him, how to grow, how to respect, how to love. Whereas, the others, no, “How can I be richer? How can I be more in darkness? How can I be more attached? How can I entertain myself more on the outside?” They have a different realization completely. In darkness, they don’t see anything.

What is real for the \textit{yogī}, it is unreal for the people of the world. And what is real for the people of the world, it is unreal for a controlled mind, a \textit{yogī}. When the mind is darkened by the clouds of senses, desires, one doesn’t see anything. But when one is in the light of God-Realization, Self-Realization, one awakens and perceives the nature of God inside of themselves and everywhere. What is unknown to the common man, what is unknown to the man of the world? The truth about the Self. They don’t know about the Self.

You see, when you look at people walking on the road, they are all walking like zombies. How many people really think that they have an \textit{ātmā}, they have a soul inside of them? Who can teach them about the soul? Nobody. They live life and they die, nobody cares. Who cares? They themselves have never cared about it, about knowing what true happiness is. “What will truly make me happy?” They have thought about it, yes, but they have brushed it under the carpet. Because the soul always reminds one from time to time, but due to their mind being so strong, they can’t. Then they build a certain reality for themselves, “Well, I am happy. I’ve been doing yoga for 20 years,” but they don’t show anything. They are so bound to the outside reality that they are never free.

But the wise ones, they are awakened in the truth, in the Light. The ignorance is like the night of darkness, and the \textit{yogī} is not in that boat. The \textit{yogī} understands what Maya  is and keeps away from that. They have this knowledge. That’s why I said, yoga is not just an exercise that you do. Yoga is also what you are doing right now. It’s a yoga of feeding the mind with the right knowledge.

\Verse[2.70]
{% --- 1. The Verse ---
  āpūryamāṇam acala-pratiṣṭhaṁ samudram āpaḥ praviśanti yadvat \\
  tadvat kāmā yaṁ praviśanti sarve sa śāntim āpnoti na kāma-kāmī
}
{% --- 2. The Word-for-Word (Clean text, no formatting) ---
  \textit{yadvat}---just as; 
  \textit{āpaḥ}---rivers; 
  \textit{praviśanti}---enter into; 
  \textit{samudram}---the ocean; 
  \textit{āpūryamāṇam}---filled; 
  \textit{acala-pratiṣṭham}---steady and immovable; 
  \textit{tadvat}---so; 
  \textit{saḥ}---the one; 
  \textit{yam}---into whom; 
  \textit{sarve}---all; 
  \textit{kāmāḥ}---desires; 
  \textit{praviśanti}---enter; 
  \textit{āpnoti}---attains; 
  \textit{śāntim}---peace; 
  \textit{na}---not; 
  \textit{kāma-kāmī}---one who desires pleasures
}
{% --- 3. The Translation ---
  Just as rivers enter into the sea, which remains full, steady and immovable, so the person into whom all desires flow attains peace, not the one who craves for pleasures.
}
\textbf{2014}---Kṛṣṇa says that, even if every desire, even if the whole world presents itself to the true \textit{yogī}, it makes no difference to him. It’s like the river entering the ocean: when the river enters the ocean, the ocean doesn’t bother about it; it doesn’t change anything in the ocean. Like that, the \textit{yogī} who is self-controlled always remains calm in every situation. He is unshakably established in God Consciousness. He still has concerns for the world, but he has discrimination and is undisturbed. He is ever peaceful. He doesn’t want more: he is satisfied with whatever God provides for Him. One has to be in that state, like the ocean. Whatever you have, be happy. God always gives you what you need. He also knows best when is the right time to give you what you need. If you learn to accept this, no matter what God will give you, it will be for your spiritual growth. Even all these desires will not have any effect on the person who is fully centred in the Divine, whose aim is clear. Nothing will affect his mental state. If he’s peaceful, there will be nothing to worry about. He knows that all is the Will of God. He knows that all comes from God. It is He who gives, and it is He who takes.

\Verse[2.71]
{% --- 1. The Verse ---
  vihāya kāmān yaḥ sarvān pumāṁś carati niḥspṛhaḥ \\
  nirmamo nirahaṅkāraḥ sa śāntim adhigacchati
}
{% --- 2. The Word-for-Word (Clean text, no formatting) ---
  \textit{vihāya}---having abandoned; 
  \textit{sarvān}---all; 
  \textit{kāmān}---desires; 
  \textit{pumān}---the person; 
  \textit{yaḥ}---who; 
  \textit{carati}---lives; 
  \textit{niḥspṛhaḥ}---without craving; 
  \textit{nirmamaḥ}---who is devoid of the sense of \enquote{mine}; 
  \textit{nirahaṅkāraḥ}---who is devoid of ego; 
  \textit{saḥ}---he; 
  \textit{adhigacchati}---attains; 
  \textit{śāntim}---[real] peace
}
{% --- 3. The Translation ---
  Having abandoned all desires, the person who lives without any craving, and does not have any notion of \enquote{I} and \enquote{mine}, attains real peace.
}
\textbf{2014}---So when the big, ‘I’, ‘I’, ‘I’, and ‘mine’, ‘mine’, ‘mine’, and ‘this belongs to me’ – when all these are removed from the mind, one becomes peaceful. The feeling of ‘I want this,’ and ‘I want that,’ brings pain to the people of the outside world. Here Kṛṣṇa says that the one who has great expectations in life, who desires many things, in a very egoistic way, doesn’t have peace of mind. Even if one wants to be peaceful, it’s very difficult, because this selfimportance, this egoistic feeling, will eat one up. What you need is humility. Because of your own self-importance, you will never be able to respect others.

Arjuna has asked Kṛṣṇa, “How does a realized soul walk in the world?”(2.54) Kṛṣṇa says, “The realized soul walks in the world, but they don’t do it for their own self-gratification, they don’t say, ‘I want this, I want that.’ They accept everything and when things are not there, they don’t bother about it. So let go of this ‘I’ and ‘mine’, because this is the cause of misery! This will be the cause of your downfall! Because of this egoistic self-importance which you give to yourself, you will never attain peace, you will not be peaceful.”

\textbf{2016}---Here Bhagavān keeps on reminding, how one will have peace. How you will have peace? By not concentrating upon the petty self, “I, I, I, I am doing this, I am doing that. That belongs to me. That’s me who is doing it.” You can talk as “I,” but don’t give importance, “Oh, I do that, I do this, I do that.”

This big “I” will bring lots of expectations in your life. It will bring, it will awaken lots of desires for many things with a very big ego. You see, when you see people talking sometimes, especially – also spiritual people. You see them talking, they are so proud of what they are doing. They don’t have any humility. If you are doing yoga, you are doing spiritual things, humility has to be there, humbleness must be there. But if there is no humility, you will never have peace. You want to experience the Love of God? It’s only by being humble. Being humble doesn’t mean you take a rope and beat yourself saying, “I am bad, I am bad, I am bad,” but your mind is concentrating on the bad. “I am good, I am good, I am good.” That’s what you should do. Then your concentration is on the good.

“Having abandoned all desires, the person who lives without any craving.” Here Bhagavān is not saying that you should just let go of everything, desire for basic needs, for survival, like food, shelter, clothing, these things you need, but it’s not that important. Here to renounce, to cut away is to have the self-control. What He is talking about extinguishing the sense of egoistic self-importance, He said to remove that from you. Remove yourself from this. Don’t run behind the sense objects. Long for something which is permanent, long for Giridhari’s Grace. Don’t long for wealth, external comfort, desire for name, fame, glory, pride, ego, vanity. All this will not make you free. It will make you more miserable, because it will keep you away from realizing yourself. All this will awaken the egoistic self-importance, will awaken this big “I-ness” inside of you. This will blind the mind with your body itself. The “mine-ness,” “I-ness,” will become predominant. This “I-ness” and “mine-ness,” that’s what keeps you away from realizing yourself. You start having self-importance. You are thinking that you are such a great personality, that the world can’t go on without you. Bhagavān is saying, that will not give you peace. This relationship binds you firmly to the senses, to the arrogance and pride, to the ego. When the mind is overpowered by that, you will never be at peace with yourself. And when you are not at peace with yourself, you forget about your true nature. And those who are attached to Giridhari, where the mind dwells on the Name, those who continuously do the japa, chanting the Divine Name, that mind reflects the divine quality. When you talk with somebody who is fully absorbed into that – you know what a true devotee who is absorbed into the Divine Name will say? He will never say, “I did that.” No, he will say, “It is the Grace of \textit{guru} and the Grace of Giridhari which has done it. It is Giridhari’s will that things have happened. It is by \textit{guru’s} Grace that I have achieved something. By \textit{guru’s} Grace that I am in this world and it is by \textit{guru’s} Grace that I will attain the Grace of Thakurji. I try my best, but I am surrendered to Him.”

Here there is also “I,” but the importance is not on the “I.” Notice that properly! When somebody who is in the ego and pride, when they talk about “I,” they are like this, “I have done this!” “That is mine!” No? But a devotee is different. He says, “Through the Grace of \textit{guru} and God, I have achieved that.” The importance is not upon “I” as the ego, but through the Grace, it’s only Giridhari and \textit{guru}.

From that point, you see that there is humility, there is humbleness. And this is where He said the \textit{guru} takes that disciple and everything is given, even Bhagavān Himself reveals Himself. And then the devotee himself, because he doesn’t claim that he is doing it, he knows that behind him there is a greater power which is doing everything, he will always have peace of mind. And his peace of mind will be reflected not only for himself, but he will become an example to the world, to everybody.

\Verse[2.72]
{% --- 1. The Verse ---
  eṣā brāhmī sthitiḥ pārtha naināṁ prāpya vimuhyati \\
  sthitvāsyām anta-kāle ’pi brahma-nirvāṇam ṛcchati
}
{% --- 2. The Word-for-Word (Clean text, no formatting) ---
  \textit{eṣā}---this [is the]; 
  \textit{brāhmī}---Brahman; 
  \textit{sthitiḥ}---state; 
  \textit{pārtha}---O Pārtha (Arjuna); 
  \textit{prāpya}---having attained; 
  \textit{enām}---which; 
  \textit{na vimuhyati}---one is not deluded; 
  \textit{sthitvā}---by abiding; 
  \textit{asyām}---in this [state]; 
  \textit{api}---even; 
  \textit{anta-kāle}---at the time of death; 
  \textit{ṛcchati}---[one] attains; 
  \textit{brahma-nirvāṇam}---the bliss of Brahman
}
{% --- 3. The Translation ---
  This is the realized state, O Pārtha, having attained which one is no longer deluded. By abiding in this state even at the hour of death, one attains the bliss of Brahman.
}
\textbf{2014}---He who succeeds in attaining Brahman during this lifetime is the best flower of humanity. He enjoys the bliss of the Divine while living a liberated life: he is a Jivan Mukta. Even he who succeeds at the last moment of his life, suddenly, as a reward for his \textit{sādhana}, or as a Grace from his \textit{guru}, if he is fixing his mind unshakably on the image of the Divine, free from egoism, attachment, and desires, he will enjoy a blissful state of existence in the next world. He will enjoy God’s Heaven, enjoy the existence of Nārāyaṇa, eternally. However, if one has not done one’s spiritual practices, it will be quite difficult to focus the attention on the Divine at the last moment. If you are very much into the outer routine, if you are very attached to the outside world, then at the last moment, your mind will focus on what you’re attached to. So how will it be possible to focus on the Divine at the last moment? It will be very difficult!

Kṛṣṇa says that for people who are focused on the outside material world, it will be very difficult to attain God if they have not practised focusing the mind, if they have not practised yoga, if they have not practised meditation, if they have not practised japa. How could they, who are centred in the ego which always wants more and more, how could they who are centred in all the negative qualities of the outside world, how could they attain God? How could they overcome all this egoism, possession, attachment, and desire at the last moment? When one at last succeeds in reaching God-realization, all this delusion, which has existed from the beginning, gets eradicated, gets removed and it will never re-appear again. “This is the realized state, O Pārtha”. Having attained the state of a \textit{yogī} who has overcome delusion, one becomes a blessing to humanity. One who has awakened the Divine within the heart, who has awakened Bhakti within the heart, who sees the Lord, who worships the Lord, who has surrendered to the Lord, is a great blessing to humanity.

Bhagavān Kṛṣṇa says that one who is centred in spirituality is a blessing for this world, because wherever that person goes, he carries the Divine’s blessing him. Wherever they go, whether they know it or not, they bring a positiveness with them. It’s like a flower: if a flower is put in the middle of a place where everything else is not beautiful, everyone’s attention will focus on that flower. Be like a flower in the middle of this world of ignorance. Spread the sweetness, the fragrance, and the beauty which God has given you! In this state, you are a blessing to humanity, you are a blessing to the world. Realizing this state, you become a Jivan Mukta. You walk in this world being free.

“By abiding in this state even at the hour of death, one attains the bliss of Brahman.” Here Kṛṣṇa says that even he who succeeds at the last moment of his life, either suddenly, or as a reward for his \textit{sādhana}, to fix his mind strongly on the Divine, on God, he becomes free from egoism, attachment, possession, and desires. Grace can transform everybody, even at the last minute. Even if throughout your life, you have practised your \textit{sādhana}, yet you still perceive all these negative qualities inside you, don’t be distressed about it. Even with doing all the spiritual exercises, all the prayers, all the japa, you don’t feel that you have reached the state that we have been talking about – don’t worry! Because at the time of passing, at the time of leaving this world, all this becomes good \texti{puṇya}! All this will help you; it will become like a blessing! At that moment, whatever you have done, if you have perceived results in your life or not; if you have been doing your practice mechanically or not, all this will become a help. This will help you at the time you are dying, at the time you are leaving this world. The \textit{sādhana} that you have done will become a blessing for your Realisation: it will open up your eyes, even if it is just for a second! Even if it is just for a moment, the moment you depart from this world – you are delivered.

This is the sweetness of the Lord. Nothing is done in vain! Even the little that you do, will bear its fruit. This is His compassion: He does everything for the sake of mankind. Kṛṣṇa says that you don’t need to be discouraged when you don’t see any results from your spiritual practices, because all these results are being collected. This good \texti{puṇya}, even if you don’t see it here, it is in Heaven, in the bank account which you have made in Heaven. And this will help you. Even if you don’t feel anything throughout your whole life when you do your \textit{sādhana}, just do it! Don’t think about it. Enjoy it when you are doing it! Don’t compare yourself with anybody to see if you are superior or inferior; don’t worry about that.

Here Kṛṣṇa says to Arjuna, “Your duty as a warrior is to fight. Fight! Don’t think about it. The whole problem starts when you start thinking. When you start thinking, you start comparing. When you start comparing, you lose everything. There is no peace of mind. And when there is no peace of mind, everything is destroyed.” This is what Lord Kṛṣṇa teaches in this chapter of the Bhagavad Gita.

\textbf{2016}---Here Bhagavān Kṛṣṇa is saying, if one has this humility, if one has this peace inside of them, one enjoys the bliss while living, and one will enjoy the Divine even after one has left this plane. When one moves away from pride and ego, one attains peace of mind. This is a natural state of the Self. It’s a state of fullness and completeness. Here is not only fullness and happiness, but completeness also. All the shackles of bondage are shattered. Delusion is evaporated, is gone, disappeared. Ignorance, which was created and conditioned by the mind, is wiped out. Here, that’s what He said; the \textit{karma} is wiped out completely. The mind becomes crystal clear, the mind is calm. Those who fix upon the Lord, the mind becomes calm, because the Lord cleanses and purifies you, and He looks after you, He cares for you, He carries you.

Then the soul starts to shine. This is where you are awake. Awakening, enlightenment, this is a state of enlightenment, this is the state of awakening. This is where bliss and the Grace of God start to reflect through you. You see the sweetness of the Name, the sweetness of being a devotee. Whatever you do on your spiritual path is not in vain. Even the little things you do bear great fruits when it is done with love and surrender, when it is done for the sake of mankind, the sake of serving. You see, when you serve, it’s not to boost one’s ego or pride. You should forget about yourself. When you pray, you should forget about yourself. Because the Lord who is looking after you, He cares for you and He will never let you down. Everybody else will let you down one day. Not only let you down, they will put you underground. But Bhagavān will not let you down, He will raise you and will put you next to Him, and you will enjoy His presence eternally.

Here ends the second chapter of the Bhagavad Gita on the knowledge of the Self.